\documentclass[UTF8,zihao=-4,openany]{ctexbook}
\usepackage[a4paper,top=2.1cm,bottom=2.2cm,left=2.2cm,right=2.2cm]{geometry}
\usepackage{amsmath,amssymb,mathtools}
\usepackage{tikz}
\usepackage{enumitem}
\usepackage{xcolor}
\usepackage{booktabs}
\usepackage{tabularx}
\usepackage{float}
\usepackage[most]{tcolorbox}
\usepackage{fancyhdr}
\usepackage{hyperref}
\usepackage{eso-pic}
\usetikzlibrary{calc,arrows.meta,patterns,decorations.pathmorphing}

\definecolor{Ink}{HTML}{1F2937}
\definecolor{Muted}{HTML}{64748B}
\definecolor{Brand}{HTML}{0B6B62}
\definecolor{BrandDark}{HTML}{0F3F3C}
\definecolor{Accent}{HTML}{C57A1B}
\definecolor{Soft}{HTML}{F2FBF8}
\definecolor{SoftWarm}{HTML}{FFF7ED}
\definecolor{ProblemBg}{HTML}{FBFCFE}
\definecolor{Answer}{HTML}{B45309}
\definecolor{Line}{HTML}{D5DEE8}
\definecolor{TitleBand}{HTML}{E8EEF5}

\hypersetup{colorlinks=true,linkcolor=BrandDark,urlcolor=BrandDark}
\setlength{\parindent}{0pt}
\setlength{\parskip}{0.45em}
\setlength{\headheight}{26pt}
\setlist[enumerate]{leftmargin=2em,itemsep=0.25em,topsep=0.25em}
\pagestyle{fancy}
\fancyhf{}
\fancyhead[L]{结构力学作业精编}
\fancyhead[R]{\leftmark}
\fancyfoot[C]{\thepage}
\renewcommand{\headrulewidth}{0.4pt}
\renewcommand{\headrule}{\hbox to\headwidth{\color{Line}\leaders\hrule height \headrulewidth\hfill}}

\tcbset{
  enhanced,
  boxrule=0.5pt,
  arc=1.2mm,
  left=8pt,
  right=8pt,
  top=8pt,
  bottom=8pt,
  fonttitle=\bfseries,
}
\newtcolorbox{problemcard}[2][]{
  breakable,
  colback=ProblemBg,
  colframe=Line,
  colbacktitle=TitleBand,
  title={#2},
  coltitle=Ink,
  borderline west={2.8pt}{0pt}{Brand},
  attach boxed title to top left={xshift=0pt,yshift=0pt},
  boxed title style={boxrule=0pt,arc=0pt,left=8pt,right=8pt,top=3pt,bottom=3pt},
  #1
}
\newtcolorbox{solutioncard}{
  breakable,
  colback=Soft,
  colframe=BrandDark,
  colbacktitle=BrandDark,
  coltitle=white,
  title={答案与解析}
}
\newtcolorbox{notecard}{
  breakable,
  colback=SoftWarm,
  colframe=Accent,
  colbacktitle=Accent,
  coltitle=white,
  title={讲解要点}
}
\newcommand{\answer}[1]{\textcolor{Answer}{\bfseries #1}}
\newcommand{\source}[1]{\textcolor{Muted}{\small #1}}
\newcommand{\figtitle}[1]{%
  \node[font=\small,text=Muted,align=center,anchor=north]
    at ($(current bounding box.south)+(0,-0.18)$) {#1};
}

\tikzset{
  member/.style={line width=0.85pt,draw=Ink,line cap=round,line join=round},
  link/.style={line width=0.75pt,draw=Ink,line cap=round,line join=round},
  pin/.style={circle,draw=Ink,fill=white,inner sep=1.4pt,line width=0.7pt},
  smallpin/.style={circle,draw=Ink,fill=white,inner sep=1.05pt,line width=0.65pt},
  rigid/.style={draw=Ink,pattern=north east lines,pattern color=Ink!45},
  ground/.style={line width=0.65pt,draw=Ink},
}

\ctexset{
  chapter={
    format={\centering\bfseries\zihao{2}},
    name={第,章},
    number=\chinese{chapter},
    aftername=\quad,
    beforeskip=18pt,
    afterskip=28pt
  },
  section={
    format={\bfseries\zihao{3}},
    beforeskip=1.4em,
    afterskip=0.8em
  }
}

\newcommand{\ground}[2]{%
  \draw[ground] ($(#1,#2)+(-0.25,0)$)--($(#1,#2)+(0.25,0)$);
  \foreach \x in {-0.18,-0.06,0.06,0.18}
    \draw[ground] ($(#1,#2)+(\x,0)$)--($(#1,#2)+(\x-0.08,-0.12)$);
}
\newcommand{\roller}[2]{%
  \node[smallpin] at (#1,#2) {};
  \draw[ground] ($(#1,#2)+(-0.22,-0.17)$)--($(#1,#2)+(0.22,-0.17)$);
  \draw[ground] ($(#1,#2)+(-0.16,-0.17)$)--($(#1,#2)+(-0.24,-0.28)$);
  \draw[ground] ($(#1,#2)+(0,-0.17)$)--($(#1,#2)+(-0.08,-0.28)$);
  \draw[ground] ($(#1,#2)+(0.16,-0.17)$)--($(#1,#2)+(0.08,-0.28)$);
}
\newcommand{\fixedbase}[2]{%
  \draw[ground] ($(#1,#2)+(-0.26,0)$)--($(#1,#2)+(0.26,0)$);
  \foreach \x in {-0.18,-0.06,0.06,0.18}
    \draw[ground] ($(#1,#2)+(\x,0)$)--($(#1,#2)+(\x-0.08,-0.12)$);
}
\newcommand{\wallroller}[2]{%
  \draw[ground] ($(#1,#2)+(-0.18,-0.28)$)--($(#1,#2)+(-0.18,0.28)$);
  \foreach \y in {-0.18,-0.06,0.06,0.18}
    \draw[ground] ($(#1,#2)+(-0.18,\y)$)--($(#1,#2)+(-0.30,\y-0.08)$);
  \node[smallpin] at (#1,#2) {};
}

\newcommand{\makecleanfront}{%
  \begin{titlepage}
  \thispagestyle{empty}
  \begin{tikzpicture}[remember picture,overlay]
    \fill[BrandDark] (current page.north west) rectangle ([yshift=-8.1cm]current page.north east);
    \fill[Soft] ([yshift=-8.1cm]current page.north west) rectangle (current page.south east);
    \draw[Brand,line width=1.2pt] ([xshift=2.1cm,yshift=-8.35cm]current page.north west)--([xshift=-2.1cm,yshift=-8.35cm]current page.north east);
    \foreach \x in {0,...,9}{
      \draw[white,opacity=0.09,line width=0.7pt] ([xshift=-1cm+\x*2.3cm,yshift=-0.5cm]current page.north west)--([xshift=4cm+\x*2.3cm,yshift=-7.9cm]current page.north west);
    }
  \end{tikzpicture}
  \vspace*{1.4cm}
  {\color{white}\zihao{4}\bfseries 结构力学基础与技巧班}\par
  \vspace{1.2cm}
  {\color{white}\zihao{0}\bfseries 作业 · 答案 · 讲解笔记}\par
  \vspace{3.4cm}
  \begin{tcolorbox}[
    width=0.82\textwidth,
    colback=white,
    colframe=Line,
    boxrule=0.6pt,
    arc=1mm,
    left=14pt,right=14pt,top=13pt,bottom=13pt
  ]
  \large
  本稿按题目、答案、解析和讲解要点重新编排；结构图统一使用 TikZ 绘制，公式与控制值保留为可编辑 LaTeX。
  \end{tcolorbox}
  \vfill
  {\color{Muted}\small 编译方式：XeLaTeX \quad 生成日期：\today}\par
  \end{titlepage}
}

\begin{document}
\makecleanfront
\frontmatter
\chapter*{整理说明}
本稿以可编辑 LaTeX 形式重构题目、答案和讲解要点。题图按原题结构重新绘制，答案中的关键控制值和符号约定统一整理，缺少推导处补充平衡方程与方法说明。
\tableofcontents
\mainmatter
\chapter{几何组成分析}

\section{判别方法速记}

\begin{notecard}
几何组成分析的核心不是套公式，而是判断体系能否作为刚片逐步并入大地。常用顺序如下：
\begin{enumerate}
  \item 先识别大地、刚片、链杆、铰和支座约束。
  \item 对平面体系，可先用计算自由度
  \[
    W=3m-2h-r
  \]
  作必要检查，其中 \(m\) 为刚片数，\(h\) 为单铰数，\(r\) 为支座链杆数或外部约束数。
  \item 再用二元体规则、两刚片规则和三刚片规则作充分性判断。
  \item 若约束方向平行且等长，常见为常变；平行但不等长，常见为瞬变。若已经几何不变但约束超过必要数，则存在多余约束。
\end{enumerate}
\end{notecard}

\section{第 1--6 题}

\begin{problemcard}{1-1 几何组成判别 \source{浙江大学 2008}}
图 1-56 所示体系的几何组成为（\quad）。

\begin{tabularx}{\textwidth}{@{}XX@{}}
A. 几何瞬变，无多余约束 & B. 几何不变 \\
C. 几何常变 & D. 几何瞬变，有多余约束
\end{tabularx}

\begin{figure}[H]
\centering
\begin{tikzpicture}[scale=0.92]
  \draw[rigid] (0,2) rectangle (3.9,2.35);
  \draw[rigid] (0,0) rectangle (3.9,0.35);
  \draw[rigid] (0,0.35) rectangle (0.38,2);
  \foreach \x/\dx in {0.75/0,1.55/0.08,2.45/0.22,3.65/0.38}{
    \node[pin] at (\x,2) {};
    \node[pin] at (\x+\dx,0.35) {};
    \draw[link] (\x,2)--(\x+\dx,0.35);
  }
  \figtitle{图 1-56：四根链杆近似平行且不等长}
\end{tikzpicture}
\end{figure}
\end{problemcard}

\begin{solutioncard}
\textbf{答案：}\answer{D}

\textbf{解析：}把上下两个阴影部分视为两块刚片。两刚片之间由 4 根链杆联系，若要形成无多余约束的几何不变体系，平面两刚片之间通常需要 3 个相互独立的约束；这里有 4 根链杆，因此从数量上看存在多余约束。

进一步看约束方向：4 根链杆方向近似平行，但长度不等，不能形成稳定的几何约束组合。按讲解笔记中的口诀，“平行且等长为常变，平行不等长为瞬变”，故该体系为几何瞬变，并且有两个多余约束。
\end{solutioncard}

\begin{problemcard}{1-2 补加支座约束 \source{中南大学 2011；合肥工业大学 2009}}
欲使图 1-57 所示体系成为无多余约束的几何不变体系，则需在 \(A\) 端加入（\quad）。

\begin{tabularx}{\textwidth}{@{}XXXX@{}}
A. 固定铰支座 & B. 固定支座 & C. 滑动铰支座 & D. 定向支座
\end{tabularx}

\begin{figure}[H]
\centering
\begin{tikzpicture}[scale=0.88]
  \coordinate (A) at (2.9,1.15);
  \coordinate (B) at (2.35,2.4);
  \coordinate (C) at (3.45,2.4);
  \coordinate (D) at (2.35,3.05);
  \coordinate (F) at (3.45,3.05);
  \coordinate (E) at (0.65,0.25);
  \coordinate (G) at (5.15,0.25);
  \draw[member] (0.6,0.25)--(0.6,3.05)--(D);
  \draw[member] (D)--(B)--(C)--(F)--(5.2,3.05)--(5.2,0.25);
  \draw[member] (B)--(A)--(C);
  \draw[member] (A)--(2.9,0.75);
  \node[pin] at (B) {};
  \node[pin] at (C) {};
  \node[pin] at (D) {};
  \node[pin] at (F) {};
  \node[pin,label=below left:\(E\)] at (E) {};
  \node[pin,label=below right:\(G\)] at (G) {};
  \roller{0.6}{0.25}
  \roller{5.2}{0.25}
  \fixedbase{2.9}{0.75}
  \node[below right] at (A) {\(A\)};
  \node[above left] at (B) {\(B\)};
  \node[above right] at (C) {\(C\)};
  \node[above] at (D) {\(D\)};
  \node[above] at (F) {\(F\)};
  \figtitle{图 1-57：在 \(A\) 处补固定端}
\end{tikzpicture}
\end{figure}
\end{problemcard}

\begin{solutioncard}
\textbf{答案：}\answer{B}

\textbf{解析：}先作自由度必要检查。原体系可取 3 个刚片、1 个单铰和 4 个外部约束，则
\[
  W=3\times 3-2\times 1-4=3 .
\]
要使体系成为几何不变体系，还需要补足 3 个独立约束。固定端正好提供两个平动约束和一个转动约束，共 3 个约束，因此候选项为固定支座。

加入固定端后，刚片 \(ABC\) 先与大地合并；随后刚片 \(DE\) 可通过链杆 \(BD\) 与铰 \(E\) 并入大地，刚片 \(FG\) 亦可用同样方式并入大地。全过程没有额外约束，故为无多余约束的几何不变体系。
\end{solutioncard}

\begin{problemcard}{1-3 两刚片规则 \source{中南大学 2009}}
图 1-58 所示体系的几何构造性质是（\quad）。

\begin{tabularx}{\textwidth}{@{}XX@{}}
A. 常变体系 & B. 瞬变体系 \\
C. 无多余联系的几何不变体系 & D. 有多余联系的几何不变体系
\end{tabularx}

\begin{figure}[H]
\centering
\begin{tikzpicture}[scale=0.78]
  \coordinate (A) at (3.1,0.7);
  \coordinate (B) at (1.55,3.0);
  \coordinate (C) at (4.75,3.0);
  \coordinate (D) at (0.55,0.35);
  \coordinate (E) at (5.75,0.35);
  \draw[member] (0.5,0.35)--(0.5,2.0)--(B)--(2.75,1.75)--(2.75,0.7);
  \draw[member] (2.75,1.75)--(C)--(5.8,2.0)--(5.8,0.35);
  \draw[link,dashed] (D)--(B);
  \draw[link,dashed] (E)--(C);
  \node[pin,label=below:\(A\)] at (A) {};
  \wallroller{3.1}{0.7}
  \node[pin,label=above:\(B\)] at (B) {};
  \node[pin,label=above:\(C\)] at (C) {};
  \node[pin,label=below:\(D\)] at (D) {};
  \node[pin,label=below:\(E\)] at (E) {};
  \roller{0.5}{0.35}
  \roller{5.8}{0.35}
  \figtitle{图 1-58：三根不共线且不交于一点的约束}
\end{tikzpicture}
\end{figure}
\end{problemcard}

\begin{solutioncard}
\textbf{答案：}\answer{C}

\textbf{解析：}将上部折线框架视为刚片 \(ABC\)，大地视为另一刚片。二者之间由三根链杆联系：\(A\)、\(BD\)、\(CE\)。这三根链杆不共线，且其作用线不交于同一点，满足两刚片规则。

由于正好提供使两刚片相对固定所需的三个独立约束，没有多余联系，因此该体系为无多余联系的几何不变体系。
\end{solutioncard}

\begin{problemcard}{1-4 常变体系与多余约束 \source{广西大学 2010}}
图 1-59 所示体系内部是 \underline{\hspace{3em}} 体系，它有 \underline{\hspace{3em}} 个多余约束。

\begin{figure}[H]
\centering
\begin{tikzpicture}[scale=0.8]
  \foreach \x in {0,1.3,2.6,3.9,5.2}{
    \node[pin] at (\x,2) {};
  }
  \foreach \x in {0,1.3,2.6,3.9,5.2}{
    \node[pin] at (\x,0) {};
  }
  \draw[member] (0,2)--(5.2,2)--(5.2,0)--(0,0)--cycle;
  \draw[member] (1.3,2)--(1.3,0);
  \draw[member] (2.6,2)--(2.6,0);
  \draw[member] (3.9,2)--(3.9,0);
  \draw[member] (0,2)--(1.3,0);
  \draw[member] (1.3,2)--(2.6,0);
  \draw[member] (3.9,2)--(5.2,0);
  \draw[member] (3.9,0)--(5.2,2);
  \wallroller{0}{0}
  \roller{5.2}{0}
  \figtitle{图 1-59：内部桁式体系}
\end{tikzpicture}
\end{figure}
\end{problemcard}

\begin{solutioncard}
\textbf{答案：}\answer{几何常变；1}

\textbf{解析：}采用两刚片规则分析。可把左侧部分看成刚片 \(ABCD\)，右侧部分看成刚片 \(EFGH\)。右侧刚片内部本身已有一个多余约束，而左右两刚片之间仅由链杆 \(CE\) 和 \(DF\) 联系。

两刚片若要组成几何不变整体，需要三根不共线且不交于一点的链杆，或等效的三个独立约束。这里只有两根链杆，不能限制两个刚片之间的全部相对运动，因此体系为几何常变。由于右侧刚片内部已有一个多余约束，故多余约束数为 1。
\end{solutioncard}

\begin{problemcard}{1-5 计算自由度 \source{哈尔滨工业大学 2011}}
图 1-60 所示体系的计算自由度 \(W=\underline{\hspace{4em}}\)。

\begin{figure}[H]
\centering
\begin{tikzpicture}[scale=0.86]
  \coordinate (P0) at (0,1.25);
  \coordinate (P1) at (1.6,1.25);
  \coordinate (P2) at (3.2,1.25);
  \coordinate (P3) at (4.6,1.25);
  \coordinate (P4) at (4.6,3.0);
  \coordinate (P5) at (6.0,1.25);
  \draw[member] (P0)--(P1)--(P2)--(P3)--(P5)--(P4)--(P3);
  \draw[member] (P2)--(3.2,0.05);
  \draw[member] (P2)--(4.6,0.05);
  \draw[member] (P4)--(P5);
  \node[pin] at (P0) {};
  \node[pin] at (P1) {};
  \node[pin] at (P2) {};
  \node[pin] at (P3) {};
  \node[pin] at (P4) {};
  \node[pin] at (P5) {};
  \wallroller{0}{1.25}
  \fixedbase{3.2}{0.05}
  \roller{4.6}{0.05}
  \node[above] at (3.2,1.35) {5};
  \node[above] at (4.05,1.35) {2};
  \node[right] at (3.75,0.65) {3};
  \node[left] at (3.2,0.65) {4};
  \node[right] at (5.2,2.0) {1};
  \figtitle{图 1-60：按刚片编号计算自由度}
\end{tikzpicture}
\end{figure}
\end{problemcard}

\begin{solutioncard}
\textbf{答案：}\answer{\(-4\)}

\textbf{解析：}按答案图所示将体系划分为 5 个刚片，体系中共有 6 个单铰，支座约束折算为 4 个外部约束，并有 1 个复铰折算为 3 个约束。于是计算自由度为
\[
  W = 5\times 3 - 6\times 2 - 1\times 3 - 4 = -4 .
\]
计算自由度为负说明约束数超过自由度数，体系中存在多余约束。注意：\(W<0\) 只是必要检查结果，最终几何组成仍需结合刚片规则判断。
\end{solutioncard}

\begin{problemcard}{1-6 二元体规则 \source{中南大学 2012}}
图 1-61 所示体系几何不变且有多余联系。（\quad）

\begin{figure}[H]
\centering
\begin{tikzpicture}[scale=0.74]
  \coordinate (A) at (0,0); \coordinate (B) at (1.4,0);
  \coordinate (C) at (1.4,1.25); \coordinate (D) at (0,1.25);
  \coordinate (E) at (1.4,2.5); \coordinate (F) at (0,2.5);
  \coordinate (G) at (2.9,2.5); \coordinate (H) at (2.9,0);
  \coordinate (I) at (1.4,3.75); \coordinate (J) at (0,3.75);
  \coordinate (K) at (-1.2,3.1); \coordinate (L) at (-1.2,2.5);
  \coordinate (M) at (-2.5,2.5);
  \draw[member] (A)--(B)--(C)--(D)--cycle;
  \draw[member] (D)--(F)--(E)--(C);
  \draw[member] (C)--(E);
  \draw[member] (F)--(J)--(I)--(E);
  \draw[member] (E)--(G)--(H);
  \draw[member] (I)--(G);
  \draw[member] (M)--(L)--(F);
  \draw[member] (M)--(K)--(J);
  \draw[member] (K)--(L);
  \draw[member] (K)--(F);
  \foreach \p/\lab in {A/A,B/B,C/C,D/D,E/E,F/F,G/G,H/H,I/I,J/J,K/K,L/L,M/M}{
    \node[pin,label=above:{\scriptsize \(\lab\)}] at (\p) {};
  }
  \wallroller{0}{0}
  \roller{0}{0}
  \roller{1.4}{0}
  \roller{2.9}{0}
  \figtitle{图 1-61：可由大地逐步添加二元体}
\end{tikzpicture}
\end{figure}
\end{problemcard}

\begin{solutioncard}
\textbf{答案：}\answer{错误，记作 \(\times\)}

\textbf{解析：}该体系可由大地开始，依次添加二元体
\[
  AB,\ ACB,\ ADC,\ DEC,\ DFE,\ EGH,\ EIG,\ FJI,\ FKJ,\ KLF,\ KML .
\]
每添加一个二元体，都是用两根不共线链杆把一个新节点稳定地接入已有几何不变部分，因此不会引入多余约束。逐步添加完成后，整体为无多余约束的几何不变体系。

题干说“几何不变且有多余联系”，前半句正确，后半句错误，所以判断题答案为 \(\times\)。
\end{solutioncard}

\section{第 7--13 题}

\begin{problemcard}{1-7 刚片并入与多余约束 \source{哈尔滨工业大学 2009}}
对图 1-62 所示体系进行几何组成分析。

\begin{figure}[H]
\centering
\begin{tikzpicture}[scale=0.9]
  \coordinate (A) at (0,0); \coordinate (B) at (1.55,0);
  \coordinate (C) at (3.1,0); \coordinate (D) at (0,2);
  \coordinate (E) at (1.55,2); \coordinate (F) at (3.1,2);
  \draw[member] (A)--(D)--(E)--(F)--(C);
  \draw[member] (B)--(E);
  \foreach \p/\lab in {A/A,B/B,C/C,D/D,E/E,F/F}{
    \node[pin,label=above:{\scriptsize \(\lab\)}] at (\p) {};
  }
  \fixedbase{0}{0}
  \fixedbase{1.55}{0}
  \fixedbase{3.1}{0}
  \figtitle{图 1-62：三立柱框架体系}
\end{tikzpicture}
\end{figure}
\end{problemcard}

\begin{solutioncard}
\textbf{结论：}\answer{有 3 个多余约束的几何不变体系}

\textbf{解析：}刚片 \(AD\)、\(BE\) 可直接并入大地，链杆 \(DE\) 因而成为一个多余约束。随后刚片 \(EFC\) 与大地通过固定端 \(C\) 和铰 \(E\) 联系，满足两刚片规则；由于固定端已提供较强约束，此处还会带来两个多余约束。

因此整体虽然几何不变，但共有 \(1+2=3\) 个多余约束。
\end{solutioncard}

\begin{problemcard}{1-8 三刚片规则 \source{武汉科技大学 2009}}
对图 1-63 所示体系作几何组成分析。

\begin{figure}[H]
\centering
\begin{tikzpicture}[scale=0.72]
  \coordinate (A) at (0,0); \coordinate (B) at (0,2.7);
  \coordinate (C) at (2.2,2.7); \coordinate (D) at (4.4,2.7);
  \coordinate (E) at (4.4,0); \coordinate (G) at (2.2,0);
  \coordinate (F) at (2.2,1.55);
  \draw[member] (A)--(B)--(C)--(D)--(E)--(G)--(A);
  \draw[member] (B)--(F)--(D);
  \draw[member] (A)--(C)--(E);
  \draw[member] (B)--(G)--(D);
  \draw[member] (C)--(F)--(G);
  \draw[member,dashed] (-1.0,-0.35)--(F)--(5.6,-0.35);
  \node[left] at (-1.0,-0.35) {\scriptsize \(O_1\)};
  \node[right] at (5.6,-0.35) {\scriptsize \(O_2\)};
  \foreach \p/\lab in {A/A,B/B,C/C,D/D,E/E,F/F,G/G}{
    \node[pin,label=above:{\scriptsize \(\lab\)}] at (\p) {};
  }
  \wallroller{0}{0}
  \roller{0}{0}
  \roller{4.4}{0}
  \figtitle{图 1-63：三铰不共线的三刚片体系}
\end{tikzpicture}
\end{figure}
\end{problemcard}

\begin{solutioncard}
\textbf{结论：}\answer{无多余约束的几何不变体系}

\textbf{解析：}按三刚片规则分析。刚片 \(ABC\) 与刚片 \(CDE\) 交于铰 \(C\)；刚片 \(ABC\) 与刚片 \(FG\) 由链杆 \(BF\)、\(AG\) 联系，其作用线交于虚铰 \(O_2\)；刚片 \(CDE\) 与刚片 \(FG\) 由链杆 \(DF\)、\(GE\) 联系，其作用线交于虚铰 \(O_1\)。

三个铰 \(C,O_1,O_2\) 不共线，满足三刚片规则，可组成一个大刚片；该大刚片再由三根不共线且不交于一点的链杆支座与大地联系，满足两刚片规则，所以体系几何不变且无多余约束。
\end{solutioncard}

\begin{problemcard}{1-9 三铰共线瞬变 \source{武汉理工大学 2008}}
对图 1-64 所示杆件体系进行几何构造分析。

\begin{figure}[H]
\centering
\begin{tikzpicture}[scale=0.78]
  \coordinate (A) at (0,2); \coordinate (B) at (1.3,2);
  \coordinate (C) at (1.3,0.85); \coordinate (D) at (4.1,2);
  \coordinate (E) at (5.4,2); \coordinate (F) at (4.1,0.85);
  \coordinate (G) at (2.7,0);
  \draw[member] (A)--(B)--(D)--(E);
  \draw[member] (B)--(C)--(F)--(D);
  \draw[member] (A)--(C);
  \draw[member] (D)--(F)--(G)--(C);
  \draw[member] (A)--(G)--(E);
  \draw[member,dashed] (E)--(6.4,2);
  \draw[member,dashed] (F)--(6.4,0.85);
  \node[right] at (6.4,2) {\scriptsize \(O_1\)};
  \foreach \p/\lab in {A/A,B/B,C/C,D/D,E/E,F/F,G/G}{
    \node[pin,label=above:{\scriptsize \(\lab\)}] at (\p) {};
  }
  \roller{0}{2}
  \roller{5.4}{2}
  \roller{2.7}{0}
  \figtitle{图 1-64：三铰共线导致瞬变}
\end{tikzpicture}
\end{figure}
\end{problemcard}

\begin{solutioncard}
\textbf{结论：}\answer{有 1 个多余约束的几何瞬变体系}

\textbf{解析：}取刚片 \(ABC\)、刚片 \(DEF\) 和大地为三刚片。刚片 \(ABC\) 与刚片 \(DEF\) 由链杆 \(BD\)、\(CF\) 联系，二者作用线平行，可视为交于无穷远铰 \(O_1\)；刚片 \(ABC\) 与大地由链杆 \(A\)、\(CG\) 联系，交于铰 \(A\)；刚片 \(DEF\) 与大地由链杆 \(FG\)、\(E\) 联系，交于铰 \(E\)。

三个等效铰共线，不满足三刚片规则，所以体系瞬变。由于约束数量超过必要数，体系同时有 1 个多余约束。
\end{solutioncard}

\begin{problemcard}{1-10 计算自由度并判别 \source{华南理工大学 2012}}
计算图 1-65 所示体系的计算自由度，并分析其几何组成。

\begin{figure}[H]
\centering
\begin{tikzpicture}[scale=0.73]
  \foreach \x in {0,1.2,2.4,3.6,4.8}{
    \node[pin] at (\x,2.2) {};
    \node[pin] at (\x,0) {};
  }
  \node[pin] at (1.2,1.1) {};
  \node[pin] at (3.6,1.1) {};
  \draw[member] (0,0)--(4.8,0)--(4.8,2.2)--(0,2.2)--cycle;
  \draw[member] (1.2,2.2)--(1.2,1.1)--(0,0);
  \draw[member] (1.2,1.1)--(2.4,2.2)--(3.6,1.1)--(4.8,2.2);
  \draw[member] (1.2,1.1)--(2.4,0)--(3.6,1.1)--(4.8,0);
  \draw[member] (0,2.2)--(1.2,1.1);
  \draw[member] (3.6,2.2)--(3.6,1.1);
  \wallroller{0}{0}
  \roller{0}{0}
  \roller{4.8}{0}
  \figtitle{图 1-65：有一个多余约束的几何不变体系}
\end{tikzpicture}
\end{figure}
\end{problemcard}

\begin{solutioncard}
\textbf{结论：}\answer{\(W=-1\)，有 1 个多余约束的几何不变体系}

\textbf{解析：}按题图统计，计算自由度为
\[
  W=12\times 2-25=-1 .
\]
再作几何组成判断：刚片 \(ABCD\) 与刚片 \(DEFB\) 通过铰 \(B\) 和铰 \(D\) 联系，相当于四个约束。两刚片规则只需三个独立约束即可使二者相对固定，因此满足几何不变条件，同时多出一个约束。
\end{solutioncard}

\begin{problemcard}{1-11 超静定次数 \source{青岛理工大学 2012}}
对图 1-66 所示平面体系进行几何组成分析，并指出结构超静定次数。

\begin{figure}[H]
\centering
\begin{tikzpicture}[scale=0.78]
  \coordinate (A) at (0,0); \coordinate (B) at (0,2.2);
  \coordinate (C) at (2.4,2.2); \coordinate (D) at (2.4,0);
  \coordinate (E) at (4.8,0); \coordinate (F) at (4.8,2.2);
  \coordinate (G) at (6.0,2.2); \coordinate (H) at (6.0,0);
  \draw[member] (A)--(B)--(C)--(F)--(G);
  \draw[member] (C)--(D);
  \draw[member] (E)--(F);
  \draw[member] (D)--(E)--(H);
  \draw[member] (B)--(D);
  \draw[member] (D)--(C);
  \draw[member] (C)--(E);
  \draw[member] (F)--(H);
  \foreach \p/\lab in {A/A,B/B,C/C,D/D,E/E,F/F,G/G,H/H}{
    \node[pin,label=above:{\scriptsize \(\lab\)}] at (\p) {};
  }
  \wallroller{0}{0}
  \roller{0}{0}
  \roller{2.4}{0}
  \roller{6.0}{0}
  \wallroller{6.0}{2.2}
  \figtitle{图 1-66：两次超静定体系}
\end{tikzpicture}
\end{figure}
\end{problemcard}

\begin{solutioncard}
\textbf{结论：}\answer{有 2 个多余约束的几何不变体系，即两次超静定}

\textbf{解析：}先去掉二元体 \(AB\) 和 \(CEG\)，不改变体系基本几何性质。剩余部分可看作带有两个多余约束的刚片 \(BCD\)，它与大地由三根不共线且不交于一点的链杆 \(D\)、\(CF\)、\(CG\) 联系，满足两刚片规则。

因此体系几何不变，同时存在两个多余约束，结构超静定次数为 2。
\end{solutioncard}

\begin{problemcard}{1-12 几何稳定性分析 \source{东南大学 2008}}
分析图 1-67 所示体系的几何稳定性，写出分析过程。

\begin{figure}[H]
\centering
\begin{tikzpicture}[scale=0.82]
  \coordinate (A) at (0,0); \coordinate (B) at (3.2,0);
  \coordinate (F) at (6.4,0); \coordinate (C) at (1.5,1.55);
  \coordinate (E) at (4.9,1.55); \coordinate (D) at (0.6,3.0);
  \coordinate (G) at (5.8,3.0);
  \draw[member] (A)--(B)--(F);
  \draw[member] (A)--(C)--(B)--(E)--(F);
  \draw[member] (C)--(D);
  \draw[member] (E)--(G);
  \draw[member] (C)--(1.5,0);
  \draw[member] (E)--(4.9,0);
  \foreach \x in {1.0,2.1,4.2,5.3}{
    \node[smallpin] at (\x,0.35) {};
    \draw[member] (\x,0.35)--(\x,0);
  }
  \foreach \p/\lab in {A/A,B/B,C/C,D/D,E/E,F/F,G/G}{
    \node[pin,label=above:{\scriptsize \(\lab\)}] at (\p) {};
  }
  \wallroller{0}{0}
  \roller{0}{0}
  \roller{6.4}{0}
  \wallroller{0.6}{3.0}
  \wallroller{5.8}{3.0}
  \figtitle{图 1-67：含一根多余杆 \(GE\)}
\end{tikzpicture}
\end{figure}
\end{problemcard}

\begin{solutioncard}
\textbf{结论：}\answer{有 1 个多余约束的几何不变体系}

\textbf{解析：}刚片 \(ABC\) 与大地由三根不共线且不交于一点的链杆支座 \(A\) 和链杆 \(CD\) 联系，满足两刚片规则，可先并入大地。随后刚片 \(FBE\) 与大地由铰 \(B\) 和链杆 \(F\) 联系，也满足两刚片规则。

在上述并入过程完成后，杆 \(GE\) 已不再是维持几何不变所必需的约束，因此为多余约束。体系为有一个多余约束的几何不变体系。
\end{solutioncard}

\begin{problemcard}{1-13 等效刚片分析 \source{山东建筑大学 2011}}
试对图 1-68 所示体系作几何组成分析。

\begin{figure}[H]
\centering
\begin{tikzpicture}[scale=0.88]
  \coordinate (P1) at (0,0);
  \coordinate (P2) at (1.3,0);
  \coordinate (P4) at (2.8,0);
  \coordinate (P5) at (3.6,-1.0);
  \coordinate (P6) at (4.8,-1.0);
  \coordinate (P3) at (5.8,0);
  \draw[member] (P1)--(P2)--(P4)--(P3);
  \draw[member] (P4)--(P5)--(P6)--(P3);
  \draw[member] (P4)--(P6);
  \draw[member] (P5)--(P3);
  \foreach \p/\lab in {P1/1,P2/2,P3/3,P4/4,P5/5,P6/6}{
    \node[pin,label=above:{\scriptsize \(\lab\)}] at (\p) {};
  }
  \roller{0}{0}
  \roller{1.3}{0}
  \roller{5.8}{0}
  \wallroller{5.8}{0}
  \node[font=\scriptsize,Muted] at (4.35,-1.55) {刚片 \(3456\) 可等效为链杆 \(34\)};
  \figtitle{图 1-68：等效为两刚片规则判断}
\end{tikzpicture}
\end{figure}
\end{problemcard}

\begin{solutioncard}
\textbf{结论：}\answer{无多余约束的几何不变体系}

\textbf{解析：}刚片 \(3456\) 可等效为链杆 \(34\)。等效后，取刚片 \(14\) 与大地分析：二者由三根不共线且不交于一点的链杆联系，满足两刚片规则。

约束数正好满足几何不变所需的独立约束数，因此体系为无多余约束的几何不变体系。
\end{solutioncard}

\section{第 14--17 题}

\begin{problemcard}{1-14 二元体去除法 \source{山东建筑大学 2008}}
试对图 1-69 所示体系作几何组成分析。

\begin{figure}[H]
\centering
\begin{tikzpicture}[scale=0.82]
  \coordinate (A) at (0,0); \coordinate (B) at (1.45,0);
  \coordinate (C) at (2.9,0); \coordinate (D) at (4.35,0);
  \coordinate (E) at (5.8,0);
  \coordinate (J) at (0,2.0); \coordinate (K) at (1.45,2.0);
  \coordinate (L) at (2.9,2.0); \coordinate (M) at (4.35,2.0);
  \coordinate (N) at (5.8,2.0);
  \draw[member] (A)--(J)--(K)--(B)--(A);
  \draw[member] (K)--(L)--(C)--(B);
  \draw[member] (L)--(M)--(D)--(C);
  \draw[member] (M)--(N)--(E)--(D);
  \draw[member] (J)--(B);
  \draw[member] (K)--(C);
  \draw[member] (M)--(C);
  \foreach \p in {A,B,C,D,E,J,K,L,M,N}{\node[pin] at (\p) {};}
  \wallroller{0}{0}
  \roller{0}{0}
  \roller{1.45}{0}
  \roller{2.9}{0}
  \wallroller{4.35}{0}
  \roller{5.8}{0}
  \figtitle{图 1-69：可先去除若干二元体}
\end{tikzpicture}
\end{figure}
\end{problemcard}

\begin{solutioncard}
\textbf{结论：}\answer{无多余约束的几何不变体系}

\textbf{解析：}先去掉二元体 \(ABC\)、\(AJ\)、\(DEF\) 和 \(GFH\) 等不改变几何性质的局部单元。剩余主刚片 \(JCDI\) 与大地由三根不共线且不交于一点的链杆支座联系，满足两刚片规则。

整个体系可由大地逐步并入，且没有超过必要数的独立约束，因此为无多余约束的几何不变体系。
\end{solutioncard}

\begin{problemcard}{1-15 计算自由度与常变体系 \source{中国地震局工程力学研究所 2010}}
计算图 1-70 的计算自由度，并作几何构造分析。

\begin{figure}[H]
\centering
\begin{tikzpicture}[scale=0.88]
  \coordinate (A) at (0,0); \coordinate (B) at (1.1,0);
  \coordinate (C) at (2.2,0); \coordinate (D) at (3.3,0);
  \coordinate (E) at (4.4,0);
  \coordinate (P) at (2.2,1.5);
  \coordinate (Q) at (3.3,1.5);
  \draw[member] (A)--(B)--(P)--(Q)--(D)--(E);
  \draw[member] (B)--(C)--(D);
  \draw[member] (B)--(Q);
  \draw[member] (P)--(D);
  \draw[member] (P)--(C);
  \draw[member] (Q)--(C);
  \foreach \p in {A,B,C,D,E,P,Q}{\node[pin] at (\p) {};}
  \roller{0}{0}
  \roller{4.4}{0}
  \figtitle{图 1-70：自由度为正的常变体系}
\end{tikzpicture}
\end{figure}
\end{problemcard}

\begin{solutioncard}
\textbf{结论：}\answer{\(W=1\)，无多余约束的几何常变体系}

\textbf{解析：}按题图统计可得计算自由度
\[
  W=7\times 3-2\times 10=1 .
\]
自由度为正，说明体系仍有独立运动可能。

再用两刚片规则分析：刚片 \(ABC\) 与刚片 \(DEF\) 仅由链杆 \(CD\)、\(BE\) 联系，独立约束不足三个，不满足两刚片规则。因此体系不是几何不变体系，而是无多余约束的几何常变体系。
\end{solutioncard}

\begin{problemcard}{1-16 三刚片共线判别 \source{西南交通大学 2008}}
对图 1-71 所示体系进行几何构造分析。

\begin{figure}[H]
\centering
\begin{tikzpicture}[scale=0.82]
  \coordinate (A) at (0,0); \coordinate (B) at (1.6,0);
  \coordinate (C) at (3.2,0); \coordinate (D) at (4.8,0);
  \coordinate (E) at (1.6,1.3); \coordinate (F) at (3.2,1.3);
  \coordinate (G) at (2.4,2.65); \coordinate (H) at (2.4,-1.25);
  \draw[member] (A)--(B)--(C)--(D);
  \draw[member] (A)--(E)--(F)--(D);
  \draw[member] (B)--(E);
  \draw[member] (C)--(F);
  \draw[member] (E)--(G)--(F);
  \draw[member] (B)--(H)--(C);
  \draw[member] (E)--(C);
  \draw[member] (B)--(F);
  \foreach \p/\lab in {A/A,B/B,C/C,D/D,E/E,F/F,G/G,H/H}{
    \node[pin,label=above:{\scriptsize \(\lab\)}] at (\p) {};
  }
  \roller{0}{0}
  \roller{4.8}{0}
  \roller{2.4}{-1.25}
  \wallroller{4.8}{0}
  \figtitle{图 1-71：三等效铰共线}
\end{tikzpicture}
\end{figure}
\end{problemcard}

\begin{solutioncard}
\textbf{结论：}\answer{有 1 个多余约束的几何瞬变体系}

\textbf{解析：}先去掉二元体 \(ABC\)。随后按三刚片规则分析：刚片 \(ADE\) 与大地由链杆 \(D\)、\(EH\) 联系，作用线交于无穷远铰 \(O_1\)；刚片 \(CFG\) 与大地由链杆 \(F\)、\(GH\) 联系，交于铰 \(O_2\)；刚片 \(ADE\) 与刚片 \(CFG\) 由链杆 \(AC\)、\(EG\) 联系，交于无穷远铰 \(O_3\)。

三个等效铰共线，不满足三刚片规则，体系瞬变；同时约束数量超过必要数，有 1 个多余约束。
\end{solutioncard}

\begin{problemcard}{1-17 计算自由度与多余约束 \source{北京工业大学 2010}}
计算图 1-72 所示体系的计算自由度，并进行几何组成分析。

\begin{figure}[H]
\centering
\begin{tikzpicture}[scale=0.85]
  \coordinate (A) at (0,3.0); \coordinate (B) at (2.3,3.0);
  \coordinate (C) at (2.3,1.55); \coordinate (D) at (4.4,1.55);
  \coordinate (E) at (2.3,0); \coordinate (P) at (0,0);
  \draw[member] (A)--(B)--(C)--(E);
  \draw[member] (C)--(D);
  \draw[member] (E)--(P);
  \foreach \p/\lab in {A/A,B/B,C/C,D/D,E/E}{
    \node[pin,label=above:{\scriptsize \(\lab\)}] at (\p) {};
  }
  \wallroller{0}{3.0}
  \roller{2.3}{0}
  \wallroller{0}{0}
  \wallroller{4.4}{1.55}
  \node[font=\scriptsize,text=Muted,anchor=west] at (4.65,2.05) {杆 \(CE\) 为多余约束};
  \figtitle{图 1-72：有一个多余约束的几何不变体系}
\end{tikzpicture}
\end{figure}
\end{problemcard}

\begin{solutioncard}
\textbf{结论：}\answer{\(W=-1\)，有 1 个多余约束的几何不变体系}

\textbf{解析：}按刚片和约束统计，计算自由度为
\[
  W=4\times 3-3\times 2-7\times 1=-1 .
\]
先将刚片 \(AB\) 并入大地；随后杆 \(CD\) 与大地由三根不共线且不交于一点的链杆 \(BC\) 和支座 \(D\) 联系，满足两刚片规则。杆 \(CE\) 在此基础上不再是维持几何不变所必需的约束，因此为多余约束。

所以体系为有一个多余约束的几何不变体系。
\end{solutioncard}


\chapter{理论力学回顾：支座反力}

\section{平衡方程速记}

\begin{notecard}
求支座反力时，先把结构作为整体隔离。能用一个力矩方程直接消去未知反力时，优先对该支座或铰点取矩；再用
\[
  \sum F_x=0,\qquad \sum F_y=0,\qquad \sum M=0
\]
补齐其余反力。分布荷载先等效为合力：均布荷载 \(q\) 作用长度 \(L\) 时合力为 \(qL\)，作用在荷载段中点；三角形荷载最大值为 \(q_0\)、底边为 \(l\) 时合力为 \(\frac12q_0l\)，作用在距最大端 \(l/3\) 处。若算得某反力为负，表示实际方向与预先假定方向相反。
\end{notecard}

\section{第 2 次作业：求支座反力}

\begin{problemcard}{2-1(a) 三角形分布荷载悬臂梁}
求图示结构的支座反力。梁 \(AB\) 长为 \(l\)，右端 \(B\) 固定，三角形分布荷载从 \(A\) 端为零线性增大到 \(B\) 端的 \(q_0\)。

\begin{figure}[H]
\centering
\begin{tikzpicture}[scale=0.95,>=Stealth]
  \draw[member] (0,0)--(4,0);
  \draw[ground] (4,-0.55)--(4,0.55);
  \foreach \y in {-0.42,-0.24,-0.06,0.12,0.30,0.48}
    \draw[ground] (4,\y)--(4.18,\y+0.10);
  \foreach \x/\h in {0.55/0.20,1.05/0.34,1.55/0.48,2.05/0.62,2.55/0.76,3.05/0.90,3.55/1.04}
    \draw[Brand,-{Stealth[length=1.8mm]}] (\x,\h)--(\x,0.08);
  \draw[Brand] (0.45,0.20)--(3.75,1.12);
  \draw[<->] (0,-0.42)--node[below] {\(l\)} (4,-0.42);
  \draw[<->] (4.25,0)--node[right] {\(q_0\)} (4.25,1.12);
  \node[below] at (0,0) {\(A\)};
  \node[below left] at (4,0) {\(B\)};
  \figtitle{(a) 固定端承受三角形分布荷载}
\end{tikzpicture}
\end{figure}
\end{problemcard}

\begin{solutioncard}
\textbf{答案：}\answer{\(F_{By}=\dfrac12q_0l\) 向上，\(M_B=\dfrac16q_0l^2\) 顺时针。}

\textbf{解析：}三角形荷载等效为一个向下合力
\[
  W=\frac12q_0l .
\]
合力作用点在距荷载强度最大端 \(B\) 为 \(l/3\) 的位置。由整体竖向力平衡，
\[
  \sum F_y=0:\qquad F_{By}-W=0
  \quad\Rightarrow\quad
  F_{By}=\frac12q_0l .
\]
对 \(B\) 点取矩，三角形荷载对 \(B\) 的力臂为 \(l/3\)，因此固定端反力矩大小为
\[
  M_B=W\cdot\frac{l}{3}
  =\frac12q_0l\cdot\frac{l}{3}
  =\frac16q_0l^2 .
\]
该反力矩方向与外荷载对 \(B\) 的转动趋势相反，按题解记为顺时针。
\end{solutioncard}

\begin{problemcard}{2-1(b) 集中力与局部均布荷载}
悬臂梁 \(AB\) 左端 \(A\) 固定，全长 \(3\,\mathrm{m}\)。右端 \(B\) 作用 \(30\,\mathrm{kN}\) 向下集中力，末端 \(CB\) 长 \(1\,\mathrm{m}\)，承受 \(15\,\mathrm{kN/m}\) 均布荷载。求固定端反力。

\begin{figure}[H]
\centering
\begin{tikzpicture}[scale=0.95,>=Stealth]
  \draw[member] (0,0)--(4.5,0);
  \draw[ground] (0,-0.50)--(0,0.50);
  \foreach \y in {-0.40,-0.22,-0.04,0.14,0.32}
    \draw[ground] (0,\y)--(-0.18,\y-0.10);
  \foreach \x in {3.15,3.45,3.75,4.05,4.35}
    \draw[Brand,-{Stealth[length=1.8mm]}] (\x,0.62)--(\x,0.08);
  \draw[Brand] (3.05,0.62)--(4.45,0.62);
  \draw[Brand,-{Stealth[length=2mm]}] (4.5,1.10)--(4.5,0.08);
  \node[above] at (3.75,0.70) {\(15\,\mathrm{kN/m}\)};
  \node[right] at (4.5,1.00) {\(30\,\mathrm{kN}\)};
  \draw[<->] (0,-0.42)--node[below] {\(3\,\mathrm{m}\)} (4.5,-0.42);
  \draw[<->] (3.0,-0.78)--node[below] {\(1\,\mathrm{m}\)} (4.5,-0.78);
  \node[below] at (0,0) {\(A\)};
  \node[below] at (3.0,0) {\(C\)};
  \node[below] at (4.5,0) {\(B\)};
  \figtitle{(b) 固定端梁}
\end{tikzpicture}
\end{figure}
\end{problemcard}

\begin{solutioncard}
\textbf{答案：}\answer{\(F_{Ay}=45\,\mathrm{kN}\) 向上，\(M_A=127.5\,\mathrm{kN\cdot m}\) 逆时针。}

\textbf{解析：}末端均布荷载的合力为
\[
  15\times 1=15\,\mathrm{kN},
\]
作用在 \(CB\) 中点，距 \(A\) 为 \(2.5\,\mathrm{m}\)。竖向平衡得
\[
  F_{Ay}-15\times1-30=0
  \quad\Rightarrow\quad
  F_{Ay}=45\,\mathrm{kN}.
\]
对 \(A\) 点取矩，外荷载使梁产生顺时针转动，固定端反力矩取相反方向：
\[
  M_A=15\times1\times2.5+30\times3
  =37.5+90
  =127.5\,\mathrm{kN\cdot m}.
\]
\end{solutioncard}

\begin{problemcard}{2-1(c) 带外伸段的简支梁}
梁上 \(A\) 端作用 \(30\,\mathrm{kN\cdot m}\) 力偶，\(C\)、\(B\) 为竖向支座；\(AC=1\,\mathrm{m}\)，\(AB=4\,\mathrm{m}\)，集中力 \(40\,\mathrm{kN}\) 作用于 \(D\)，且 \(DB=1.5\,\mathrm{m}\)。求 \(C\)、\(B\) 处支反力。

\begin{figure}[H]
\centering
\begin{tikzpicture}[scale=0.95,>=Stealth]
  \draw[member] (0,0)--(5.0,0);
  \roller{1.25}{0}
  \roller{5.0}{0}
  \draw[Brand,-{Stealth[length=2mm]}] (3.15,1.0)--(3.15,0.08);
  \node[right] at (3.15,0.86) {\(40\,\mathrm{kN}\)};
  \draw[Brand,-{Stealth[length=2mm]}] (0.20,0.55) arc[start angle=65,end angle=300,radius=0.36];
  \node[above left] at (0.15,0.55) {\(30\,\mathrm{kN\cdot m}\)};
  \draw[<->] (0,-0.55)--node[below] {\(1\,\mathrm{m}\)} (1.25,-0.55);
  \draw[<->] (0,-0.90)--node[below] {\(4\,\mathrm{m}\)} (5,-0.90);
  \draw[<->] (3.15,-0.55)--node[below] {\(1.5\,\mathrm{m}\)} (5,-0.55);
  \node[below] at (0,0) {\(A\)};
  \node[below] at (1.25,0) {\(C\)};
  \node[below] at (3.15,0) {\(D\)};
  \node[below] at (5,0) {\(B\)};
  \figtitle{(c) 对 \(C\) 点取矩可直接求 \(B\) 反力}
\end{tikzpicture}
\end{figure}
\end{problemcard}

\begin{solutioncard}
\textbf{答案：}\answer{\(F_B=10\,\mathrm{kN}\) 向上，\(F_C=30\,\mathrm{kN}\) 向上。}

\textbf{解析：}取整体隔离体。对 \(C\) 点取矩时，\(F_C\) 的力矩为零，\(B\) 到 \(C\) 的距离为 \(3\,\mathrm{m}\)，\(40\,\mathrm{kN}\) 到 \(C\) 的距离为 \(1.5\,\mathrm{m}\)。按题解的符号约定列式：
\[
  F_B\times3+30-40\times1.5=0 .
\]
故
\[
  F_B=\frac{40\times1.5-30}{3}=10\,\mathrm{kN}.
\]
再由竖向力平衡
\[
  F_C+F_B-40=0
  \quad\Rightarrow\quad
  F_C=40-10=30\,\mathrm{kN}.
\]
\end{solutioncard}

\begin{problemcard}{2-1(d) 均布荷载与集中力偶}
简支梁 \(AB\) 全长 \(10\,\mathrm{m}\)，全梁受 \(0.2\,\mathrm{kN/m}\) 均布荷载；\(C\) 点距 \(B\) 为 \(2\,\mathrm{m}\)，在 \(C\) 点作用 \(4\,\mathrm{kN\cdot m}\) 力偶。求两端支反力。

\begin{figure}[H]
\centering
\begin{tikzpicture}[scale=0.88,>=Stealth]
  \draw[member] (0,0)--(6,0);
  \roller{0}{0}
  \roller{6}{0}
  \foreach \x in {0.25,0.65,...,5.85}
    \draw[Brand,-{Stealth[length=1.7mm]}] (\x,0.62)--(\x,0.08);
  \draw[Brand] (0.15,0.62)--(5.95,0.62);
  \draw[Brand,-{Stealth[length=2mm]}] (4.80,0.85) arc[start angle=120,end angle=-160,radius=0.36];
  \node[above] at (2.25,0.80) {\(0.2\,\mathrm{kN/m}\)};
  \node[above right] at (4.80,0.98) {\(4\,\mathrm{kN\cdot m}\)};
  \draw[<->] (0,-0.50)--node[below] {\(10\,\mathrm{m}\)} (6,-0.50);
  \draw[<->] (4.8,-0.86)--node[below] {\(2\,\mathrm{m}\)} (6,-0.86);
  \node[below] at (0,0) {\(A\)};
  \node[below] at (4.8,0) {\(C\)};
  \node[below] at (6,0) {\(B\)};
  \figtitle{(d) 全跨均布荷载}
\end{tikzpicture}
\end{figure}
\end{problemcard}

\begin{solutioncard}
\textbf{答案：}\answer{\(F_B=1.4\,\mathrm{kN}\) 向上，\(F_A=0.6\,\mathrm{kN}\) 向上。}

\textbf{解析：}全梁均布荷载合力为
\[
  W=0.2\times10=2\,\mathrm{kN},
\]
作用在跨中，距 \(A\) 为 \(5\,\mathrm{m}\)。对 \(A\) 点取矩：
\[
  F_B\times10-4-0.2\times10\times5=0 .
\]
于是
\[
  F_B=\frac{4+10}{10}=1.4\,\mathrm{kN}.
\]
竖向力平衡给出
\[
  F_A+F_B-2=0
  \quad\Rightarrow\quad
  F_A=2-1.4=0.6\,\mathrm{kN}.
\]
该题的关键是把集中力偶直接作为力矩项加入方程，不需要给它再乘距离。
\end{solutioncard}

\begin{problemcard}{2-1(e) 两端力偶与集中力}
简支梁 \(AB\) 长 \(3\,\mathrm{m}\)，\(A\) 端作用 \(6\,\mathrm{kN\cdot m}\) 力偶，\(B\) 端作用 \(20\,\mathrm{kN\cdot m}\) 力偶；\(C\) 点距 \(B\) 为 \(1\,\mathrm{m}\)，有 \(20\,\mathrm{kN}\) 集中力向下。求支反力。

\begin{figure}[H]
\centering
\begin{tikzpicture}[scale=0.95,>=Stealth]
  \draw[member] (0,0)--(4.2,0);
  \roller{0}{0}
  \roller{4.2}{0}
  \draw[Brand,-{Stealth[length=2mm]}] (0.25,0.62) arc[start angle=60,end angle=300,radius=0.35];
  \draw[Brand,-{Stealth[length=2mm]}] (4.00,0.62) arc[start angle=120,end angle=-160,radius=0.35];
  \draw[Brand,-{Stealth[length=2mm]}] (2.8,0.95)--(2.8,0.08);
  \node[above left] at (0.30,0.72) {\(6\,\mathrm{kN\cdot m}\)};
  \node[above right] at (3.95,0.98) {\(20\,\mathrm{kN\cdot m}\)};
  \node[left] at (2.75,1.05) {\(20\,\mathrm{kN}\)};
  \draw[<->] (0,-0.52)--node[below] {\(3\,\mathrm{m}\)} (4.2,-0.52);
  \draw[<->] (2.8,-0.88)--node[below] {\(1\,\mathrm{m}\)} (4.2,-0.88);
  \node[below] at (0,0) {\(A\)};
  \node[below] at (2.8,0) {\(C\)};
  \node[below] at (4.2,0) {\(B\)};
  \figtitle{(e) 集中力偶不乘距离}
\end{tikzpicture}
\end{figure}
\end{problemcard}

\begin{solutioncard}
\textbf{答案：}\answer{\(F_B=22\,\mathrm{kN}\) 向上；若向上为正，\(F_A=-2\,\mathrm{kN}\)，即 \(A\) 处实际向下，大小为 \(2\,\mathrm{kN}\)。}

\textbf{解析：}集中力 \(20\,\mathrm{kN}\) 作用点距 \(A\) 为 \(2\,\mathrm{m}\)。对 \(A\) 点取矩，按题解中力偶方向写为
\[
  F_B\times3-6-20\times2-20=0 .
\]
解得
\[
  F_B=\frac{6+40+20}{3}=22\,\mathrm{kN}.
\]
由竖向平衡
\[
  F_A+F_B-20=0
  \quad\Rightarrow\quad
  F_A=20-22=-2\,\mathrm{kN}.
\]
因此若以向上为正，\(F_A=-2\,\mathrm{kN}\)；题解图按反力实际方向标注时，可记为 \(A\) 处 \(2\,\mathrm{kN}\) 向下。若只列反力大小，应同时注明方向，避免把符号含义混淆。
\end{solutioncard}

\begin{problemcard}{2-1(f) 伸臂梁端部均布荷载}
梁 \(AB\) 总长 \(5a\)，在 \(A\) 与 \(C\) 处支承，\(CB=a\)。伸臂段 \(CB\) 上作用强度为 \(q\) 的均布荷载。求 \(A\)、\(C\) 两处支反力。

\begin{figure}[H]
\centering
\begin{tikzpicture}[scale=0.88,>=Stealth]
  \draw[member] (0,0)--(6,0);
  \roller{0}{0}
  \roller{4.8}{0}
  \foreach \x in {5.0,5.25,5.50,5.75}
    \draw[Brand,-{Stealth[length=1.7mm]}] (\x,0.62)--(\x,0.08);
  \draw[Brand] (4.85,0.62)--(6.0,0.62);
  \node[above] at (5.4,0.72) {\(q\)};
  \draw[<->] (0,-0.52)--node[below] {\(5a\)} (6,-0.52);
  \draw[<->] (4.8,-0.88)--node[below] {\(a\)} (6,-0.88);
  \node[below] at (0,0) {\(A\)};
  \node[below] at (4.8,0) {\(C\)};
  \node[below] at (6,0) {\(B\)};
  \figtitle{(f) 右端伸臂段受均布荷载}
\end{tikzpicture}
\end{figure}
\end{problemcard}

\begin{solutioncard}
\textbf{答案：}\answer{\(F_C=\dfrac98qa\) 向上；若向上为正，\(F_A=-\dfrac18qa\)，即 \(A\) 处实际向下，大小为 \(\dfrac18qa\)。}

\textbf{解析：}伸臂段上的均布荷载合力为 \(qa\)，作用点在 \(CB\) 中点。由题图尺寸，\(AC=4a\)，合力作用点距 \(A\) 为
\[
  4a+\frac{a}{2}=4.5a .
\]
对 \(A\) 点取矩：
\[
  F_C\cdot4a-qa\cdot4.5a=0
  \quad\Rightarrow\quad
  F_C=\frac{4.5}{4}qa=\frac98qa .
\]
再由竖向力平衡
\[
  F_A+F_C-qa=0
  \quad\Rightarrow\quad
  F_A=qa-\frac98qa=-\frac18qa .
\]
负号说明 \(A\) 处反力实际方向向下，这也是伸臂梁端部荷载常见的“压起”现象。
\end{solutioncard}

\begin{problemcard}{2-1(g) 局部均布荷载悬臂梁}
悬臂梁 \(AB\) 左端 \(A\) 固定，全长 \(2.5\,\mathrm{m}\)，右端 \(CB\) 长 \(1\,\mathrm{m}\)，承受 \(5\,\mathrm{kN/m}\) 均布荷载。求固定端反力。

\begin{figure}[H]
\centering
\begin{tikzpicture}[scale=0.95,>=Stealth]
  \draw[member] (0,0)--(4.5,0);
  \draw[ground] (0,-0.50)--(0,0.50);
  \foreach \y in {-0.40,-0.22,-0.04,0.14,0.32}
    \draw[ground] (0,\y)--(-0.18,\y-0.10);
  \foreach \x in {2.75,3.10,3.45,3.80,4.15}
    \draw[Brand,-{Stealth[length=1.7mm]}] (\x,0.62)--(\x,0.08);
  \draw[Brand] (2.7,0.62)--(4.5,0.62);
  \node[above] at (3.55,0.72) {\(5\,\mathrm{kN/m}\)};
  \draw[<->] (0,-0.50)--node[below] {\(2.5\,\mathrm{m}\)} (4.5,-0.50);
  \draw[<->] (2.7,-0.86)--node[below] {\(1\,\mathrm{m}\)} (4.5,-0.86);
  \node[below] at (0,0) {\(A\)};
  \node[below] at (2.7,0) {\(C\)};
  \node[below] at (4.5,0) {\(B\)};
  \figtitle{(g) 合力距固定端 \(2\,\mathrm{m}\)}
\end{tikzpicture}
\end{figure}
\end{problemcard}

\begin{solutioncard}
\textbf{答案：}\answer{\(F_{Ay}=5\,\mathrm{kN}\) 向上，\(M_A=10\,\mathrm{kN\cdot m}\) 逆时针。}

\textbf{解析：}均布荷载合力为
\[
  W=5\times1=5\,\mathrm{kN}.
\]
荷载段长 \(1\,\mathrm{m}\)，其中点距 \(B\) 为 \(0.5\,\mathrm{m}\)，因此距固定端 \(A\) 为 \(2.5-0.5=2\,\mathrm{m}\)。竖向平衡得
\[
  F_{Ay}=5\,\mathrm{kN}.
\]
对 \(A\) 点取矩：
\[
  M_A=5\times2=10\,\mathrm{kN\cdot m}.
\]
反力矩方向与荷载产生的顺时针转动相反，故为逆时针。
\end{solutioncard}

\begin{problemcard}{2-1(h) 集中力与集中力偶悬臂梁}
悬臂梁 \(AB\) 左端固定，全长 \(2\,\mathrm{m}\)。\(C\) 点距 \(A\) 为 \(1\,\mathrm{m}\)，在 \(C\) 点作用 \(15\,\mathrm{kN}\) 向下集中力和 \(10\,\mathrm{kN\cdot m}\) 力偶。求固定端反力。

\begin{figure}[H]
\centering
\begin{tikzpicture}[scale=0.95,>=Stealth]
  \draw[member] (0,0)--(4.4,0);
  \draw[ground] (0,-0.50)--(0,0.50);
  \foreach \y in {-0.40,-0.22,-0.04,0.14,0.32}
    \draw[ground] (0,\y)--(-0.18,\y-0.10);
  \draw[Brand,-{Stealth[length=2mm]}] (2.2,1.05)--(2.2,0.08);
  \draw[Brand,-{Stealth[length=2mm]}] (2.45,0.80) arc[start angle=120,end angle=-160,radius=0.36];
  \node[above left] at (2.2,1.02) {\(15\,\mathrm{kN}\)};
  \node[above right] at (2.35,0.80) {\(10\,\mathrm{kN\cdot m}\)};
  \draw[<->] (0,-0.50)--node[below] {\(1\,\mathrm{m}\)} (2.2,-0.50);
  \draw[<->] (0,-0.86)--node[below] {\(2\,\mathrm{m}\)} (4.4,-0.86);
  \node[below] at (0,0) {\(A\)};
  \node[below] at (2.2,0) {\(C\)};
  \node[below] at (4.4,0) {\(B\)};
  \figtitle{(h) 集中力与集中力偶叠加}
\end{tikzpicture}
\end{figure}
\end{problemcard}

\begin{solutioncard}
\textbf{答案：}\answer{\(F_{Ay}=15\,\mathrm{kN}\) 向上，\(M_A=25\,\mathrm{kN\cdot m}\) 逆时针。}

\textbf{解析：}竖向外力只有 \(15\,\mathrm{kN}\) 向下，因此
\[
  F_{Ay}=15\,\mathrm{kN}.
\]
对 \(A\) 点取矩时，集中力贡献 \(15\times1=15\,\mathrm{kN\cdot m}\)，题给集中力偶贡献 \(10\,\mathrm{kN\cdot m}\)。所以固定端反力矩大小为
\[
  M_A=10+15\times1=25\,\mathrm{kN\cdot m}.
\]
方向为抵抗外荷载转动的逆时针方向。
\end{solutioncard}

\begin{problemcard}{2-1(i) 对称均布荷载简支梁}
简支梁 \(AB\) 全长 \(5a\)。中部 \(CD\) 长 \(3a\)，承受强度为 \(q\) 的均布荷载；\(AC=a\)，\(DB=a\)。求 \(A\)、\(B\) 两处支反力。

\begin{figure}[H]
\centering
\begin{tikzpicture}[scale=0.85,>=Stealth]
  \draw[member] (0,0)--(6,0);
  \roller{0}{0}
  \roller{6}{0}
  \foreach \x in {1.25,1.65,...,4.75}
    \draw[Brand,-{Stealth[length=1.7mm]}] (\x,0.62)--(\x,0.08);
  \draw[Brand] (1.2,0.62)--(4.8,0.62);
  \node[above] at (3,0.72) {\(q\)};
  \draw[<->] (0,-0.50)--node[below] {\(5a\)} (6,-0.50);
  \draw[<->] (1.2,-0.86)--node[below] {\(3a\)} (4.8,-0.86);
  \draw[<->] (0,-1.18)--node[below] {\(a\)} (1.2,-1.18);
  \node[below] at (0,0) {\(A\)};
  \node[below] at (1.2,0) {\(C\)};
  \node[below] at (4.8,0) {\(D\)};
  \node[below] at (6,0) {\(B\)};
  \figtitle{(i) 荷载关于跨中对称}
\end{tikzpicture}
\end{figure}
\end{problemcard}

\begin{solutioncard}
\textbf{答案：}\answer{\(F_A=F_B=1.5qa\) 向上。}

\textbf{解析：}均布荷载合力为
\[
  W=q\cdot3a=3qa,
\]
由于荷载段位于跨中，合力作用点距 \(A\) 为 \(2.5a\)。对 \(A\) 点取矩：
\[
  F_B\cdot5a-3qa\cdot2.5a=0,
\]
所以
\[
  F_B=\frac{3qa\cdot2.5a}{5a}=1.5qa.
\]
再由竖向平衡
\[
  F_A+F_B-3qa=0
  \quad\Rightarrow\quad
  F_A=1.5qa.
\]
也可直接利用对称性得到两端反力相等，各承担总荷载的一半。
\end{solutioncard}

\begin{problemcard}{2-1(j) 两个集中力偶作用的简支梁}
简支梁 \(AB\) 长 \(3a\)，\(C\) 点距 \(A\) 为 \(a\)。\(C\) 点作用力偶 \(M_e\)，\(B\) 点作用力偶 \(2M_e\)，方向如图。求两端竖向支反力。

\begin{figure}[H]
\centering
\begin{tikzpicture}[scale=0.9,>=Stealth]
  \draw[member] (0,0)--(5.4,0);
  \roller{0}{0}
  \roller{5.4}{0}
  \draw[Brand,-{Stealth[length=2mm]}] (1.75,0.70) arc[start angle=140,end angle=-40,radius=0.43];
  \draw[Brand,-{Stealth[length=2mm]}] (5.05,0.70) arc[start angle=140,end angle=320,radius=0.43];
  \node[above] at (1.8,0.78) {\(M_e\)};
  \node[above] at (5.0,0.78) {\(2M_e\)};
  \draw[<->] (0,-0.50)--node[below] {\(a\)} (1.8,-0.50);
  \draw[<->] (0,-0.86)--node[below] {\(3a\)} (5.4,-0.86);
  \node[below] at (0,0) {\(A\)};
  \node[below] at (1.8,0) {\(C\)};
  \node[below] at (5.4,0) {\(B\)};
  \figtitle{(j) 仅有力偶时两支反力方向相反}
\end{tikzpicture}
\end{figure}
\end{problemcard}

\begin{solutioncard}
\textbf{答案：}\answer{\(F_B=\dfrac{M_e}{3a}\) 向上，\(F_A=\dfrac{M_e}{3a}\) 向下。}

\textbf{解析：}全结构没有竖向外力，只有两个外力偶，因此两端竖向反力必然大小相等、方向相反。对 \(A\) 点取矩，按题解方向列式：
\[
  F_B\cdot3a+M_e-2M_e=0 .
\]
解得
\[
  F_B=\frac{M_e}{3a}.
\]
由竖向平衡
\[
  F_A+F_B=0
  \quad\Rightarrow\quad
  F_A=-\frac{M_e}{3a}.
\]
负号表示若假定 \(F_A\) 向上，则实际方向向下，大小为 \(\frac{M_e}{3a}\)。
\end{solutioncard}

\begin{problemcard}{2-1(k) 带中间铰的刚架}
图示刚架左柱 \(AD\) 承受水平均布荷载 \(q\)，总高度为 \(2a\)，上横梁中点 \(C\) 为铰。底部 \(A\) 与右侧 \(B\) 为支座，尺寸如图。求 \(A\)、\(B\) 处支座反力。

\begin{figure}[H]
\centering
\begin{tikzpicture}[scale=0.88,>=Stealth]
  \coordinate (A) at (0,0);
  \coordinate (D) at (0,3.2);
  \coordinate (C) at (2.0,3.2);
  \coordinate (E) at (4.0,3.2);
  \coordinate (B) at (4.0,1.6);
  \draw[member] (A)--(D)--(C)--(E)--(B);
  \node[pin] at (C) {};
  \roller{0}{0}
  \wallroller{4.0}{1.6}
  \foreach \y in {0.35,0.70,...,2.95}
    \draw[Brand,-{Stealth[length=1.7mm]}] (-0.55,\y)--(-0.05,\y);
  \draw[Brand] (-0.55,0.20)--(-0.55,3.10);
  \node[left] at (-0.60,1.65) {\(q\)};
  \draw[<->] (0,-0.58)--node[below] {\(a\)} (2,-0.58);
  \draw[<->] (2,-0.58)--node[below] {\(a\)} (4,-0.58);
  \draw[<->] (4.55,0)--node[right] {\(a\)} (4.55,1.6);
  \draw[<->] (4.55,1.6)--node[right] {\(a\)} (4.55,3.2);
  \node[below] at (A) {\(A\)};
  \node[above left] at (D) {\(D\)};
  \node[above] at (C) {\(C\)};
  \node[right] at (B) {\(B\)};
  \figtitle{(k) 取整体与 \(BC\) 部分隔离体}
\end{tikzpicture}
\end{figure}
\end{problemcard}

\begin{solutioncard}
\textbf{答案：}\answer{\(F_{By}=\dfrac23qa\)，\(F_{Bx}=\dfrac23qa\)，\(F_{Ay}=\dfrac23qa\)，\(F_{Ax}=\dfrac43qa\)。方向按图示反力方向标注。}

\textbf{解析：}左柱上的水平均布荷载合力为
\[
  q\cdot2a=2qa,
\]
作用在柱高的中点，距 \(A\) 的竖向距离为 \(a\)。取整体隔离体对 \(A\) 点取矩：
\[
  F_{By}\cdot2a+F_{Bx}\cdot a=q\cdot2a\cdot a .
\]
再取 \(BC\) 部分隔离体，对铰 \(C\) 点取矩，有
\[
  F_{By}\cdot a=F_{Bx}\cdot a,
\]
即
\[
  F_{By}=F_{Bx}.
\]
联立可得
\[
  3aF_{By}=2qa^2
  \quad\Rightarrow\quad
  F_{By}=F_{Bx}=\frac23qa .
\]
最后由整体平衡得到 \(A\) 处反力：
\[
  \sum F_y=0:\quad F_{Ay}=F_{By}=\frac23qa,
\]
\[
  \sum F_x=0:\quad F_{Ax}+F_{Bx}=2qa
  \quad\Rightarrow\quad
  F_{Ax}=2qa-\frac23qa=\frac43qa .
\]
\end{solutioncard}

\begin{problemcard}{2-1(l) 门式刚架局部均布荷载}
门式刚架宽 \(10\,\mathrm{m}\)，高 \(6\,\mathrm{m}\)，顶部中点 \(C\) 为铰。右半跨顶部受 \(2\,\mathrm{kN/m}\) 均布荷载，荷载长度 \(5\,\mathrm{m}\)。求两支座反力。

\begin{figure}[H]
\centering
\begin{tikzpicture}[scale=0.62,>=Stealth]
  \coordinate (A) at (0,0);
  \coordinate (D) at (0,6);
  \coordinate (C) at (5,6);
  \coordinate (E) at (10,6);
  \coordinate (B) at (10,0);
  \draw[member] (A)--(D)--(C)--(E)--(B);
  \node[pin] at (C) {};
  \wallroller{0}{0}
  \wallroller{10}{0}
  \foreach \x in {5.6,6.4,...,9.6}
    \draw[Brand,-{Stealth[length=1.8mm]}] (\x,7.05)--(\x,6.08);
  \draw[Brand] (5.0,7.05)--(10,7.05);
  \node[above] at (7.5,7.15) {\(2\,\mathrm{kN/m}\)};
  \draw[<->] (0,-1.0)--node[below] {\(5\,\mathrm{m}\)} (5,-1.0);
  \draw[<->] (5,-1.0)--node[below] {\(5\,\mathrm{m}\)} (10,-1.0);
  \draw[<->] (10.9,0)--node[right] {\(6\,\mathrm{m}\)} (10.9,6);
  \node[below] at (A) {\(A\)};
  \node[below] at (B) {\(B\)};
  \node[above] at (C) {\(C\)};
  \figtitle{(l) 右半跨顶部均布荷载}
\end{tikzpicture}
\end{figure}
\end{problemcard}

\begin{solutioncard}
\textbf{答案：}\answer{\(F_{By}=7.5\,\mathrm{kN}\)，\(F_{Bx}=\dfrac{25}{12}\,\mathrm{kN}\)，\(F_{Ay}=2.5\,\mathrm{kN}\)，\(F_{Ax}=\dfrac{25}{12}\,\mathrm{kN}\)。方向按题解图的支座反力方向。}

\textbf{解析：}右半跨均布荷载的合力为
\[
  W=2\times5=10\,\mathrm{kN},
\]
作用于右半跨中点，距 \(A\) 为 \(7.5\,\mathrm{m}\)。取整体对 \(A\) 点取矩：
\[
  F_{By}\times10=2\times5\times7.5,
\]
故
\[
  F_{By}=7.5\,\mathrm{kN}.
\]
再取 \(BC\) 部分隔离体对铰 \(C\) 点取矩。\(B\) 到 \(C\) 的水平距离为 \(5\,\mathrm{m}\)，竖向距离为 \(6\,\mathrm{m}\)，右半跨荷载合力到 \(C\) 的距离为 \(2.5\,\mathrm{m}\)，于是
\[
  F_{By}\times5=F_{Bx}\times6+2\times5\times2.5 .
\]
代入 \(F_{By}=7.5\,\mathrm{kN}\)，得
\[
  F_{Bx}=\frac{7.5\times5-25}{6}
  =\frac{12.5}{6}
  =\frac{25}{12}\,\mathrm{kN}.
\]
竖向平衡：
\[
  F_{Ay}+F_{By}=10
  \quad\Rightarrow\quad
  F_{Ay}=2.5\,\mathrm{kN}.
\]
水平平衡中没有外部水平荷载，因此两侧水平反力大小相等、方向相反，按题解方向均记大小为
\[
  F_{Ax}=F_{Bx}=\frac{25}{12}\,\mathrm{kN}.
\]
\end{solutioncard}

\begin{problemcard}{2-1(m) 矩形刚架侧向均布荷载}
矩形刚架宽 \(l\)、高 \(h\)。左柱 \(AB\) 受向右的水平均布荷载 \(q\)，底部 \(A\)、\(D\) 为支座。求 \(A\)、\(D\) 处支座反力。

\begin{figure}[H]
\centering
\begin{tikzpicture}[scale=0.9,>=Stealth]
  \coordinate (A) at (0,0);
  \coordinate (B) at (0,3.2);
  \coordinate (C) at (4.2,3.2);
  \coordinate (D) at (4.2,0);
  \draw[member] (A)--(B)--(C)--(D);
  \roller{0}{0}
  \wallroller{4.2}{0}
  \foreach \y in {0.45,0.85,...,3.05}
    \draw[Brand,-{Stealth[length=1.7mm]}] (-0.55,\y)--(-0.04,\y);
  \draw[Brand] (-0.55,0.25)--(-0.55,3.15);
  \node[left] at (-0.62,1.65) {\(q\)};
  \draw[<->] (0,-0.62)--node[below] {\(l\)} (4.2,-0.62);
  \draw[<->] (4.75,0)--node[right] {\(h\)} (4.75,3.2);
  \node[below] at (A) {\(A\)};
  \node[above left] at (B) {\(B\)};
  \node[above right] at (C) {\(C\)};
  \node[below] at (D) {\(D\)};
  \figtitle{(m) 侧向均布荷载刚架}
\end{tikzpicture}
\end{figure}
\end{problemcard}

\begin{solutioncard}
\textbf{答案：}\answer{\(F_{Dy}=\dfrac{qh^2}{2l}\) 向上，\(F_{Ay}=\dfrac{qh^2}{2l}\) 向下，\(F_{Dx}=qh\) 向左。}

\textbf{解析：}左柱水平均布荷载的合力为
\[
  W=qh,
\]
作用于柱高的中点，距 \(A\) 的竖向距离为 \(h/2\)。取整体对 \(A\) 点取矩：
\[
  F_{Dy}\cdot l=q h\cdot\frac{h}{2},
\]
故
\[
  F_{Dy}=\frac{qh^2}{2l}.
\]
竖向外荷载为零，所以两端竖向支反力大小相等、方向相反：
\[
  F_{Ay}+F_{Dy}=0
  \quad\Rightarrow\quad
  F_{Ay}=-\frac{qh^2}{2l}.
\]
因此 \(A\) 处实际竖向反力向下，大小为 \(\frac{qh^2}{2l}\)。水平方向平衡给出
\[
  F_{Dx}=qh,
\]
方向与水平分布荷载相反，即向左。
\end{solutioncard}

\chapter{材料力学回顾：剪力图与弯矩图}

\section*{作图规则速记}
\addcontentsline{toc}{section}{作图规则速记}

\begin{notecard}
材料力学回顾部分的重点是把支座反力、剪力图和弯矩图连成一套检查链。作图时可按以下顺序：
\begin{enumerate}
  \item 先用整体平衡求支座反力和固定端反力矩。
  \item 从左到右画剪力图：集中力使剪力图突跳，均布荷载使剪力图为斜线，线性分布荷载使剪力图为二次曲线。
  \item 由 \(dM/dx=F_Q\) 画弯矩图：剪力为常数时弯矩为斜线，剪力为一次函数时弯矩为二次曲线，剪力为二次函数时弯矩为三次曲线。
  \item 集中力偶只使弯矩图突跳，不改变剪力图。
  \item 下列图中弯矩正负以讲义常用的“受拉侧”说明为准；数值标在图上，便于和原题答案逐项核对。
\end{enumerate}
\end{notecard}

\section{一、画剪力图和弯矩图}

\begin{problemcard}{3-1 画图题：前 6 个基本梁}
画出图示各梁的剪力图 \(Q\) 和弯矩图 \(M\)。本题的六个结构与第 2 次作业中求支座反力的前六个结构相同；本章在已求反力的基础上继续作内力图。

\begin{figure}[H]
\centering
\begin{tikzpicture}[scale=0.74,>=Stealth]
  % (a)
  \begin{scope}[shift={(0,5.2)}]
    \draw[member] (0,0)--(3.0,0);
    \draw[ground] (3,-0.45)--(3,0.45);
    \foreach \y in {-0.34,-0.16,0.02,0.20,0.38} \draw[ground] (3,\y)--(3.15,\y+0.08);
    \foreach \x/\h in {0.45/0.18,0.9/0.32,1.35/0.46,1.8/0.60,2.25/0.74,2.7/0.88}
      \draw[Brand,-{Stealth[length=1.5mm]}] (\x,\h)--(\x,0.07);
    \draw[Brand] (0.35,0.18)--(2.85,0.92);
    \node[below] at (0,0) {\(A\)};
    \node[below left] at (3,0) {\(B\)};
    \node at (1.5,-0.70) {(a)};
  \end{scope}
  % (b)
  \begin{scope}[shift={(5.2,5.2)}]
    \draw[member] (0,0)--(3.4,0);
    \draw[ground] (0,-0.43)--(0,0.43);
    \foreach \y in {-0.32,-0.14,0.04,0.22} \draw[ground] (0,\y)--(-0.15,\y-0.08);
    \foreach \x in {2.45,2.7,2.95,3.2} \draw[Brand,-{Stealth[length=1.5mm]}] (\x,0.56)--(\x,0.07);
    \draw[Brand] (2.35,0.56)--(3.35,0.56);
    \draw[Brand,-{Stealth[length=1.8mm]}] (3.4,0.92)--(3.4,0.07);
    \node[above] at (2.85,0.65) {\scriptsize \(15\,\mathrm{kN/m}\)};
    \node[right] at (3.4,0.80) {\scriptsize \(30\,\mathrm{kN}\)};
    \node[below] at (0,0) {\(A\)};
    \node[below] at (2.35,0) {\(C\)};
    \node[below] at (3.4,0) {\(B\)};
    \node at (1.7,-0.70) {(b)};
  \end{scope}
  % (c)
  \begin{scope}[shift={(0,2.55)}]
    \draw[member] (0,0)--(4.0,0);
    \roller{1.0}{0}
    \roller{4.0}{0}
    \draw[Brand,-{Stealth[length=1.8mm]}] (2.5,0.90)--(2.5,0.07);
    \draw[Brand,-{Stealth[length=1.8mm]}] (0.18,0.48) arc[start angle=65,end angle=300,radius=0.30];
    \node[above left] at (0.18,0.55) {\scriptsize \(30\,\mathrm{kN\cdot m}\)};
    \node[right] at (2.5,0.75) {\scriptsize \(40\,\mathrm{kN}\)};
    \node[below] at (0,0) {\(A\)};
    \node[below] at (1.0,0) {\(C\)};
    \node[below] at (2.5,0) {\(D\)};
    \node[below] at (4.0,0) {\(B\)};
    \node at (2.0,-0.70) {(c)};
  \end{scope}
  % (d)
  \begin{scope}[shift={(5.2,2.55)}]
    \draw[member] (0,0)--(4.0,0);
    \roller{0}{0}
    \roller{4.0}{0}
    \foreach \x in {0.25,0.55,...,3.85} \draw[Brand,-{Stealth[length=1.4mm]}] (\x,0.56)--(\x,0.07);
    \draw[Brand] (0.15,0.56)--(3.95,0.56);
    \draw[Brand,-{Stealth[length=1.8mm]}] (3.20,0.76) arc[start angle=120,end angle=-160,radius=0.28];
    \node[above] at (1.75,0.66) {\scriptsize \(0.2\,\mathrm{kN/m}\)};
    \node[above right] at (3.2,0.82) {\scriptsize \(4\,\mathrm{kN\cdot m}\)};
    \node[below] at (0,0) {\(A\)};
    \node[below] at (3.2,0) {\(C\)};
    \node[below] at (4.0,0) {\(B\)};
    \node at (2.0,-0.70) {(d)};
  \end{scope}
  % (e)
  \begin{scope}[shift={(0,0)}]
    \draw[member] (0,0)--(4.0,0);
    \roller{0}{0}
    \roller{4.0}{0}
    \draw[Brand,-{Stealth[length=1.8mm]}] (0.20,0.50) arc[start angle=60,end angle=300,radius=0.28];
    \draw[Brand,-{Stealth[length=1.8mm]}] (3.80,0.50) arc[start angle=120,end angle=-160,radius=0.28];
    \draw[Brand,-{Stealth[length=1.8mm]}] (2.65,0.90)--(2.65,0.07);
    \node[above left] at (0.28,0.58) {\scriptsize \(6\,\mathrm{kN\cdot m}\)};
    \node[above right] at (3.55,0.62) {\scriptsize \(20\,\mathrm{kN\cdot m}\)};
    \node[above] at (2.65,0.92) {\scriptsize \(20\,\mathrm{kN}\)};
    \node[below] at (0,0) {\(A\)};
    \node[below] at (2.65,0) {\(C\)};
    \node[below] at (4.0,0) {\(B\)};
    \node at (2.0,-0.70) {(e)};
  \end{scope}
  % (f)
  \begin{scope}[shift={(5.2,0)}]
    \draw[member] (0,0)--(4.2,0);
    \roller{0}{0}
    \roller{3.35}{0}
    \foreach \x in {3.55,3.75,3.95,4.15} \draw[Brand,-{Stealth[length=1.4mm]}] (\x,0.56)--(\x,0.07);
    \draw[Brand] (3.4,0.56)--(4.2,0.56);
    \node[above] at (3.8,0.66) {\scriptsize \(q\)};
    \node[below] at (0,0) {\(A\)};
    \node[below] at (3.35,0) {\(C\)};
    \node[below] at (4.2,0) {\(B\)};
    \node at (2.1,-0.70) {(f)};
  \end{scope}
  \figtitle{第 3 次作业第一部分题图：结构示意}
\end{tikzpicture}
\end{figure}
\end{problemcard}

\begin{solutioncard}
\textbf{(a) 三角形荷载悬臂梁}

由整体平衡已得
\[
  F_{By}=\frac12q_0l,\qquad
  M_B=-\frac16q_0l^2\quad(\text{上侧受拉}).
\]
若从自由端 \(A\) 起取坐标 \(x\)，荷载集度为 \(q(x)=q_0x/l\)，则
\[
  Q(x)=-\int_0^x\frac{q_0s}{l}\,ds=-\frac{q_0x^2}{2l},\qquad
  M(x)=\int_0^xQ(s)\,ds=-\frac{q_0x^3}{6l}.
\]
故 \(Q_A=M_A=0\)，固定端 \(B\) 处 \(Q_B=-\frac12q_0l\)，\(M_B=-\frac16q_0l^2\)。

\begin{figure}[H]
\centering
\begin{tikzpicture}[scale=0.95,>=Stealth]
  \begin{scope}
    \draw[member] (0,0)--(3.4,0);
    \draw[Brand,thick] (0,0) .. controls (1.1,-0.08) and (2.2,-0.42) .. (3.4,-1.0);
    \draw (3.4,0)--(3.4,-1.0);
    \node[right] at (3.4,-1.0) {\(\frac12q_0l\)};
    \node[below] at (1.7,-1.12) {\(Q\) 图};
  \end{scope}
  \begin{scope}[shift={(5.0,0)}]
    \draw[member] (0,0)--(3.4,0);
    \draw[Brand,thick] (0,0) .. controls (1.0,-0.04) and (2.3,-0.25) .. (3.4,-1.0);
    \draw (3.4,0)--(3.4,-1.0);
    \node[right] at (3.4,-1.0) {\(\frac16q_0l^2\)};
    \node[below] at (1.7,-1.12) {\(M\) 图};
  \end{scope}
\end{tikzpicture}
\end{figure}
\end{solutioncard}

\begin{solutioncard}
\textbf{(b) 末端集中力与局部均布荷载悬臂梁}

反力为
\[
  F_{Ay}=15\times1+30=45\,\mathrm{kN},\qquad
  M_A=-15\times1\times2.5-30\times3=-127.5\,\mathrm{kN\cdot m}.
\]
弯矩图上侧受拉。关键值为 \(M_A=127.5\,\mathrm{kN\cdot m}\)，在 \(C\) 处剩余荷载对截面的力矩为 \(37.5\,\mathrm{kN\cdot m}\)，到 \(B\) 端为零；剪力图在 \(AC\) 段为 \(45\,\mathrm{kN}\)，在 \(CB\) 段线性降至 \(30\,\mathrm{kN}\)，再由端部集中力突降为零。

\begin{figure}[H]
\centering
\begin{tikzpicture}[scale=0.88,>=Stealth]
  \begin{scope}
    \draw[member] (0,0)--(4.2,0);
    \draw[Brand,thick] (0,-1.15)--(3.0,-0.34)--(4.2,0);
    \draw (0,0)--(0,-1.15);
    \draw (3.0,0)--(3.0,-0.34);
    \node[left] at (0,-1.15) {\(127.5\)};
    \node[below] at (3.0,-0.34) {\(37.5\)};
    \node[below] at (2.1,-1.36) {\(M\) 图 \((\mathrm{kN\cdot m})\)};
  \end{scope}
  \begin{scope}[shift={(5.5,0)}]
    \draw[member] (0,0)--(4.2,0);
    \draw[Brand,thick] (0,0.85)--(3.0,0.85)--(4.2,0.55);
    \draw (0,0)--(0,0.85);
    \draw (4.2,0.55)--(4.2,0);
    \node[above] at (0,0.85) {\(45\)};
    \node[above] at (3.0,0.85) {\(45\)};
    \node[above] at (4.2,0.55) {\(30\)};
    \node[below] at (2.1,-0.28) {\(Q\) 图 \((\mathrm{kN})\)};
  \end{scope}
\end{tikzpicture}
\end{figure}
\end{solutioncard}

\begin{solutioncard}
\textbf{(c) 带外伸段的简支梁}

由第 2 章平衡结果
\[
  F_B=10\,\mathrm{kN},\qquad F_C=30\,\mathrm{kN}.
\]
剪力在 \(C\) 点由零突跳到 \(+30\,\mathrm{kN}\)，在 \(D\) 点受 \(40\,\mathrm{kN}\) 集中力后变为 \(-10\,\mathrm{kN}\)，到 \(B\) 点回零。弯矩图受 \(A\) 端力偶影响，在 \(AC\) 段保持 \(30\,\mathrm{kN\cdot m}\)，随后线性变化，\(D\) 附近出现 \(-15\,\mathrm{kN\cdot m}\) 控制值，至 \(B\) 点为零。

\begin{figure}[H]
\centering
\begin{tikzpicture}[scale=0.88,>=Stealth]
  \begin{scope}
    \draw[member] (0,0)--(4.0,0);
    \draw[Brand,thick] (0,0.75)--(1.0,0.75)--(2.5,-0.38)--(4.0,0);
    \draw[densely dashed] (1.0,0.75)--(4.0,0);
    \draw (0,0)--(0,0.75);
    \node[above] at (0.55,0.75) {\(30\)};
    \node[below] at (2.5,-0.38) {\(15\)};
    \node[below] at (2.0,-0.78) {\(M\) 图 \((\mathrm{kN\cdot m})\)};
  \end{scope}
  \begin{scope}[shift={(5.3,0)}]
    \draw[member] (0,0)--(4.0,0);
    \draw[Brand,thick] (1.0,0.70)--(2.5,0.70)--(2.5,-0.32)--(4.0,-0.32);
    \draw (1.0,0)--(1.0,0.70);
    \draw (4.0,-0.32)--(4.0,0);
    \node[above] at (1.7,0.70) {\(30\)};
    \node[below] at (3.3,-0.32) {\(10\)};
    \node[below] at (2.0,-0.78) {\(Q\) 图 \((\mathrm{kN})\)};
  \end{scope}
\end{tikzpicture}
\end{figure}
\end{solutioncard}

\begin{solutioncard}
\textbf{(d) 全跨均布荷载与集中力偶}

整体平衡得
\[
  F_B\times10-4-0.2\times10\times5=0,\qquad
  F_B=1.4\,\mathrm{kN},\quad F_A=0.6\,\mathrm{kN}.
\]
集中力偶使 \(C\) 点两侧弯矩突变。由 \(C\) 左侧截开取 \(AC\) 段，
\[
  M_C^L=0.6\times8-0.2\times8\times4=-1.6\,\mathrm{kN\cdot m},
\]
由 \(C\) 结点力矩平衡得
\[
  M_C^R=2.4\,\mathrm{kN\cdot m}.
\]
剪力图为斜线，从 \(A\) 点 \(+0.6\,\mathrm{kN}\) 线性降到 \(B\) 点支座前的 \(-1.4\,\mathrm{kN}\)。

\begin{figure}[H]
\centering
\begin{tikzpicture}[scale=0.88,>=Stealth]
  \begin{scope}
    \draw[member] (0,0)--(4.8,0);
    \draw[Brand,thick] (0,0) parabola bend (2.4,-0.55) (3.85,-0.34);
    \draw[Brand,thick] (3.85,-0.34)--(3.85,0.60);
    \draw[Brand,thick] (3.85,0.60) parabola bend (4.25,0.32) (4.8,0);
    \node[above] at (3.88,0.62) {\(2.4\)};
    \node[below] at (3.55,-0.35) {\(1.6\)};
    \node[below] at (2.4,-0.86) {\(M\) 图 \((\mathrm{kN\cdot m})\)};
  \end{scope}
  \begin{scope}[shift={(6.0,0)}]
    \draw[member] (0,0)--(4.8,0);
    \draw[Brand,thick] (0,0.45)--(4.8,-0.80);
    \draw (0,0)--(0,0.45);
    \draw (4.8,-0.80)--(4.8,0);
    \node[above] at (0,0.45) {\(0.6\)};
    \node[below] at (4.8,-0.80) {\(1.4\)};
    \node[below] at (2.4,-1.05) {\(Q\) 图 \((\mathrm{kN})\)};
  \end{scope}
\end{tikzpicture}
\end{figure}
\end{solutioncard}

\begin{solutioncard}
\textbf{(e) 两端集中力偶与集中力}

整体平衡为
\[
  F_B\times3-6-20\times2-20=0,
\]
故
\[
  F_B=22\,\mathrm{kN}.
\]
若竖向向上为正，则
\[
  F_A+F_B-20=0\quad\Rightarrow\quad F_A=-2\,\mathrm{kN},
\]
即 \(A\) 处实际反力向下，大小为 \(2\,\mathrm{kN}\)。剪力图在 \(AC\) 段为 \(-2\,\mathrm{kN}\)，过集中力后为 \(-22\,\mathrm{kN}\)，到 \(B\) 点回零；弯矩图在集中力偶处突跳，并以图中控制值 \(6,2,20\) 标注。

\begin{figure}[H]
\centering
\begin{tikzpicture}[scale=0.88,>=Stealth]
  \begin{scope}
    \draw[member] (0,0)--(4.2,0);
    \draw[Brand,thick] (0,-0.18)--(2.8,-0.06)--(4.2,0.92);
    \draw (0,0)--(0,-0.18);
    \draw (4.2,0)--(4.2,0.92);
    \node[left] at (0,-0.18) {\(6\)};
    \node[below] at (2.8,-0.06) {\(2\)};
    \node[right] at (4.2,0.92) {\(20\)};
    \node[below] at (2.1,-0.70) {\(M\) 图 \((\mathrm{kN\cdot m})\)};
  \end{scope}
  \begin{scope}[shift={(5.5,0)}]
    \draw[member] (0,0)--(4.2,0);
    \draw[Brand,thick] (0,-0.18)--(2.8,-0.18)--(2.8,-0.82)--(4.2,-0.82);
    \draw (0,0)--(0,-0.18);
    \draw (4.2,-0.82)--(4.2,0);
    \node[below] at (1.4,-0.18) {\(2\)};
    \node[below] at (3.5,-0.82) {\(22\)};
    \node[below] at (2.1,-1.05) {\(Q\) 图 \((\mathrm{kN})\)};
  \end{scope}
\end{tikzpicture}
\end{figure}
\end{solutioncard}

\begin{solutioncard}
\textbf{(f) 伸臂梁端部均布荷载}

由整体平衡
\[
  F_C\cdot4a-qa\cdot4.5a=0,
\]
得
\[
  F_C=\frac98qa,\qquad F_A=-\frac18qa.
\]
所以 \(A\) 处反力实际向下。剪力图在 \(AC\) 段为 \(-\frac18qa\)，到 \(C\) 点由支反力突跳为 \(qa\)，在 \(CB\) 受均布荷载线性降至零。弯矩图在 \(C\) 点达到控制值 \(\frac12qa^2\)，伸臂段为二次曲线收至 \(B\) 点零弯矩。

\begin{figure}[H]
\centering
\begin{tikzpicture}[scale=0.88,>=Stealth]
  \begin{scope}
    \draw[member] (0,0)--(5.0,0);
    \draw[Brand,thick] (0,0)--(4.0,0.70);
    \draw[Brand,thick] (4.0,0.70) .. controls (4.35,0.55) and (4.75,0.15) .. (5.0,0);
    \node[above] at (4.0,0.70) {\(\frac12qa^2\)};
    \node[below] at (2.5,-0.35) {\(M\) 图};
  \end{scope}
  \begin{scope}[shift={(6.0,0)}]
    \draw[member] (0,0)--(5.0,0);
    \draw[Brand,thick] (0,-0.25)--(4.0,-0.25)--(4.0,0.78)--(5.0,0);
    \draw (0,0)--(0,-0.25);
    \node[below] at (2.0,-0.25) {\(\frac18qa\)};
    \node[above] at (4.12,0.78) {\(qa\)};
    \node[below] at (2.5,-0.70) {\(Q\) 图};
  \end{scope}
\end{tikzpicture}
\end{figure}
\end{solutioncard}

\begin{notecard}
\textbf{讲解笔记补充：第一部分的分段式与作图检查}

\begin{enumerate}
  \item 对 (a)，若从自由端 \(A\) 向固定端 \(B\) 取坐标 \(x\)，线性荷载为 \(q(x)=q_0x/l\)，因此
  \[
    Q(x)=-\frac{q_0}{2l}x^2,\qquad
    M(x)=-\frac{q_0}{6l}x^3.
  \]
  剪力图为二次曲线，弯矩图为三次曲线，且自由端均为零。

  \item 对 (b)，从固定端 \(A\) 起取 \(x\)，讲解笔记中给出的分段式为
  \[
    Q(x)=
    \begin{cases}
      45, & 0\le x<2,\\
      75-15x, & 2\le x\le3,
    \end{cases}
  \]
  \[
    M(x)=
    \begin{cases}
      45x-127.5, & 0\le x<2,\\
      -157.5+75x-7.5x^2, & 2\le x\le3.
    \end{cases}
  \]
  因而 \(A\) 端控制值为 \(127.5\,\mathrm{kN\cdot m}\)，\(C\) 处为 \(37.5\,\mathrm{kN\cdot m}\)，到 \(B\) 端收为零。

  \item 对 (c)，先由平衡得 \(F_C=30\,\mathrm{kN}\)、\(F_B=10\,\mathrm{kN}\)。由于 \(AC\) 段无竖向剪力，弯矩保持 \(30\,\mathrm{kN\cdot m}\)；\(D\) 点受 \(40\,\mathrm{kN}\) 集中力后，剪力从 \(+30\,\mathrm{kN}\) 变为 \(-10\,\mathrm{kN}\)，弯矩图折线斜率随之改变，故在 \(D\) 附近出现 \(-15\,\mathrm{kN\cdot m}\) 控制值。

  \item 对 (d)，均布荷载使剪力图为一条斜线。弯矩图在 \(C\) 点因集中力偶突跳；按受拉侧标注时，\(C\) 左侧控制值为 \(1.6\,\mathrm{kN\cdot m}\)，\(C\) 右侧控制值为 \(2.4\,\mathrm{kN\cdot m}\)。这正是“集中力偶只改变弯矩图、不改变剪力图”的典型检查点。

  \item 对 (e)，由 \(AC\) 段隔离体可得 \(Q_{AC}=-2\,\mathrm{kN}\)、\(M_C=2\,\mathrm{kN\cdot m}\)。\(AC\) 段无分布荷载，故 \(M\) 图为斜线；过集中力后 \(Q\) 突跳到 \(-22\,\mathrm{kN}\)，再结合两端力偶得到 \(M\) 图的 \(6,2,20\) 三个控制值。

  \item 对 (f)，从右端伸臂 \(CB\) 隔离体可直接得到 \(M_C=\frac12qa^2\)。\(AC\) 段剪力为 \(-\frac18qa\)，到 \(C\) 点由支座反力突跳为 \(qa\)，随后在 \(CB\) 段受均布荷载线性降为零；因此 \(C\) 点既是剪力突跳点，也是弯矩控制点。
\end{enumerate}
\end{notecard}

\section{二、利用 \texorpdfstring{\(q\)、\(Q\)、\(M\)}{q、Q、M} 的微分关系作图}

\begin{notecard}
本组题的作图检查关系为
\[
  \frac{dQ}{dx}=-q(x),\qquad \frac{dM}{dx}=Q(x).
\]
因此：无分布荷载区间剪力图为水平线、弯矩图为斜线；均布荷载区间剪力图为斜线、弯矩图为抛物线；集中力使剪力突跳；集中力偶使弯矩突跳。
\end{notecard}

\begin{problemcard}{3-2 由微分关系作剪力图和弯矩图}
试利用荷载集度、剪力和弯矩之间的微分关系，作下列各梁的剪力图和弯矩图。

\begin{figure}[H]
\centering
\begin{tikzpicture}[scale=0.68,>=Stealth]
  % a
  \begin{scope}[shift={(0,5.3)}]
    \draw[member] (0,0)--(3.2,0);
    \draw[ground] (0,-0.45)--(0,0.45);
    \foreach \y in {-0.34,-0.16,0.02,0.20,0.38} \draw[ground] (0,\y)--(-0.15,\y-0.08);
    \foreach \x in {2.15,2.40,2.65,2.90,3.15} \draw[Brand,-{Stealth[length=1.4mm]}] (\x,0.55)--(\x,0.07);
    \draw[Brand] (2.05,0.55)--(3.2,0.55);
    \node[above] at (2.65,0.66) {\scriptsize \(5\,\mathrm{kN/m}\)};
    \node[below] at (0,0) {\(A\)};
    \node[below] at (2.05,0) {\(C\)};
    \node[below] at (3.2,0) {\(B\)};
    \node at (1.6,-0.65) {(a)};
  \end{scope}
  % b
  \begin{scope}[shift={(5.3,5.3)}]
    \draw[member] (0,0)--(3.2,0);
    \draw[ground] (0,-0.45)--(0,0.45);
    \foreach \y in {-0.34,-0.16,0.02,0.20,0.38} \draw[ground] (0,\y)--(-0.15,\y-0.08);
    \draw[Brand,-{Stealth[length=1.8mm]}] (1.6,0.95)--(1.6,0.07);
    \draw[Brand,-{Stealth[length=1.8mm]}] (1.88,0.70) arc[start angle=120,end angle=-160,radius=0.28];
    \node[above left] at (1.6,0.96) {\scriptsize \(15\,\mathrm{kN}\)};
    \node[above right] at (1.82,0.75) {\scriptsize \(10\,\mathrm{kN\cdot m}\)};
    \node[below] at (0,0) {\(A\)};
    \node[below] at (1.6,0) {\(C\)};
    \node[below] at (3.2,0) {\(B\)};
    \node at (1.6,-0.65) {(b)};
  \end{scope}
  % c
  \begin{scope}[shift={(0,2.65)}]
    \draw[member] (0,0)--(4.2,0);
    \roller{0}{0}
    \roller{4.2}{0}
    \foreach \x in {0.95,1.25,...,3.25} \draw[Brand,-{Stealth[length=1.35mm]}] (\x,0.55)--(\x,0.07);
    \draw[Brand] (0.85,0.55)--(3.35,0.55);
    \node[above] at (2.1,0.66) {\scriptsize \(q\)};
    \node[below] at (0,0) {\(A\)};
    \node[below] at (0.85,0) {\(C\)};
    \node[below] at (3.35,0) {\(D\)};
    \node[below] at (4.2,0) {\(B\)};
    \node at (2.1,-0.65) {(c)};
  \end{scope}
  % d
  \begin{scope}[shift={(5.3,2.65)}]
    \draw[member] (0,0)--(4.2,0);
    \roller{0}{0}
    \roller{4.2}{0}
    \draw[Brand,-{Stealth[length=1.8mm]}] (1.45,0.65) arc[start angle=140,end angle=-40,radius=0.32];
    \draw[Brand,-{Stealth[length=1.8mm]}] (3.95,0.66) arc[start angle=140,end angle=320,radius=0.32];
    \node[above] at (1.45,0.72) {\scriptsize \(M_e\)};
    \node[above] at (3.95,0.72) {\scriptsize \(2M_e\)};
    \node[below] at (0,0) {\(A\)};
    \node[below] at (1.4,0) {\(C\)};
    \node[below] at (4.2,0) {\(B\)};
    \node at (2.1,-0.65) {(d)};
  \end{scope}
  % e
  \begin{scope}[shift={(0,0)}]
    \draw[member] (0,0)--(4.2,0);
    \roller{0}{0}
    \roller{4.2}{0}
    \foreach \x in {0.3,0.6,...,3.9} \draw[Brand,-{Stealth[length=1.35mm]}] (\x,0.55)--(\x,0.07);
    \draw[Brand] (0.15,0.55)--(4.05,0.55);
    \draw[Brand,-{Stealth[length=1.8mm]}] (0.25,0.78) arc[start angle=60,end angle=300,radius=0.28];
    \draw[Brand,-{Stealth[length=1.8mm]}] (3.95,0.78) arc[start angle=120,end angle=-160,radius=0.28];
    \node[above] at (2.1,0.66) {\scriptsize \(4\,\mathrm{kN/m}\)};
    \node[above left] at (0.25,0.86) {\scriptsize \(4\,\mathrm{kN\cdot m}\)};
    \node[above right] at (3.85,0.86) {\scriptsize \(20\,\mathrm{kN\cdot m}\)};
    \node at (2.1,-0.65) {(e)};
  \end{scope}
  % f
  \begin{scope}[shift={(5.3,0)}]
    \draw[member] (0,0)--(4.2,0);
    \roller{0}{0}
    \roller{4.2}{0}
    \draw[Brand,-{Stealth[length=1.8mm]}] (1.05,0.85)--(1.05,0.07);
    \draw[Brand,-{Stealth[length=1.8mm]}] (2.1,-0.80)--(2.1,-0.07);
    \draw[Brand,-{Stealth[length=1.8mm]}] (3.15,0.85)--(3.15,0.07);
    \node[above] at (1.05,0.86) {\scriptsize \(F\)};
    \node[below] at (2.1,-0.86) {\scriptsize \(\frac{11}{8}F\)};
    \node[above] at (3.15,0.86) {\scriptsize \(F\)};
    \node[below] at (0,0) {\(A\)};
    \node[below] at (4.2,0) {\(B\)};
    \node at (2.1,-1.22) {(f)};
  \end{scope}
  \figtitle{第二部分题图 \((a)\)--\((f)\)}
\end{tikzpicture}
\end{figure}

\begin{figure}[H]
\centering
\begin{tikzpicture}[scale=0.76,>=Stealth]
  % g
  \begin{scope}
    \draw[member] (0,0)--(4.0,0);
    \roller{0}{0}
    \roller{4.0}{0}
    \foreach \x in {0.25,0.5,...,1.75} \draw[Brand,-{Stealth[length=1.35mm]}] (\x,0.55)--(\x,0.07);
    \draw[Brand] (0.15,0.55)--(1.85,0.55);
    \node[above] at (1.0,0.66) {\scriptsize \(2\,\mathrm{kN/m}\)};
    \node[below] at (0,0) {\(A\)};
    \node[below] at (2.0,0) {\(C\)};
    \node[below] at (4.0,0) {\(B\)};
    \node at (2.0,-0.65) {(g)};
  \end{scope}
  % h
  \begin{scope}[shift={(5.4,0)}]
    \draw[member] (0,0)--(5.0,0);
    \roller{0}{0}
    \roller{5.0}{0}
    \foreach \x in {0.25,0.55,...,4.75} \draw[Brand,-{Stealth[length=1.25mm]}] (\x,0.50)--(\x,0.07);
    \draw[Brand] (0.15,0.50)--(4.85,0.50);
    \draw[Brand,-{Stealth[length=1.8mm]}] (1.65,1.05)--(1.65,0.07);
    \draw[Brand,-{Stealth[length=1.8mm]}] (3.35,1.05)--(3.35,0.07);
    \node[above] at (1.05,0.61) {\scriptsize \(10\,\mathrm{kN/m}\)};
    \node[above] at (1.65,1.05) {\scriptsize \(250\,\mathrm{kN}\)};
    \node[above] at (3.35,1.05) {\scriptsize \(250\,\mathrm{kN}\)};
    \node[below] at (0,0) {\(A\)};
    \node[below] at (5.0,0) {\(B\)};
    \node at (2.5,-0.65) {(h)};
  \end{scope}
  \figtitle{第二部分题图 \((g)\)--\((h)\)}
\end{tikzpicture}
\end{figure}
\end{problemcard}

\begin{solutioncard}
\textbf{(a) 端部均布荷载悬臂梁}

整体平衡得
\[
  F_{Ay}=5\times1=5\,\mathrm{kN},\qquad
  M_A=-5\times1\times2=-10\,\mathrm{kN\cdot m}.
\]
在无荷载的 \(AC\) 段，剪力为常数、弯矩为斜线；在受均布荷载的 \(CB\) 段，剪力线性降为零，弯矩为二次曲线。关键值为 \(M_A=10\,\mathrm{kN\cdot m}\)、\(M_C=2.5\,\mathrm{kN\cdot m}\)、\(Q_A=Q_C=5\,\mathrm{kN}\)。

\begin{figure}[H]
\centering
\begin{tikzpicture}[scale=0.88,>=Stealth]
  \begin{scope}
    \draw[member] (0,0)--(4.2,0);
    \draw[Brand,thick] (0,1.0)--(2.8,0.25);
    \draw[Brand,thick] (2.8,0.25) .. controls (3.4,0.15) and (3.9,0.03) .. (4.2,0);
    \draw (0,0)--(0,1.0);
    \node[left] at (0,1.0) {\(10\)};
    \node[above] at (2.8,0.25) {\(2.5\)};
    \node[below] at (2.1,-0.35) {\(M\) 图};
  \end{scope}
  \begin{scope}[shift={(5.4,0)}]
    \draw[member] (0,0)--(4.2,0);
    \draw[Brand,thick] (0,0.75)--(2.8,0.75)--(4.2,0);
    \draw (0,0)--(0,0.75);
    \node[above] at (0,0.75) {\(5\)};
    \node[above] at (2.8,0.75) {\(5\)};
    \node[below] at (2.1,-0.35) {\(Q\) 图};
  \end{scope}
\end{tikzpicture}
\end{figure}
\end{solutioncard}

\begin{solutioncard}
\textbf{(b) 集中力与集中力偶悬臂梁}

整体平衡得
\[
  F_{Ay}=15\,\mathrm{kN},\qquad
  M_A=-10-15\times1=-25\,\mathrm{kN\cdot m}.
\]
\(AC\) 段剪力为 \(15\,\mathrm{kN}\)，弯矩线性变化；集中力偶使 \(C\) 点弯矩突跳，\(CB\) 段无剪力、无弯矩。

\begin{figure}[H]
\centering
\begin{tikzpicture}[scale=0.9,>=Stealth]
  \begin{scope}
    \draw[member] (0,0)--(4.0,0);
    \draw[Brand,thick] (0,1.0)--(2.0,0.42);
    \draw[Brand,thick] (2.0,0.42)--(2.0,0)--(4.0,0);
    \node[left] at (0,1.0) {\(25\)};
    \node[above] at (2.0,0.42) {\(10\)};
    \node[below] at (2.0,-0.35) {\(M\) 图};
  \end{scope}
  \begin{scope}[shift={(5.2,0)}]
    \draw[member] (0,0)--(4.0,0);
    \draw[Brand,thick] (0,0.75)--(2.0,0.75)--(2.0,0)--(4.0,0);
    \draw (0,0)--(0,0.75);
    \node[above] at (1.0,0.75) {\(15\)};
    \node[below] at (2.0,-0.35) {\(Q\) 图};
  \end{scope}
\end{tikzpicture}
\end{figure}
\end{solutioncard}

\begin{solutioncard}
\textbf{(c) 对称局部均布荷载简支梁}

由平衡方程
\[
  F_B\cdot5a-q\cdot3a\cdot2.5a=0
\]
得
\[
  F_A=F_B=\frac32qa.
\]
剪力图在 \(AC\) 段为 \(+\frac32qa\)，在 \(CD\) 段线性降至 \(-\frac32qa\)，在 \(DB\) 段保持负常数。弯矩图在 \(C,D\) 两点均为 \(\frac32qa^2\)，跨中最大为 \(\frac{21}{8}qa^2\)。

\begin{figure}[H]
\centering
\begin{tikzpicture}[scale=0.88,>=Stealth]
  \begin{scope}
    \draw[member] (0,0)--(5,0);
    \draw[Brand,thick] (0,0)--(1,0.55);
    \draw[Brand,thick] (1,0.55) .. controls (2.0,0.95) and (3.0,0.95) .. (4,0.55);
    \draw[Brand,thick] (4,0.55)--(5,0);
    \node[above] at (1,0.55) {\(\frac32qa^2\)};
    \node[above] at (2.5,0.95) {\(\frac{21}{8}qa^2\)};
    \node[above] at (4,0.55) {\(\frac32qa^2\)};
    \node[below] at (2.5,-0.35) {\(M\) 图};
  \end{scope}
  \begin{scope}[shift={(6.0,0)}]
    \draw[member] (0,0)--(5,0);
    \draw[Brand,thick] (0,0.70)--(1,0.70)--(4,-0.70)--(5,-0.70);
    \draw (0,0)--(0,0.70);
    \draw (5,-0.70)--(5,0);
    \node[above] at (0.5,0.70) {\(\frac32qa\)};
    \node[below] at (4.5,-0.70) {\(\frac32qa\)};
    \node[below] at (2.5,-1.0) {\(Q\) 图};
  \end{scope}
\end{tikzpicture}
\end{figure}
\end{solutioncard}

\begin{solutioncard}
\textbf{(d) 两个集中力偶作用的简支梁}

取整体对 \(A\) 点取矩：
\[
  F_B\cdot3a+M_e-2M_e=0,
\]
故
\[
  F_B=\frac{M_e}{3a},\qquad F_A=-\frac{M_e}{3a}.
\]
剪力图为全跨常值 \(-M_e/(3a)\)。弯矩图在 \(C\) 点因 \(M_e\) 突跳：\(C\) 左侧为 \(\frac13M_e\)，突跳后为 \(\frac43M_e\)，至 \(B\) 点为 \(2M_e\)。

\begin{figure}[H]
\centering
\begin{tikzpicture}[scale=0.9,>=Stealth]
  \begin{scope}
    \draw[member] (0,0)--(4.5,0);
    \draw[Brand,thick] (0,0)--(1.5,0.25)--(1.5,0.75)--(4.5,1.05);
    \draw (4.5,0)--(4.5,1.05);
    \node[above left,font=\scriptsize] at (1.45,0.27) {\(\frac13M_e\)};
    \node[above,font=\scriptsize] at (1.90,0.78) {\(\frac43M_e\)};
    \node[right] at (4.5,1.05) {\(2M_e\)};
    \node[below] at (2.25,-0.35) {\(M\) 图};
  \end{scope}
  \begin{scope}[shift={(5.8,0)}]
    \draw[member] (0,0)--(4.5,0);
    \draw[Brand,thick] (0,-0.45)--(4.5,-0.45);
    \draw (0,0)--(0,-0.45);
    \draw (4.5,-0.45)--(4.5,0);
    \node at (2.25,-0.23) {\(\frac{M_e}{3a}\)};
    \node[below] at (2.25,-0.88) {\(Q\) 图};
  \end{scope}
\end{tikzpicture}
\end{figure}
\end{solutioncard}

\begin{solutioncard}
\textbf{(e) 全跨均布荷载与两端力偶}

对 \(A\) 点取矩：
\[
  F_B\cdot4=4\times4\times2+4+20,
\]
得
\[
  F_B=14\,\mathrm{kN},\qquad F_A=2\,\mathrm{kN}.
\]
剪力图从 \(+2\,\mathrm{kN}\) 线性降至 \(-14\,\mathrm{kN}\)。弯矩图为抛物线，并受两端力偶影响，按答案图标出端部控制值 \(4\,\mathrm{kN\cdot m}\) 与 \(20\,\mathrm{kN\cdot m}\)，跨内控制值约 \(8\,\mathrm{kN\cdot m}\)。

\begin{figure}[H]
\centering
\begin{tikzpicture}[scale=0.88,>=Stealth]
  \begin{scope}
    \draw[member] (0,0)--(4.4,0);
    \draw[Brand,thick] (0,0.28) .. controls (1.4,0.10) and (3.0,0.50) .. (4.4,1.0);
    \draw (0,0)--(0,0.28);
    \draw (4.4,0)--(4.4,1.0);
    \node[left] at (0,0.28) {\(4\)};
    \node[above] at (2.2,0.42) {\(8\)};
    \node[right] at (4.4,1.0) {\(20\)};
    \node[below] at (2.2,-0.35) {\(M\) 图};
  \end{scope}
  \begin{scope}[shift={(5.7,0)}]
    \draw[member] (0,0)--(4.4,0);
    \draw[Brand,thick] (0,0.38)--(4.4,-0.85);
    \draw (0,0)--(0,0.38);
    \draw (4.4,-0.85)--(4.4,0);
    \node[above] at (0,0.38) {\(+2\)};
    \node[below] at (4.4,-0.85) {\(14\)};
    \node[below] at (2.2,-1.12) {\(Q\) 图};
  \end{scope}
\end{tikzpicture}
\end{figure}
\end{solutioncard}

\begin{solutioncard}
\textbf{(f) 三个集中力作用的简支梁}

整体平衡：
\[
  F_B\cdot4a+\frac{11}{8}F\cdot2a=F\cdot a+F\cdot3a,
\]
故
\[
  F_B=\frac{5}{16}F,\qquad F_A=\frac{5}{16}F.
\]
剪力图依次为
\[
  +\frac{5}{16}F,\quad -\frac{11}{16}F,\quad +\frac{11}{16}F,\quad -\frac{5}{16}F.
\]
弯矩图由这些常剪力段积分得到折线图，控制值为两侧 \(\frac{5}{16}Fa\)，中部 \(\frac{3}{8}Fa\)。

\begin{figure}[H]
\centering
\begin{tikzpicture}[scale=0.88,>=Stealth]
  \begin{scope}
    \draw[member] (0,0)--(4.8,0);
    \draw[Brand,thick] (0,0)--(1.2,-0.45)--(2.4,0.70)--(3.6,-0.45)--(4.8,0);
    \node[below] at (1.2,-0.45) {\(\frac{5}{16}Fa\)};
    \node[above] at (2.4,0.70) {\(\frac38Fa\)};
    \node[below] at (3.6,-0.45) {\(\frac{5}{16}Fa\)};
    \node[below] at (2.4,-0.85) {\(M\) 图};
  \end{scope}
  \begin{scope}[shift={(6.0,0)}]
    \draw[member] (0,0)--(4.8,0);
    \draw[Brand,thick] (0,0.40)--(1.2,0.40)--(1.2,-0.70)--(2.4,-0.70)--(2.4,0.70)--(3.6,0.70)--(3.6,-0.40)--(4.8,-0.40);
    \node[above] at (0.6,0.40) {\(\frac{5}{16}F\)};
    \node[below] at (1.8,-0.70) {\(\frac{11}{16}F\)};
    \node[above] at (3.0,0.70) {\(\frac{11}{16}F\)};
    \node[below] at (4.2,-0.40) {\(\frac{5}{16}F\)};
    \node[below] at (2.4,-1.02) {\(Q\) 图};
  \end{scope}
\end{tikzpicture}
\end{figure}
\end{solutioncard}

\begin{solutioncard}
\textbf{(g) 左半跨均布荷载简支梁}

对 \(A\) 点取矩：
\[
  F_B\cdot2=2\times1\times0.5,
\]
故
\[
  F_B=\frac12\,\mathrm{kN},\qquad F_A=\frac32\,\mathrm{kN}.
\]
在左侧 \(1\,\mathrm{m}\) 均布荷载段，
\[
  Q(x)=\frac32-2x,\qquad M(x)=\frac32x-x^2.
\]
所以 \(M_C=\frac12\,\mathrm{kN\cdot m}\)，最大弯矩在 \(x=0.75\,\mathrm{m}\) 处，为 \(\frac{9}{16}\,\mathrm{kN\cdot m}\)。

\begin{figure}[H]
\centering
\begin{tikzpicture}[scale=0.88,>=Stealth]
  \begin{scope}
    \draw[member] (0,0)--(4,0);
    \draw[Brand,thick] (0,0) .. controls (0.8,0.62) and (1.2,0.55) .. (2,0.45)--(4,0);
    \node[above] at (1.0,0.62) {\(\frac{9}{16}\)};
    \node[above] at (2,0.45) {\(\frac12\)};
    \node[below] at (2,-0.35) {\(M\) 图};
  \end{scope}
  \begin{scope}[shift={(5.2,0)}]
    \draw[member] (0,0)--(4,0);
    \draw[Brand,thick] (0,0.75)--(2,-0.25)--(4,-0.25);
    \draw (0,0)--(0,0.75);
    \draw (4,-0.25)--(4,0);
    \node[above] at (0,0.75) {\(\frac32\)};
    \node[below] at (3,-0.25) {\(\frac12\)};
    \node[below] at (2,-0.70) {\(Q\) 图};
  \end{scope}
\end{tikzpicture}
\end{figure}
\end{solutioncard}

\begin{solutioncard}
\textbf{(h) 对称集中力与全跨均布荷载}

取整体对 \(A\) 点取矩：
\[
  F_B\cdot6=10\times6\times3+250\times2+250\times4,
\]
得
\[
  F_B=280\,\mathrm{kN},\qquad F_A=280\,\mathrm{kN}.
\]
剪力图从 \(+280\,\mathrm{kN}\) 开始，均布荷载使其线性下降；在两个 \(250\,\mathrm{kN}\) 集中力处各突降一次。关键值依次为 \(280,260,10,-10,-260,-280\)。弯矩图关于跨中对称，在 \(x=2\,\mathrm{m}\) 和 \(x=4\,\mathrm{m}\) 处为 \(540\,\mathrm{kN\cdot m}\)，跨中最大为 \(545\,\mathrm{kN\cdot m}\)。

\begin{figure}[H]
\centering
\begin{tikzpicture}[scale=0.86,>=Stealth]
  \begin{scope}
    \draw[member] (0,0)--(5.4,0);
    \draw[Brand,thick] (0,0) .. controls (1.2,0.75) and (1.8,0.86) .. (2.0,0.82);
    \draw[Brand,thick] (2.0,0.82) .. controls (2.5,0.90) and (3.1,0.90) .. (3.4,0.82);
    \draw[Brand,thick] (3.4,0.82) .. controls (4.0,0.72) and (5.0,0.25) .. (5.4,0);
    \node[above,font=\scriptsize] at (1.85,0.82) {\(540\)};
    \node[above,font=\scriptsize] at (2.70,1.04) {\(545\)};
    \node[above,font=\scriptsize] at (3.55,0.82) {\(540\)};
    \node[below] at (2.7,-0.35) {\(M\) 图 \((\mathrm{kN\cdot m})\)};
  \end{scope}
  \begin{scope}[shift={(6.5,0)}]
    \draw[member] (0,0)--(5.4,0);
    \draw[Brand,thick] (0,0.90)--(1.8,0.78)--(1.8,0.04)--(2.7,0)--(3.6,-0.04)--(3.6,-0.78)--(5.4,-0.90);
    \draw (0,0)--(0,0.90);
    \draw (5.4,-0.90)--(5.4,0);
    \node[above] at (0,0.90) {\(280\)};
    \node[above] at (1.8,0.78) {\(260\)};
    \node[above] at (2.35,0.08) {\(10\)};
    \node[below] at (3.15,-0.08) {\(10\)};
    \node[below] at (3.6,-0.78) {\(260\)};
    \node[below] at (5.4,-0.90) {\(280\)};
    \node[below] at (2.7,-1.18) {\(Q\) 图 \((\mathrm{kN})\)};
  \end{scope}
\end{tikzpicture}
\end{figure}
\end{solutioncard}

\begin{notecard}
\textbf{讲解笔记补充：第二部分的微分关系检查}

\begin{enumerate}
  \item (a)、(b) 都是悬臂梁。无荷载段剪力为常数、弯矩为斜线；受均布荷载段剪力为斜线、弯矩为抛物线。(a) 的关键值是 \(M_A=10\,\mathrm{kN\cdot m}\)、\(M_C=2.5\,\mathrm{kN\cdot m}\)；(b) 的 \(CB\) 段无剪力、无弯矩，集中力偶只造成 \(C\) 处弯矩突跳。

  \item (c) 的要点是对称性与叠加。两端反力均为 \(\frac32qa\)，在 \(C,D\) 处弯矩均为 \(\frac32qa^2\)。中段均布荷载会在这个线性底图上叠加抛物线，因此跨中最大值为 \(\frac{21}{8}qa^2\)。

  \item (d) 没有横向分布荷载，剪力图全跨为常值 \(-M_e/(3a)\)。两个集中力偶使弯矩图在 \(C\) 与 \(B\) 处突跳；作图时先画由常剪力产生的斜线，再在力偶位置按力偶大小跳跃。

  \item (e) 全跨均布荷载使剪力线性下降，端部力偶只改弯矩图。由平衡得 \(F_A=2\,\mathrm{kN}\)、\(F_B=14\,\mathrm{kN}\)，弯矩图受两端力偶控制，端部标为 \(4\,\mathrm{kN\cdot m}\)、\(20\,\mathrm{kN\cdot m}\)，跨内约 \(8\,\mathrm{kN\cdot m}\) 为检查值。

  \item (f) 可利用对称性快速校核：两端反力均为 \(\frac{5}{16}F\)。剪力依次为 \(+\frac{5}{16}F,-\frac{11}{16}F,+\frac{11}{16}F,-\frac{5}{16}F\)，弯矩控制值为两侧 \(\frac{5}{16}Fa\)、中部 \(\frac38Fa\)。

  \item (g) 左半跨均布荷载下，先求 \(F_A=\frac32\,\mathrm{kN}\)、\(F_B=\frac12\,\mathrm{kN}\)。令 \(Q(x)=\frac32-2x=0\) 可得最大弯矩位置 \(x=0.75\,\mathrm{m}\)，最大值 \(\frac{9}{16}\,\mathrm{kN\cdot m}\)。

  \item (h) 结构与荷载关于跨中对称，故 \(F_A=F_B=280\,\mathrm{kN}\)。剪力关键值 \(280,260,10,-10,-260,-280\) 是逐段检查图形最可靠的顺序；弯矩图关于跨中对称，\(x=2\,\mathrm{m}\)、\(x=4\,\mathrm{m}\) 处为 \(540\,\mathrm{kN\cdot m}\)，跨中为 \(545\,\mathrm{kN\cdot m}\)。
\end{enumerate}
\end{notecard}

\section{三、具有中间铰的梁}

\begin{notecard}
带中间铰的梁作图时，先利用铰处弯矩为零把结构拆成两部分。若某一段为二力杆，则该段只有轴力，不产生剪力和弯矩；若铰的一侧还有支座和外荷载，则先对该侧隔离体取矩求铰力，再回到另一侧画内力图。
\end{notecard}

\begin{problemcard}{3-3 具有中间铰的梁}
试作下列具有中间铰的梁的剪力图和弯矩图。

\begin{figure}[H]
\centering
\begin{tikzpicture}[scale=0.88,>=Stealth]
  \begin{scope}
    \draw[ground] (0,-0.50)--(0,0.55);
    \foreach \y in {-0.40,-0.20,0,0.20,0.40}
      \draw[ground] (0,\y)--(-0.15,\y-0.10);
    \draw[member] (0,0)--(1.8,0);
    \node[smallpin] at (1.8,0) {};
    \draw[member] (1.8,0)--(5.4,0);
    \roller{3.6}{0}
    \foreach \x in {0.35,0.65,...,1.55}
      \draw[Brand,-{Stealth[length=1.5mm]}] (\x,0.62)--(\x,0.08);
    \draw[Brand] (0.22,0.62)--(1.65,0.62);
    \draw[Brand,-{Stealth[length=2mm]}] (5.4,0.95)--(5.4,0.08);
    \node[above] at (0.92,0.68) {\(q\)};
    \node[right] at (5.4,0.85) {\(F=qa\)};
    \node[left] at (0,0) {\(A\)};
    \node[below] at (1.8,0) {\(C\)};
    \node[below] at (3.6,0) {\(B\)};
    \node[below] at (5.4,0) {\(D\)};
    \draw[<->] (0,-0.62)--(1.8,-0.62) node[midway,below] {\(a\)};
    \draw[<->] (1.8,-0.62)--(3.6,-0.62) node[midway,below] {\(a\)};
    \draw[<->] (3.6,-0.62)--(5.4,-0.62) node[midway,below] {\(a\)};
    \node at (2.7,-1.15) {(a)};
  \end{scope}

  \begin{scope}[shift={(0,-3.1)}]
    \node[smallpin] at (0,0) {};
    \draw[ground] (-0.24,-0.25)--(0.24,-0.25);
    \draw[ground] (-0.18,-0.25)--(-0.28,-0.38);
    \draw[ground] (0,-0.25)--(-0.10,-0.38);
    \draw[ground] (0.18,-0.25)--(0.08,-0.38);
    \draw[member] (0,0)--(1.8,0);
    \node[smallpin] at (1.8,0) {};
    \draw[member] (1.8,0)--(5.4,0);
    \roller{3.6}{0}
    \roller{5.4}{0}
    \draw[Brand,-{Stealth[length=2mm]}] (1.8,0.95)--(1.8,0.08);
    \draw[Brand,-{Stealth[length=2mm]}] (5.15,0.52) arc[start angle=140,end angle=-210,radius=0.32];
    \node[above] at (1.8,0.95) {\(F=qa\)};
    \node[above right] at (5.35,0.56) {\(M_e=qa^2\)};
    \node[below] at (0,0) {\(A\)};
    \node[below] at (1.8,0) {\(C\)};
    \node[below] at (3.6,0) {\(B\)};
    \node[below] at (5.4,0) {\(D\)};
    \draw[<->] (0,-0.62)--(1.8,-0.62) node[midway,below] {\(a\)};
    \draw[<->] (1.8,-0.62)--(3.6,-0.62) node[midway,below] {\(a\)};
    \draw[<->] (3.6,-0.62)--(5.4,-0.62) node[midway,below] {\(a\)};
    \node at (2.7,-1.15) {(b)};
  \end{scope}
  \figtitle{第三部分题图：中间铰梁结构示意}
\end{tikzpicture}
\end{figure}
\end{problemcard}

\begin{solutioncard}
\textbf{(a) 固端梁与右侧伸臂段通过中间铰连接}

先取右侧 \(CD\) 部分隔离体。设铰 \(C\) 对右侧隔离体的竖向作用为 \(F_C\)，对 \(B\) 点取矩：
\[
  F_C\cdot a=qa\cdot a,
\]
故
\[
  F_C=qa.
\]
该力在右侧隔离体上向下，在左侧 \(AC\) 段上则向上。右侧竖向平衡还给出
\[
  F_B-F_C-qa=0\quad\Rightarrow\quad F_B=2qa.
\]

再取左侧 \(AC\) 段。铰力 \(F_C=qa\) 向上，均布荷载合力 \(qa\) 向下，故固定端竖向反力为零；对 \(A\) 点取矩得
\[
  M_A=qa\cdot a-\frac12qa\cdot a=\frac12qa^2,
\]
下侧受拉。若 \(x\) 从 \(A\) 向 \(C\) 量取，则
\[
  Q_{AC}(x)=-qx,\qquad
  M_{AC}(x)=\frac12qa^2-\frac12qx^2,
  \quad 0\le x\le a.
\]
于是 \(C\) 铰处 \(M_C=0\)。右侧 \(CD\) 段中，\(CB\) 段剪力为 \(-qa\)，到 \(B\) 点由支反力 \(2qa\) 突跳为 \(+qa\)，再在 \(D\) 点由端部集中力 \(qa\) 回零；弯矩图在 \(B\) 点达到控制值 \(qa^2\)，两端 \(C,D\) 均为零。

\begin{figure}[H]
\centering
\begin{tikzpicture}[scale=0.90,>=Stealth]
  \begin{scope}
    \draw[member] (0,0)--(5.4,0);
    \draw[Brand,thick] (0,0.62) parabola bend (1.8,0) (1.8,0);
    \draw[Brand,thick] (1.8,0)--(3.6,1.05)--(5.4,0);
    \node[above] at (0,0.62) {\(\frac12qa^2\)};
    \node[above] at (3.6,1.05) {\(qa^2\)};
    \node[below] at (2.7,-0.35) {\(M\) 图};
  \end{scope}
  \begin{scope}[shift={(0,-2.2)}]
    \draw[member] (0,0)--(5.4,0);
    \draw[Brand,thick] (0,0)--(1.8,-0.72)--(1.8,-0.72)--(3.6,-0.72)--(3.6,0.72)--(5.4,0.72);
    \draw[Brand,thick] (5.4,0.72)--(5.4,0);
    \node[below] at (1.55,-0.72) {\(qa\)};
    \node[above] at (4.5,0.72) {\(qa\)};
    \node[below] at (2.7,-1.05) {\(Q\) 图};
  \end{scope}
\end{tikzpicture}
\end{figure}
\end{solutioncard}

\begin{solutioncard}
\textbf{(b) 左段为二力杆的中间铰梁}

左段 \(AC\) 两端均为铰，且杆上没有横向荷载和力偶，因此 \(AC\) 是二力杆，只承受轴力：
\[
  Q_{AC}=0,\qquad M_{AC}=0.
\]

只需分析右侧 \(CBD\) 部分。集中力 \(F=qa\) 作用在 \(C\) 点，\(D\) 点有集中力偶 \(M_e=qa^2\)。由 \(C\) 点右侧隔离体可知，\(B\) 点竖向支反力承担该集中力，
\[
  F_B=qa,\qquad F_D=0,
\]
且 \(F_B\cdot a=M_e\) 正好平衡端部力偶。故 \(CB\) 段剪力为 \(-qa\)，过 \(B\) 点后剪力回零；弯矩图从 \(C\) 点零弯矩线性增至 \(B\) 点控制值 \(qa^2\)，在 \(BD\) 段保持常值，至 \(D\) 点受端部力偶突跳为零。

\begin{figure}[H]
\centering
\begin{tikzpicture}[scale=0.92,>=Stealth]
  \begin{scope}
    \draw[member] (0,0)--(5.4,0);
    \draw[Brand,thick] (1.8,0)--(3.6,0.88)--(5.4,0.88)--(5.4,0);
    \node[above] at (4.5,0.88) {\(qa^2\)};
    \node[below] at (2.7,-0.35) {\(M\) 图};
  \end{scope}
  \begin{scope}[shift={(0,-1.9)}]
    \draw[member] (0,0)--(5.4,0);
    \draw[Brand,thick] (1.8,0)--(1.8,-0.70)--(3.6,-0.70)--(3.6,0);
    \node[below] at (2.7,-0.70) {\(qa\)};
    \node[below] at (2.7,-1.05) {\(Q\) 图};
  \end{scope}
\end{tikzpicture}
\end{figure}
\end{solutioncard}

\chapter{静定梁和刚架}

\begin{notecard}
静定梁和刚架的内力图作图顺序为：先判断附属部分与基本部分，再从附属部分开始求支座反力；遇到铰结点时，铰两侧弯矩为零但剪力、轴力可以传递；遇到集中力偶时，弯矩图发生突跳，剪力图不变。本章题图均按原题结构重新绘制。
\end{notecard}

\section{一、基础判断与局部杆件内力}

\begin{problemcard}{2-1 截面弯矩符号判断}
图示结构中，截面 \(A\) 的弯矩以下侧受拉为正。判断该截面弯矩。
\[
  \text{选项：}\quad
  \mathrm{A.}\ M,\qquad
  \mathrm{B.}\ -M,\qquad
  \mathrm{C.}\ -2M,\qquad
  \mathrm{D.}\ 0.
\]
\begin{figure}[H]
\centering
\begin{tikzpicture}[scale=0.9,>=Stealth]
  \draw[member] (0,1.4)--(2.2,1.4)--(4.4,1.4);
  \draw[member] (2.2,1.4)--(2.2,0);
  \wallroller{0}{1.4}
  \roller{4.4}{1.4}
  \roller{0}{0}
  \node[smallpin] at (2.2,0) {};
  \draw[->,Brand,thick] (2.0,1.78) arc[start angle=160,end angle=-165,radius=0.34];
  \node[above] at (2.2,1.9) {\(M\)};
  \node[above] at (0.42,1.4) {\(A\)};
  \node[above] at (2.2,1.4) {\(C\)};
  \node[above] at (4.4,1.4) {\(B\)};
  \node[right] at (2.2,0) {\(D\)};
  \draw[<->] (2.35,0)--(2.35,1.4) node[midway,right] {\(h\)};
  \draw[<->] (0,-0.45)--(2.2,-0.45) node[midway,below] {\(l\)};
  \draw[<->] (2.2,-0.45)--(4.4,-0.45) node[midway,below] {\(l\)};
\end{tikzpicture}
\end{figure}
\end{problemcard}

\begin{solutioncard}
答案为 \answer{D，即 \(M_A=0\)}。

竖杆 \(CD\) 两端均为铰，且杆上没有横向荷载和集中力偶，因此 \(CD\) 为二力杆，只传递轴力。左侧 \(AC\) 段在 \(A\) 处为靠墙滚支，不能提供端弯矩；截面 \(A\) 取在支承附近时，该截面没有由外力形成的弯矩臂，故
\[
  M_A=0.
\]
集中力偶 \(M\) 只使 \(C\) 附近的弯矩图发生突跳，不会使无力臂的 \(A\) 截面产生弯矩。
\end{solutioncard}

\begin{problemcard}{2-2 杆端剪力判断}
图示结构中，求 \(BA\) 杆的剪力 \(F_{QBA}\)。
\begin{figure}[H]
\centering
\begin{tikzpicture}[scale=0.82,>=Stealth]
  \draw[member] (0,0)--(2,0)--(2,2.5)--(4,2.5)--(4,0)--(5.8,0);
  \draw[member] (5.8,0)--(5.8,-0.1);
  \roller{0}{0}
  \roller{2}{0}
  \roller{4}{0}
  \node[smallpin] at (2,2.5) {};
  \node[smallpin] at (4,2.5) {};
  \foreach \x in {4.25,4.55,...,5.65}
    \draw[Brand,-{Stealth[length=1.5mm]}] (\x,0.65)--(\x,0.08);
  \draw[Brand] (4.2,0.65)--(5.75,0.65);
  \node[above] at (5.0,0.65) {\(q\)};
  \node[below] at (0,0) {\(A\)};
  \node[below] at (2,0) {\(B\)};
  \node[left] at (2,2.5) {\(C\)};
  \node[right] at (4,2.5) {\(D\)};
  \node[below] at (4,0) {\(E\)};
  \node[below] at (5.8,0) {\(F\)};
  \draw[<->] (0,-0.5)--(2,-0.5) node[midway,below] {\(l\)};
  \draw[<->] (2,-0.5)--(4,-0.5) node[midway,below] {\(l\)};
  \draw[<->] (4,-0.5)--(5.8,-0.5) node[midway,below] {\(l\)};
  \draw[<->] (6.2,0)--(6.2,2.5) node[midway,right] {\(l\)};
\end{tikzpicture}
\end{figure}
\end{problemcard}

\begin{solutioncard}
\[
  F_{QBA}=0.
\]
\answer{答案：\(F_{QBA}=0\)。}
右侧 \(DE\) 杆两端为铰且杆上没有横向荷载，可视为二力杆。左侧 \(AB\) 杆上也没有分布荷载、集中力或力偶，且两端均通过铰支承传力。对 \(AB\) 杆隔离，杆内只可能有轴向力；横向平衡给出
\[
  \sum F_\perp=0\quad\Rightarrow\quad F_{QBA}=0.
\]
因此 \(BA\) 杆不画剪力图，弯矩图也为零。
\end{solutioncard}

\begin{problemcard}{2-3 刚架中 \(AB\) 杆内力}
图示刚架中，求 \(AB\) 杆的弯矩图与剪力图。
\begin{figure}[H]
\centering
\begin{tikzpicture}[scale=0.62,>=Stealth]
  \draw[member] (0,0)--(0,4)--(8,4)--(8,0);
  \draw[member] (4,0)--(4,4);
  \draw[member] (8,2)--(10,2)--(10,0);
  \roller{0}{0}
  \roller{4}{0}
  \roller{8}{0}
  \roller{10}{0}
  \node[smallpin] at (0,2) {};
  \node[smallpin] at (8,2) {};
  \foreach \y in {0.5,1.0,1.5}
    \draw[Brand,-{Stealth[length=1.5mm]}] (-0.85,\y)--(-0.05,\y);
  \draw[Brand] (-0.85,0.25)--(-0.85,1.75);
  \node[left] at (-0.85,1.0) {\(20\,\mathrm{kN/m}\)};
  \draw[Brand,->,thick] (0.35,4.35) arc[start angle=20,end angle=335,radius=0.34];
  \node[above] at (0.9,4.25) {\(10\,\mathrm{kN\cdot m}\)};
  \draw[Brand,->,thick] (0.35,2.0)--(1.5,2.0) node[right] {\(5\,\mathrm{kN\cdot m}\)};
  \draw[Brand,-{Stealth[length=2mm]}] (9,3.0)--(9,2.05) node[midway,right] {\(30\,\mathrm{kN}\)};
  \draw[<->] (0,-0.55)--(4,-0.55) node[midway,below] {\(4\,\mathrm{m}\)};
  \draw[<->] (4,-0.55)--(8,-0.55) node[midway,below] {\(8\,\mathrm{m}\)};
  \draw[<->] (8,-0.55)--(10,-0.55) node[midway,below] {\(2\,\mathrm{m}\)};
  \draw[<->] (10.55,0)--(10.55,2) node[midway,right] {\(4\,\mathrm{m}\)};
  \draw[<->] (10.55,2)--(10.55,4) node[midway,right] {\(4\,\mathrm{m}\)};
  \node[left] at (0,4) {\(A\)};
  \node[left] at (0,2) {\(B\)};
  \node[below] at (4,0) {\(C\)};
  \node[below] at (8,0) {\(D\)};
\end{tikzpicture}
\end{figure}
\end{problemcard}

\begin{solutioncard}
先求左侧水平分布荷载合力：
\[
  F_x=20\times4=80\,\mathrm{kN}.
\]
该合力作用在左下 \(4\,\mathrm{m}\) 荷载段的中点。对中间支承附近取矩，可得到支承传来的水平力为 \(80\,\mathrm{kN}\)。于是 \(AB\) 杆没有沿杆方向的分布荷载，其剪力为常值，弯矩为直线图。按原答案的控制值：
\[
  Q_{AB}=80\,\mathrm{kN},\qquad
  M_A=0,\quad M_B=320\,\mathrm{kN\cdot m}.
\]
其中
\[
  M_B=80\times8-80\times4=320\,\mathrm{kN\cdot m}.
\]
\begin{figure}[H]
\centering
\begin{tikzpicture}[scale=0.9,>=Stealth]
  \begin{scope}
    \draw[member] (0,0)--(0,2.6);
    \draw[Brand,thick] (0,0)--(1.2,2.6);
    \draw (0,2.6)--(1.2,2.6);
    \node[right] at (1.2,2.6) {\(320\)};
    \node[below] at (0.6,-0.25) {\(M_{AB}\) 图};
  \end{scope}
  \begin{scope}[shift={(3.4,0)}]
    \draw[member] (0,0)--(0,2.6);
    \draw[Brand,thick] (0.85,0)--(0.85,2.6);
    \draw (0,0)--(0.85,0);
    \draw (0,2.6)--(0.85,2.6);
    \node[right] at (0.85,1.3) {\(80\)};
    \node[below] at (0.45,-0.25) {\(Q_{AB}\) 图};
  \end{scope}
\end{tikzpicture}
\end{figure}
若取相反截面正号，则图形方向相反，但控制值不变。
\end{solutioncard}

\begin{problemcard}{2-4 刚架杆端弯矩}
图示刚架承受顶梁均布荷载 \(q\)，底部两端支承距离如图。求 \(AD\) 杆截面 \(D\) 的弯矩 \(M_{DA}\)，以内侧受拉为正。
\begin{figure}[H]
\centering
\begin{tikzpicture}[scale=0.78,>=Stealth]
  \draw[member] (0,0)--(1.0,4.0)--(5.0,4.0)--(6.0,0);
  \roller{0}{0}
  \wallroller{6.0}{0}
  \foreach \x in {1.25,1.65,...,4.75}
    \draw[Brand,-{Stealth[length=1.4mm]}] (\x,4.62)--(\x,4.08);
  \draw[Brand] (1.05,4.62)--(4.95,4.62);
  \node[above] at (3.0,4.62) {\(q\)};
  \node[below] at (0,0) {\(A\)};
  \node[below] at (6.0,0) {\(B\)};
  \node[above] at (1.0,4.0) {\(D\)};
  \node[above] at (3.0,4.0) {\(C\)};
  \node[above] at (5.0,4.0) {\(E\)};
  \draw[<->] (0,-0.55)--(1,-0.55) node[midway,below] {\(a\)};
  \draw[<->] (1,-0.55)--(3,-0.55) node[midway,below] {\(2a\)};
  \draw[<->] (3,-0.55)--(5,-0.55) node[midway,below] {\(2a\)};
  \draw[<->] (5,-0.55)--(6,-0.55) node[midway,below] {\(a\)};
  \draw[<->] (6.55,0)--(6.55,4) node[midway,right] {\(4a\)};
\end{tikzpicture}
\end{figure}
\end{problemcard}

\begin{solutioncard}
由整体竖向平衡与对称性，两个支座竖向反力均为
\[
  F_{Ay}=F_{By}=2qa.
\]
对左柱 \(AD\) 的 \(D\) 截面取隔离体。底部反力 \(2qa\) 到截面 \(D\) 的水平距离为 \(a\)，外侧水平反力按整体平衡抵消后，截面弯矩控制值为
\[
  M_{DA}=2qa\cdot a=2qa^2.
\]
题目规定以内侧受拉为正，故 \answer{\(M_{DA}=+2qa^2\)}。这一题的关键是先用整体对称求支座竖向反力，再对单侧柱取截面。
\end{solutioncard}

\begin{problemcard}{2-5 多跨梁等效弯矩图}
试作图示结构的弯矩图。题中各相邻标距均为 \(2\,\mathrm{m}\)。
\begin{figure}[H]
\centering
\begin{tikzpicture}[scale=0.76,>=Stealth]
  \draw[member] (0,0)--(7,0);
  \wallroller{0}{0}
  \node[smallpin] at (1,0) {};
  \node[smallpin] at (3,0) {};
  \roller{4}{0}
  \node[smallpin] at (5,0) {};
  \roller{5.5}{0}
  \node[above] at (1,0) {\(B\)};
  \node[above] at (3,0) {\(C\)};
  \node[above] at (4,0) {\(D\)};
  \node[above] at (5,0) {\(E\)};
  \node[above] at (5.5,0) {\(F\)};
  \node[above] at (7,0) {\(G\)};
  \draw[Brand,-{Stealth[length=2mm]}] (1,0.95)--(1,0.08) node[midway,left] {\(4\,\mathrm{kN}\)};
  \draw[Brand,-{Stealth[length=2mm]}] (2,1.05)--(2,0.08) node[midway,right] {\(10\,\mathrm{kN}\)};
  \foreach \x in {5.7,6.0,...,6.9}
    \draw[Brand,-{Stealth[length=1.4mm]}] (\x,0.72)--(\x,0.08);
  \draw[Brand] (5.6,0.72)--(7.0,0.72);
  \node[above] at (6.35,0.72) {\(6\,\mathrm{kN/m}\)};
  \foreach \x in {0,1,...,7}
    \draw[<->] (\x,-0.55)--(\x+1,-0.55);
  \node[below] at (3.5,-0.75) {每段 \(2\,\mathrm{m}\)};
\end{tikzpicture}
\end{figure}
\end{problemcard}

\begin{solutioncard}
先识别层次：含内铰的 \(BC\) 段为附属部分，左右基本部分分别承担附属部分传来的铰力。按附属部分 \(BC\) 平衡可求出传到两侧铰的竖向力，再分别在 \(AB\) 与 \(CDFG\) 上作弯矩图。

按原答案的控制值整理为：
\[
  M_A=18\,\mathrm{kN\cdot m},\qquad
  M_E=-5\,\mathrm{kN\cdot m},\qquad
  M_F=12\,\mathrm{kN\cdot m}.
\]
其中 \(B,C\) 为铰点，故
\[
  M_B=M_C=0.
\]
\begin{figure}[H]
\centering
\begin{tikzpicture}[scale=0.86,>=Stealth]
  \draw[member] (0,0)--(7,0);
  \draw[Brand,thick] (0,1.15)--(1,0)--(2,-0.45)--(3,0)--(4,0.55)--(5,0)--(5.5,-0.35)--(7,0.80);
  \draw (0,0)--(0,1.15);
  \node[left] at (0,1.15) {\(18\)};
  \node[below] at (2,-0.45) {\(10\)};
  \node[above] at (4,0.55) {\(10\)};
  \node[below] at (5.5,-0.35) {\(5\)};
  \node[right] at (7,0.80) {\(12\)};
  \node[below] at (3.5,-0.85) {\(M\) 图（\(\mathrm{kN\cdot m}\)）};
\end{tikzpicture}
\end{figure}
检查时应确保铰点处弯矩为零；均布荷载区弯矩为二次曲线，其余无分布荷载区为直线。
\end{solutioncard}

\begin{problemcard}{2-6 多跨梁弯矩图、剪力图与 \(AB\) 杆平衡}
作图示多跨梁的弯矩图和剪力图，并给出 \(AB\) 杆的受力平衡图。各跨均为 \(2\,\mathrm{m}\)。
\begin{figure}[H]
\centering
\begin{tikzpicture}[scale=0.74,>=Stealth]
  \draw[member] (0,0)--(8,0);
  \wallroller{0}{0}
  \roller{1}{0}
  \node[smallpin] at (2,0) {};
  \node[smallpin] at (4,0) {};
  \roller{5.5}{0}
  \node[smallpin] at (6.5,0) {};
  \wallroller{8}{0}
  \draw[Brand,->,thick] (1.2,0.55) arc[start angle=160,end angle=-165,radius=0.32];
  \node[above] at (1.2,0.85) {\(40\,\mathrm{kN\cdot m}\)};
  \foreach \x in {2.2,2.55,...,4.2}
    \draw[Brand,-{Stealth[length=1.4mm]}] (\x,0.68)--(\x,0.08);
  \draw[Brand] (2.1,0.68)--(4.3,0.68);
  \node[above] at (3.2,0.68) {\(10\,\mathrm{kN/m}\)};
  \draw[Brand,-{Stealth[length=2mm]}] (5.0,1.0)--(5.0,0.08) node[midway,right] {\(20\,\mathrm{kN}\)};
  \node[below] at (0,0) {\(B\)};
  \node[below] at (2,0) {\(A\)};
  \draw[<->] (0,-0.55)--(8,-0.55) node[midway,below] {\(8\times2\,\mathrm{m}\)};
\end{tikzpicture}
\end{figure}
\end{problemcard}

\begin{solutioncard}
从右端附属段向左递推。均布荷载合力为
\[
  10\times4=40\,\mathrm{kN},
\]
集中力为 \(20\,\mathrm{kN}\)。由附属段平衡求得传给基本部分的等效竖向力，再叠加左端 \(40\,\mathrm{kN\cdot m}\) 力偶。按答案图控制值可写为
\[
  M_A=40\,\mathrm{kN\cdot m},\qquad
  M_D=120\,\mathrm{kN\cdot m}.
\]
\[
  M_D=20\times4+20\times2=120\,\mathrm{kN\cdot m}.
\]
\begin{figure}[H]
\centering
\begin{tikzpicture}[scale=0.78,>=Stealth]
  \begin{scope}
    \draw[member] (0,0)--(8,0);
    \draw[Brand,thick] (0,0.70)--(1.0,0.35)--(2,0)--(3.3,-0.25)--(4.3,0.10)--(5.5,1.20)--(6.5,0)--(8,0);
    \node[left] at (0,0.70) {\(80\)};
    \node[above] at (5.5,1.20) {\(120\)};
    \node[below] at (4,-0.75) {\(M\) 图};
  \end{scope}
  \begin{scope}[shift={(0,-2.2)}]
    \draw[member] (0,0)--(8,0);
    \draw[Brand,thick] (0,0.65)--(2,0.65)--(2,-0.35)--(4,-0.35)--(4,0.35)--(5.5,0.35)--(5.5,-0.65)--(8,-0.65);
    \node[left] at (0,0.65) {\(40\)};
    \node[right] at (8,-0.65) {\(20\)};
    \node[below] at (4,-1.0) {\(Q\) 图};
  \end{scope}
\end{tikzpicture}
\end{figure}

\(AB\) 杆的受力平衡图可将 \(40\,\mathrm{kN\cdot m}\) 力偶、端部竖向力和水平约束力单独列出。检查条件为
\[
  \sum X=0,\qquad \sum Y=0,\qquad \sum M_A=0.
\]
\end{solutioncard}

\section{二、弯矩图、剪力图与内力图}

\begin{problemcard}{2-7 梁的内力图与弯矩图}
求图示梁的内力，并作弯矩图。左端为水平约束，梁上有两个 \(40\,\mathrm{kN}\) 集中力，右半跨有 \(40\,\mathrm{kN/m}\) 均布荷载。
\begin{figure}[H]
\centering
\begin{tikzpicture}[scale=0.8,>=Stealth]
  \draw[member] (0,0)--(7,0);
  \wallroller{0}{0}
  \roller{4}{0}
  \roller{7}{0}
  \draw[Brand,-{Stealth[length=2mm]}] (1,1.0)--(1,0.08) node[midway,left] {\(40\)};
  \draw[Brand,-{Stealth[length=2mm]}] (2,1.0)--(2,0.08) node[midway,right] {\(40\)};
  \foreach \x in {4.25,4.55,...,6.85}
    \draw[Brand,-{Stealth[length=1.4mm]}] (\x,0.72)--(\x,0.08);
  \draw[Brand] (4.1,0.72)--(6.95,0.72);
  \node[above] at (5.5,0.72) {\(40\,\mathrm{kN/m}\)};
  \draw[<->] (0,-0.55)--(1,-0.55) node[midway,below] {\(2m\)};
  \draw[<->] (1,-0.55)--(2,-0.55) node[midway,below] {\(2m\)};
  \draw[<->] (2,-0.55)--(4,-0.55) node[midway,below] {\(4m\)};
  \draw[<->] (4,-0.55)--(7,-0.55) node[midway,below] {\(4m\)};
\end{tikzpicture}
\end{figure}
\end{problemcard}

\begin{solutioncard}
均布荷载合力为 \(40\times4=160\,\mathrm{kN}\)。右跨按简支受均布荷载处理，可先得两端等效竖向力各 \(80\,\mathrm{kN}\)，再与左侧两个集中力叠加。答案图的弯矩控制值为
\[
  M_A=0,\qquad M_B=240\,\mathrm{kN\cdot m},\qquad M_C=0,\qquad M_D=0.
\]
其中
\[
  M_B=80\times2+40\times2\times1=240\,\mathrm{kN\cdot m}.
\]
\begin{figure}[H]
\centering
\begin{tikzpicture}[scale=0.84,>=Stealth]
  \draw[member] (0,0)--(7,0);
  \draw[Brand,thick] (0,0)--(1,0.35)--(2,0.70)--(4,1.20)..controls(5.2,0.45)and(6.2,0.10)..(7,0);
  \node[above] at (1,0.35) {\(80\)};
  \node[above] at (2,0.70) {\(160\)};
  \node[above] at (4,1.20) {\(240\)};
  \node[below] at (3.5,-0.35) {\(M\) 图};
\end{tikzpicture}
\end{figure}
作图检查：集中力处剪力突变，弯矩图斜率突变；均布荷载段弯矩为抛物线，并在右端支座处回零。
\end{solutioncard}

\begin{problemcard}{2-8 刚架弯矩图}
绘出图示刚架的弯矩图。各水平段长度为 \(a\)，各竖向杆高度为 \(a\)，图中水平荷载均为 \(F_p\)。
\begin{figure}[H]
\centering
\begin{tikzpicture}[scale=0.82,>=Stealth]
  \draw[member] (0,0)--(0,1.6)--(2,1.6)--(2,0)--(4,0)--(4,1.6)--(6,1.6)--(6,0);
  \fixedbase{0}{0}
  \roller{2}{0}
  \wallroller{4}{0.8}
  \roller{6}{0}
  \foreach \p in {(0,1.6),(2,1.6),(4,1.6)}
    \draw[Brand,-{Stealth[length=2mm]}] \p++(0,0.65)--++(0.95,0) node[midway,above] {\(F_p\)};
  \node[below] at (0,0) {\(A\)};
  \node[below] at (2,0) {\(C\)};
  \node[below] at (4,0) {\(D\)};
  \draw[<->] (0,-0.55)--(2,-0.55) node[midway,below] {\(a\)};
  \draw[<->] (2,-0.55)--(4,-0.55) node[midway,below] {\(a\)};
  \draw[<->] (4,-0.55)--(6,-0.55) node[midway,below] {\(a\)};
\end{tikzpicture}
\end{figure}
\end{problemcard}

\begin{solutioncard}
左、右两段可分别取隔离体。左端固定端的控制弯矩由上部水平力与中部传力叠加：
\[
  M_A=2F_pa-F_pa=F_pa.
\]
中间铰支承处弯矩为零，右侧短柱段按同样方式得到端部控制值。弯矩图在各直杆无分布荷载处均为直线。
\begin{figure}[H]
\centering
\begin{tikzpicture}[scale=0.84,>=Stealth]
  \draw[member] (0,0)--(0,1.6)--(2,1.6)--(2,0)--(4,0)--(4,1.6)--(6,1.6)--(6,0);
  \draw[Brand,thick] (0.45,0)--(0.45,1.6)--(2,1.6);
  \draw[Brand,thick] (4,0)--(4.45,1.6)--(6,1.6);
  \node[right] at (0.45,0.8) {\(F_pa\)};
  \node[right] at (4.45,1.0) {\(F_pa\)};
  \node[below] at (3,-0.35) {\(M\) 图};
\end{tikzpicture}
\end{figure}
检查要点：铰点弯矩必须为零；水平力只通过与杆高 \(a\) 的力臂形成 \(F_pa\) 级别的弯矩。
\end{solutioncard}

\begin{problemcard}{2-9 刚架弯矩图}
作图示静定刚架的弯矩图。
\begin{figure}[H]
\centering
\begin{tikzpicture}[scale=0.78,>=Stealth]
  \draw[member] (0,0)--(0,3)--(2,3)--(2,1.5)--(4,1.5)--(4,0);
  \roller{0}{0}
  \roller{4}{0}
  \draw[Brand,->,thick] (2,1.2) arc[start angle=-80,end angle=250,radius=0.32];
  \node[below] at (2,0.95) {\(3\,\mathrm{kN\cdot m}\)};
  \draw[<->] (0,-0.55)--(2,-0.55) node[midway,below] {\(2m\)};
  \draw[<->] (2,-0.55)--(4,-0.55) node[midway,below] {\(2m\)};
  \draw[<->] (-0.55,0)--(-0.55,1.5) node[midway,left] {\(1m\)};
  \draw[<->] (-0.55,1.5)--(-0.55,3) node[midway,left] {\(2m\)};
\end{tikzpicture}
\end{figure}
\end{problemcard}

\begin{solutioncard}
由整体对 \(B\) 点取矩可得左支反力
\[
  F_{Ay}\cdot4=3+2\times2\times\frac32=12,\qquad F_{Ay}=3\,\mathrm{kN}.
\]
再逐段隔离可得答案图控制值：
\[
  M_{DC}=2\,\mathrm{kN\cdot m},\qquad
  F_{Bx}=1\,\mathrm{kN},\qquad
  M_{ED}=5\,\mathrm{kN\cdot m}.
\]
\begin{figure}[H]
\centering
\begin{tikzpicture}[scale=0.82,>=Stealth]
  \draw[member] (0,0)--(0,3)--(2,3)--(2,1.5)--(4,1.5)--(4,0);
  \draw[Brand,thick] (0,0)--(0.55,3)--(2,3)--(2.45,1.5)--(4.55,1.5)--(4,0);
  \node[right] at (2.45,1.5) {\(2\)};
  \node[right] at (4.55,1.5) {\(5\)};
  \node[below] at (2,-0.35) {\(M\) 图};
\end{tikzpicture}
\end{figure}
\end{solutioncard}

\begin{problemcard}{2-10 门式刚架弯矩图}
绘制图示结构的弯矩图。刚架顶梁有均布荷载，左侧顶端有水平集中力。
\begin{figure}[H]
\centering
\begin{tikzpicture}[scale=0.72,>=Stealth]
  \draw[member] (0,0)--(0,4)--(3,4)--(3,3)--(6,3)--(6,0);
  \roller{0}{0}
  \roller{6}{0}
  \node[smallpin] at (3,4) {};
  \draw[Brand,-{Stealth[length=2mm]}] (-0.8,4)--(0,4) node[midway,above] {\(20\)};
  \foreach \x in {3.2,3.55,...,5.8}
    \draw[Brand,-{Stealth[length=1.4mm]}] (\x,3.65)--(\x,3.05);
  \draw[Brand] (3.05,3.65)--(5.95,3.65);
  \node[above] at (4.5,3.65) {\(10\,\mathrm{kN/m}\)};
  \draw[<->] (0,-0.55)--(3,-0.55) node[midway,below] {\(3m\)};
  \draw[<->] (3,-0.55)--(6,-0.55) node[midway,below] {\(3m\)};
  \draw[<->] (6.55,0)--(6.55,3) node[midway,right] {\(6m\)};
\end{tikzpicture}
\end{figure}
\end{problemcard}

\begin{solutioncard}
先取右侧隔离体，由均布荷载合力与支座反力平衡得中间竖向反力控制值
\[
  F_{OD}=40\,\mathrm{kN}.
\]
答案弯矩图主要控制值为
\[
  M_A=120\,\mathrm{kN\cdot m},\quad
  M_D=165\,\mathrm{kN\cdot m},\quad
  M_E=170\,\mathrm{kN\cdot m}.
\]
\begin{figure}[H]
\centering
\begin{tikzpicture}[scale=0.78,>=Stealth]
  \draw[member] (0,0)--(0,4)--(3,4)--(3,3)--(6,3)--(6,0);
  \draw[Brand,thick] (0.55,0)--(0.55,4)--(3,4.35)--(3,3)--(6,3.35)--(6.55,0);
  \node[right] at (0.55,1.8) {\(120\)};
  \node[above] at (3,4.35) {\(165\)};
  \node[right] at (6.55,2.0) {\(170\)};
  \node[below] at (3,-0.35) {\(M\) 图};
\end{tikzpicture}
\end{figure}
均布荷载作用段的弯矩边线应为抛物线；本图以控制值示意，正式作图时将顶梁荷载段画成二次曲线。
\end{solutioncard}

\begin{problemcard}{2-11 中间铰梁弯矩图}
绘图示结构的弯矩图。
\begin{figure}[H]
\centering
\begin{tikzpicture}[scale=0.78,>=Stealth]
  \draw[member] (0,0)--(0,2.5)--(2,2.5)--(2,2.5)--(4,2.5);
  \draw[member] (0,2.5)--(0,0);
  \roller{0}{0}
  \roller{4}{2.5}
  \node[smallpin] at (2,2.5) {};
  \foreach \x in {0.35,0.7,...,3.65}
    \draw[Brand,-{Stealth[length=1.3mm]}] (\x,3.15)--(\x,2.58);
  \draw[Brand] (0.2,3.15)--(3.8,3.15);
  \node[above] at (2,3.15) {\(q\)};
  \draw[Brand,-{Stealth[length=1.8mm]}] (-0.75,2.5)--(0,2.5) node[midway,above] {\(ql\)};
  \draw[<->] (0,-0.55)--(2,-0.55) node[midway,below] {\(l\)};
  \draw[<->] (2,-0.55)--(4,-0.55) node[midway,below] {\(l\)};
\end{tikzpicture}
\end{figure}
\end{problemcard}

\begin{solutioncard}
中间铰处弯矩为零。先对右侧跨取矩可得铰力
\[
  F_{RC}\,l-ql\cdot\frac{l}{2}=ql^2,\qquad
  F_{RC}=\frac32ql.
\]
再回到左侧，得到右端控制弯矩
\[
  M_D=ql^2,\qquad M_B=2ql^2.
\]
弯矩图在均布荷载段为抛物线，在无分布荷载段为直线。
\end{solutioncard}

\begin{problemcard}{2-12 由弯矩图求剪力图与轴力图}
已知图示刚架的弯矩图，求剪力图和轴力图。已知 \(F_p=10\,\mathrm{kN}\)，端部力偶 \(M=40\,\mathrm{kN\cdot m}\)。
\begin{figure}[H]
\centering
\begin{tikzpicture}[scale=0.78,>=Stealth]
  \draw[member] (0,0)--(0,4)--(2,4)--(4,4)--(4,0);
  \fixedbase{0}{0}
  \fixedbase{4}{0}
  \node[smallpin] at (2,4) {};
  \draw[Brand,thick] (0,3.7)--(2,4.8)--(2,4)--(4,4.55)--(4,3.7);
  \node[above] at (0.6,4.1) {\(20\)};
  \node[above] at (1.8,4.8) {\(40\)};
  \node[above] at (3.2,4.55) {\(20\)};
  \draw[Brand,-{Stealth[length=2mm]}] (-0.75,2)--(0,2) node[midway,above] {\(F_p\)};
  \draw[Brand,->,thick] (2.25,4.25) arc[start angle=30,end angle=330,radius=0.3];
  \node[right] at (2.55,4.25) {\(M\)};
  \node[below] at (2,-0.35) {已知 \(M\) 图};
\end{tikzpicture}
\end{figure}
\end{problemcard}

\begin{solutioncard}
由弯矩图斜率求剪力：
\[
  Q=\frac{\Delta M}{\Delta s}.
\]
两侧柱高均为 \(4\,\mathrm{m}\)，顶梁两半跨均为 \(2\,\mathrm{m}\)，故各段剪力控制值为 \(10\,\mathrm{kN}\) 量级。结合外力 \(F_p=10\,\mathrm{kN}\)，轴力图在两柱中为 \(10\,\mathrm{kN}\)，顶梁段轴力由柱端水平剪力平衡确定。
\begin{figure}[H]
\centering
\begin{tikzpicture}[scale=0.78,>=Stealth]
  \begin{scope}
    \draw[member] (0,0)--(0,3)--(2,3)--(4,3)--(4,0);
    \draw[Brand,thick] (0.35,0)--(0.35,3)--(2,3.35)--(4,3)--(4.35,0);
    \node[right] at (0.35,1.5) {\(10\)};
    \node[above] at (2,3.35) {\(10\)};
    \node[below] at (2,-0.35) {\(Q\) 图};
  \end{scope}
  \begin{scope}[shift={(6,0)}]
    \draw[member] (0,0)--(0,3)--(2,3)--(4,3)--(4,0);
    \draw[Brand,thick] (-0.35,0)--(-0.35,3)--(4.35,3)--(4.35,0);
    \node[left] at (-0.35,1.5) {\(10\)};
    \node[above] at (2,3) {\(10\)};
    \node[right] at (4.35,1.5) {\(10\)};
    \node[below] at (2,-0.35) {\(N\) 图};
  \end{scope}
\end{tikzpicture}
\end{figure}
检查：弯矩图为直线的杆段剪力为常值；弯矩图水平的杆段剪力为零。
\end{solutioncard}

\begin{problemcard}{2-13 刚架弯矩图}
作图示刚架的弯矩图。
\begin{figure}[H]
\centering
\begin{tikzpicture}[scale=0.78,>=Stealth]
  \draw[member] (0,2.4)--(2,2.4)--(4,2.4)--(4,0);
  \draw[member] (2,2.4)--(2,0);
  \roller{4}{0}
  \wallroller{4}{2.4}
  \node[smallpin] at (2,2.4) {};
  \draw[Brand,-{Stealth[length=2mm]}] (0,3.2)--(0,2.45) node[midway,left] {\(F_p=ql\)};
  \foreach \x in {2.35,2.7,...,3.75}
    \draw[Brand,-{Stealth[length=1.4mm]}] (\x,3.0)--(\x,2.48);
  \draw[Brand] (2.25,3.0)--(3.9,3.0);
  \node[above] at (3.1,3.0) {\(q\)};
  \draw[<->] (0,-0.55)--(2,-0.55) node[midway,below] {\(l\)};
  \draw[<->] (2,-0.55)--(4,-0.55) node[midway,below] {\(l\)};
  \draw[<->] (4.55,0)--(4.55,2.4) node[midway,right] {\(l\)};
\end{tikzpicture}
\end{figure}
\end{problemcard}

\begin{solutioncard}
由整体竖向平衡可得
\[
  F_{Ay}=2ql.
\]
对上部横梁和右侧柱分别取隔离体，答案控制值为
\[
  M_D=\frac12 ql^2.
\]
中间铰处弯矩为零；顶梁受均布荷载段为抛物线，左伸臂段为直线。
\end{solutioncard}

\begin{problemcard}{2-14 刚架弯矩图}
作图示结构的弯矩图。
\begin{figure}[H]
\centering
\begin{tikzpicture}[scale=0.78,>=Stealth]
  \draw[member] (0,0)--(2,3)--(4,3)--(4,0)--(0,0);
  \fixedbase{4}{0}
  \draw[Brand,-{Stealth[length=2mm]}] (0,-0.75)--(0,0) node[midway,left] {\(3\,\mathrm{kN}\)};
  \foreach \x in {0.2,0.55,...,1.85}
    \draw[Brand,-{Stealth[length=1.4mm]}] (\x,3.65)--(\x,3.05);
  \draw[Brand] (0,3.65)--(2,3.65);
  \node[above] at (1,3.65) {\(2\,\mathrm{kN/m}\)};
  \draw[Brand,->,thick] (4,3.45) arc[start angle=90,end angle=-250,radius=0.32];
  \node[right] at (4.35,3.35) {\(4\,\mathrm{kN\cdot m}\)};
  \draw[<->] (0,-0.55)--(2,-0.55) node[midway,below] {\(2m\)};
  \draw[<->] (2,-0.55)--(4,-0.55) node[midway,below] {\(2m\)};
  \draw[<->] (4.55,0)--(4.55,3) node[midway,right] {\(3m\)};
\end{tikzpicture}
\end{figure}
\end{problemcard}

\begin{solutioncard}
均布荷载合力为 \(2\times2=4\,\mathrm{kN}\)。按原答案逐段取矩：
\[
  M_B=4+2\times2\times1=8\,\mathrm{kN\cdot m},
\]
\[
  M_C=3\times4+2\times2\times2-4=20\,\mathrm{kN\cdot m}.
\]
弯矩图控制值为 \(4,8,20\,\mathrm{kN\cdot m}\)，斜杆和无分布荷载杆段均为直线。
\end{solutioncard}

\begin{problemcard}{2-15 刚架内力图}
作图示结构的内力图。
\begin{figure}[H]
\centering
\begin{tikzpicture}[scale=0.68,>=Stealth]
  \draw[member] (0,0)--(0,2)--(1,2)--(1,5)--(7,5)--(7,2)--(8.5,2)--(8.5,0);
  \roller{0}{0}
  \roller{1}{0}
  \roller{7}{0}
  \wallroller{8.5}{0}
  \foreach \y in {0.4,0.8,1.2,1.6}
    \draw[Brand,-{Stealth[length=1.4mm]}] (-0.75,\y)--(-0.05,\y);
  \draw[Brand] (-0.75,0.25)--(-0.75,1.85);
  \node[left] at (-0.75,1.0) {\(10\,\mathrm{kN/m}\)};
  \draw[Brand,->,thick] (1.2,5.35) arc[start angle=160,end angle=-165,radius=0.32];
  \node[above] at (1.8,5.35) {\(30\,\mathrm{kN\cdot m}\)};
  \draw[Brand,-{Stealth[length=2mm]}] (8.5,3.2)--(8.5,2.05) node[midway,right] {\(10\,\mathrm{kN}\)};
  \draw[Brand,-{Stealth[length=2mm]}] (9.5,2)--(8.55,2) node[midway,above] {\(30\,\mathrm{kN}\)};
  \draw[<->] (1,-0.55)--(7,-0.55) node[midway,below] {\(8m\)};
  \draw[<->] (8.95,2)--(8.95,5) node[midway,right] {\(4m\)};
\end{tikzpicture}
\end{figure}
\end{problemcard}

\begin{solutioncard}
先取右侧附属杆段，由端部 \(10\,\mathrm{kN}\) 竖向力和 \(30\,\mathrm{kN}\) 水平力求出传给主体刚架的端力。左侧均布水平荷载合力为
\[
  10\times4=40\,\mathrm{kN},
\]
对底部取矩得
\[
  10\times4\times2=80\,\mathrm{kN\cdot m}.
\]
答案图中剪力、轴力的主要控制值为
\[
  40,\quad 30,\quad 8.75,\quad 48.75\quad(\mathrm{kN}),
\]
弯矩图主要控制值为
\[
  30,\quad 40,\quad 80,\quad 120\quad(\mathrm{kN\cdot m}).
\]
作图时按照“先附属段，后主体刚架”的顺序，把附属段传来的力作为主体刚架外力。
\end{solutioncard}

\begin{problemcard}{2-16 刚架内力与弯矩图}
求图示结构的内力，并作弯矩图。
\begin{figure}[H]
\centering
\begin{tikzpicture}[scale=0.78,>=Stealth]
  \draw[member] (0,0)--(0,1.6)--(2,1.6)--(4,1.6)--(6,1.6)--(6,0);
  \draw[member] (4,1.6)--(4,0);
  \wallroller{0}{0}
  \roller{4}{0}
  \roller{6}{0}
  \draw[Brand,-{Stealth[length=2mm]}] (-0.85,0.8)--(0,0.8) node[midway,above] {\(F_p\)};
  \draw[Brand,-{Stealth[length=2mm]}] (2,2.45)--(2,1.68) node[midway,left] {\(F_p\)};
  \draw[Brand,->,thick] (6,2.0) arc[start angle=80,end angle=-260,radius=0.35];
  \node[right] at (6.4,2.0) {\(F_pa\)};
  \draw[<->] (0,-0.55)--(2,-0.55) node[midway,below] {\(a\)};
  \draw[<->] (2,-0.55)--(4,-0.55) node[midway,below] {\(a\)};
  \draw[<->] (4,-0.55)--(6,-0.55) node[midway,below] {\(2a\)};
  \draw[<->] (6.55,0)--(6.55,1.6) node[midway,right] {\(2a\)};
\end{tikzpicture}
\end{figure}
\end{problemcard}

\begin{solutioncard}
对左侧 \(ABC\) 部分取隔离体，由竖向平衡有铰力平衡；对 \(B\) 点取矩可得左侧控制弯矩。由整体水平方向平衡：
\[
  F_{Ex}=\frac12F_p.
\]
右端力偶使右柱端弯矩出现 \(F_pa\) 控制值，且在无横向荷载杆段弯矩为直线。答案控制值为
\[
  M_{DC}=F_pa,\qquad M_C=M_B=0.
\]
作图检查：铰点弯矩为零；集中力偶只改变弯矩图，不改变剪力图。
\end{solutioncard}

\section{三、综合刚架与反推荷载}

\begin{problemcard}{2-17 刚架弯矩图}
作图示结构的弯矩图。左侧附属柱受水平均布荷载，顶梁受竖向均布荷载。
\begin{figure}[H]
\centering
\begin{tikzpicture}[scale=0.68,>=Stealth]
  \draw[member] (0,0)--(0,3)--(2.5,3)--(5.5,3)--(5.5,0);
  \draw[member] (-2,0)--(-2,1.5)--(0,1.5);
  \roller{-2}{0}
  \roller{0}{0}
  \roller{5.5}{0}
  \node[smallpin] at (0,1.5) {};
  \foreach \y in {0.35,0.7,1.05,1.4}
    \draw[Brand,-{Stealth[length=1.4mm]}] (-2.7,\y)--(-2.05,\y);
  \draw[Brand] (-2.7,0.2)--(-2.7,1.55);
  \node[left] at (-2.7,0.95) {\(10\,\mathrm{kN/m}\)};
  \foreach \x in {0.3,0.65,...,5.25}
    \draw[Brand,-{Stealth[length=1.4mm]}] (\x,3.65)--(\x,3.05);
  \draw[Brand] (0.15,3.65)--(5.35,3.65);
  \node[above] at (2.75,3.65) {\(20\,\mathrm{kN/m}\)};
  \draw[<->] (-2,-0.55)--(0,-0.55) node[midway,below] {\(6m\)};
  \draw[<->] (0,-0.55)--(2.75,-0.55) node[midway,below] {\(6m\)};
  \draw[<->] (2.75,-0.55)--(5.5,-0.55) node[midway,below] {\(6m\)};
\end{tikzpicture}
\end{figure}
\end{problemcard}

\begin{solutioncard}
从右侧主体刚架取矩求竖向反力。按原答案手写过程，
\[
  F_y\cdot12=60\times7.5+20\times12\times6,
\]
故
\[
  F_y=168.75\,\mathrm{kN}.
\]
再对下段水平平衡与节点取矩，得到
\[
  F_x=72.5\,\mathrm{kN},\qquad F_{Dx}=-17.5\,\mathrm{kN}.
\]
弯矩图控制值整理为
\[
  67.5,\quad 105,\quad 360,\quad 652.5\quad(\mathrm{kN\cdot m}).
\]
\begin{figure}[H]
\centering
\begin{tikzpicture}[scale=0.68,>=Stealth]
  \draw[member] (0,0)--(0,3)--(2.5,3)--(5.5,3)--(5.5,0);
  \draw[member] (-2,0)--(-2,1.5)--(0,1.5);
  \draw[Brand,thick] (-2.5,0)--(-2.25,1.5)--(0.45,1.5)--(0.65,3)--(2.5,3.8)--(5.9,3)--(5.9,0);
  \node[left] at (-2.5,0.8) {\(67.5\)};
  \node[above] at (2.5,3.8) {\(360\)};
  \node[right] at (5.9,1.6) {\(652.5\)};
  \node[below] at (2.6,-0.35) {\(M\) 图（\(\mathrm{kN\cdot m}\)）};
\end{tikzpicture}
\end{figure}
检查：左侧附属柱的荷载先等效到铰点，主体刚架再承受该铰力；顶梁均布荷载段弯矩应为抛物线。
\end{solutioncard}

\begin{problemcard}{2-18 刚架弯矩图}
作图示结构的弯矩图。
\begin{figure}[H]
\centering
\begin{tikzpicture}[scale=0.7,>=Stealth]
  \draw[member] (0,0)--(0,3)--(2,3)--(4,3)--(4,0);
  \draw[member] (4,3)--(6,3)--(6,1.5)--(6,0);
  \roller{0}{0}
  \roller{4}{0}
  \roller{6}{0}
  \node[smallpin] at (2,3) {};
  \foreach \x in {0.25,0.6,...,1.9}
    \draw[Brand,-{Stealth[length=1.3mm]}] (\x,3.58)--(\x,3.05);
  \draw[Brand] (0.1,3.58)--(2,3.58);
  \node[above] at (1,3.58) {\(10\,\mathrm{kN/m}\)};
  \draw[Brand,-{Stealth[length=2mm]}] (4.9,4.0)--(4.9,3.08) node[midway,right] {\(30\,\mathrm{kN}\)};
  \draw[Brand,-{Stealth[length=2mm]}] (6.85,1.5)--(6.05,1.5) node[midway,above] {\(40\,\mathrm{kN}\)};
  \draw[<->] (0,-0.55)--(2,-0.55) node[midway,below] {\(2m\)};
  \draw[<->] (2,-0.55)--(4,-0.55) node[midway,below] {\(2m\)};
  \draw[<->] (4,-0.55)--(6,-0.55) node[midway,below] {\(2m\)};
\end{tikzpicture}
\end{figure}
\end{problemcard}

\begin{solutioncard}
右侧 \(EFGH\) 杆段可先作为附属部分处理。水平 \(40\,\mathrm{kN}\) 作用在高 \(4\,\mathrm{m}\) 处，等效到连接处后，再求主体刚架支座反力。由原答案推得
\[
  F_{Hx}=25\,\mathrm{kN},\qquad
  F_{Hy}=35\,\mathrm{kN},\qquad
  F_{Bx}=75\,\mathrm{kN}.
\]
弯矩图控制值可整理为
\[
  10,\quad20,\quad30,\quad60\quad(\mathrm{kN\cdot m}).
\]
均布荷载作用段画抛物线，其余杆段画直线；中间铰处弯矩为零。
\end{solutioncard}

\begin{problemcard}{2-19 由弯矩图判定荷载}
已知结构在集中力和集中力偶作用下的弯矩图，试确定结构所受荷载，并绘出 \(BE\) 段隔离体受力平衡图。
\begin{figure}[H]
\centering
\begin{tikzpicture}[scale=0.72,>=Stealth]
  \draw[member] (0,0)--(0,3)--(2,3)--(4,3)--(4,0);
  \draw[member] (2,3)--(2.55,3.35)--(3.2,3.0);
  \roller{0}{0}
  \wallroller{4}{0}
  \node[smallpin] at (2.55,3.35) {};
  \draw[Brand,thick] (0.35,0)--(0.35,3)--(2,4.0)--(2.55,3.35)--(4.35,3)--(4.35,0);
  \node[left] at (0.35,2.2) {\(2qa^2\)};
  \node[above] at (2,4.0) {\(3qa^2\)};
  \node[right] at (4.35,2.3) {\(qa^2\)};
  \draw[<->] (0,-0.55)--(4,-0.55) node[midway,below] {\(2a\)};
  \draw[<->] (4.55,0)--(4.55,3) node[midway,right] {\(2a\)};
  \node[below] at (2,-0.9) {已知 \(M\) 图};
\end{tikzpicture}
\end{figure}
\end{problemcard}

\begin{solutioncard}
由弯矩图斜率变化可判断：弯矩图曲线段对应均布荷载，折线斜率突变对应集中力，弯矩突跳对应集中力偶。对左柱弯矩图的三角形面积关系：
\[
  \frac12 F_x(2a)=qa^2+2qa^2,
\]
故
\[
  F_x=3qa.
\]
由顶部抛物线段反推均布荷载强度，得
\[
  q_{\mathrm{top}}=2q.
\]
\(BE\) 段隔离体应标出铰力、由弯矩图斜率得到的剪力，以及相邻杆端传来的集中力偶；其平衡校核为
\[
  \sum X=0,\qquad \sum Y=0,\qquad \sum M_B=0.
\]
\end{solutioncard}

\begin{problemcard}{2-20 由弯矩图反推荷载、剪力图和轴力图}
已知图示静定刚架的弯矩图，其中曲线段为二次抛物线。试绘图标出作用于刚架上的荷载，并作剪力图和轴力图。
\begin{figure}[H]
\centering
\begin{tikzpicture}[scale=0.76,>=Stealth]
  \draw[member] (0,0)--(0,4)--(2,4)--(4,4)--(4,0);
  \fixedbase{0}{0}
  \wallroller{4}{0}
  \node[smallpin] at (2,4) {};
  \draw[Brand,thick] (0.35,0)--(0.35,4)--(2,4.55)--(4.35,4)--(4.35,0);
  \node[left] at (0.35,2) {\(2.5\)};
  \node[above] at (2,4.55) {\(3.5\)};
  \node[right] at (4.35,1.8) {\(4.5\)};
  \draw[<->] (0,-0.55)--(2,-0.55) node[midway,below] {\(2m\)};
  \draw[<->] (2,-0.55)--(4,-0.55) node[midway,below] {\(2m\)};
  \draw[<->] (4.55,0)--(4.55,2) node[midway,right] {\(2m\)};
  \draw[<->] (4.55,2)--(4.55,4) node[midway,right] {\(2m\)};
  \node[below] at (2,-0.9) {已知 \(M\) 图（\(\mathrm{kN\cdot m}\)）};
\end{tikzpicture}
\end{figure}
\end{problemcard}

\begin{solutioncard}
由弯矩图斜率突变处确定集中力：
\[
  \frac12 F\cdot4=1+5,\qquad F=6\,\mathrm{kN}.
\]
由二次抛物线段的弯矩差确定均布荷载：
\[
  \frac18 q\cdot4^2=5-1,\qquad q=2\,\mathrm{kN/m}.
\]
又由右柱端弯矩得
\[
  M_B=2\,\mathrm{kN\cdot m}.
\]
对底部支座取矩、列水平平衡，可得
\[
  F_{Bx}=4.5\,\mathrm{kN},\qquad F_{Ax}=2.5\,\mathrm{kN}.
\]
\begin{figure}[H]
\centering
\begin{tikzpicture}[scale=0.76,>=Stealth]
  \begin{scope}
    \draw[member] (0,0)--(0,3)--(2,3)--(4,3)--(4,0);
    \draw[Brand,thick] (-0.35,0)--(-0.35,3)--(2,3.35)--(4.35,3)--(4.35,0);
    \node[left] at (-0.35,1.4) {\(2.5\)};
    \node[above] at (2,3.35) {\(3.5\)};
    \node[right] at (4.35,1.4) {\(4.5\)};
    \node[below] at (2,-0.35) {\(Q\) 图};
  \end{scope}
  \begin{scope}[shift={(6,0)}]
    \draw[member] (0,0)--(0,3)--(2,3)--(4,3)--(4,0);
    \draw[Brand,thick] (-0.3,0)--(-0.3,3)--(4.3,3)--(4.3,0);
    \node[left] at (-0.3,1.5) {\(1\)};
    \node[above] at (2,3) {\(3.5\)};
    \node[right] at (4.3,1.5) {\(1\)};
    \node[below] at (2,-0.35) {\(N\) 图};
  \end{scope}
\end{tikzpicture}
\end{figure}
\end{solutioncard}

\begin{problemcard}{2-21 不经计算绘弯矩图}
不经计算，绘制图示结构的弯矩图。
\begin{figure}[H]
\centering
\begin{tikzpicture}[scale=0.8,>=Stealth]
  \draw[member] (0,0)--(3,0)--(6,0)--(11,0);
  \wallroller{0}{0}
  \roller{3}{0}
  \node[smallpin] at (6,0) {};
  \roller{11}{0}
  \draw[Brand,->,thick] (5.7,0.45) arc[start angle=160,end angle=-165,radius=0.34];
  \node[above] at (5.9,0.75) {\(10\,\mathrm{kN\cdot m}\)};
  \draw[<->] (0,-0.55)--(3,-0.55) node[midway,below] {\(3m\)};
  \draw[<->] (3,-0.55)--(6,-0.55) node[midway,below] {\(3m\)};
  \draw[<->] (6,-0.55)--(11,-0.55) node[midway,below] {\(5m\)};
\end{tikzpicture}
\end{figure}
\end{problemcard}

\begin{solutioncard}
该结构没有横向分布荷载，因此各直杆段弯矩图斜率不变。集中力偶只使弯矩图在该处突跳 \(10\,\mathrm{kN\cdot m}\)，剪力图不发生突变。左端 \(AB\) 段无外荷载且两端约束不形成弯矩，弯矩为零；\(BC\)、\(CD\) 两段为直线图。
\begin{figure}[H]
\centering
\begin{tikzpicture}[scale=0.86,>=Stealth]
  \draw[member] (0,0)--(3,0)--(6,0)--(11,0);
  \draw[Brand,thick] (0,0)--(3,0)--(6,0.70)--(6,0.20)--(11,0);
  \node[above] at (6,0.70) {\(10\)};
  \node[below] at (5.5,-0.35) {\(M\) 图};
\end{tikzpicture}
\end{figure}
\end{solutioncard}

\begin{problemcard}{2-22 刚架弯矩图}
试求解图示刚架，并绘制弯矩图。
\begin{figure}[H]
\centering
\begin{tikzpicture}[scale=0.64,>=Stealth]
  \draw[member] (0,0)--(0,3)--(2.5,3)--(2.5,5)--(5,5)--(5,3)--(7.5,3)--(7.5,0);
  \draw[member] (2.5,3)--(5,3);
  \roller{0}{0}
  \wallroller{2.5}{0}
  \roller{5}{0}
  \roller{7.5}{0}
  \node[smallpin] at (3.75,5) {};
  \draw[Brand,-{Stealth[length=2mm]}] (2.5,3.55)--(3.55,3.55) node[midway,above] {\(F_p\)};
  \draw[Brand,-{Stealth[length=2mm]}] (5,3.55)--(6.05,3.55) node[midway,above] {\(F_p\)};
  \draw[Brand,->,thick] (0,3.45) arc[start angle=90,end angle=-250,radius=0.34];
  \draw[Brand,->,thick] (7.5,3.45) arc[start angle=90,end angle=-250,radius=0.34];
  \node[left] at (0,3.45) {\(M\)};
  \node[right] at (7.5,3.45) {\(M\)};
  \draw[<->] (0,-0.55)--(2.5,-0.55) node[midway,below] {\(3m\)};
  \draw[<->] (2.5,-0.55)--(5,-0.55) node[midway,below] {\(3m\)};
  \draw[<->] (5,-0.55)--(7.5,-0.55) node[midway,below] {\(3m\)};
\end{tikzpicture}
\end{figure}
\end{problemcard}

\begin{solutioncard}
结构与荷载关于中间跨对称。先取左右两侧附属部分，水平 \(F_p\) 在柱顶产生 \(F_pa\) 级弯矩；再由对称条件求出 \(A,B\) 两处水平反力。对左侧隔离体取矩：
\[
  M_F=F_p\cdot5-F_p\cdot2=3F_p.
\]
因此弯矩图的主要控制值为 \(3F_p\)，中间铰处弯矩为零；外侧小跨受端部力偶影响，弯矩图在端部突跳。
\begin{figure}[H]
\centering
\begin{tikzpicture}[scale=0.64,>=Stealth]
  \draw[member] (0,0)--(0,3)--(2.5,3)--(2.5,5)--(5,5)--(5,3)--(7.5,3)--(7.5,0);
  \draw[member] (2.5,3)--(5,3);
  \draw[Brand,thick] (0.45,0)--(0.45,3)--(2.5,3.65)--(3.75,5)--(5,3.65)--(7.05,3)--(7.05,0);
  \node[right] at (0.45,1.5) {\(M\)};
  \node[above] at (3.75,5) {\(0\)};
  \node[above] at (2.5,3.65) {\(3F_p\)};
  \node[above] at (5,3.65) {\(3F_p\)};
  \node[below] at (3.75,-0.35) {\(M\) 图};
\end{tikzpicture}
\end{figure}
\end{solutioncard}

\begin{notecard}
第 4 次作业讲解笔记补充：静定刚架作图时，最容易错在“先后顺序”。含中间铰、伸臂或附属杆的结构，先取附属部分；由弯矩图反推荷载时，看三件事：直线斜率代表剪力，抛物线曲率代表均布荷载，竖向突跳代表集中力偶。所有控制值最后都要用 \(\sum X=0,\sum Y=0,\sum M=0\) 回代校核。
\end{notecard}

\chapter{静定桁架}

\begin{notecard}
静定桁架求解默认约定为：杆件轴力以受拉为正、受压为负。解题顺序通常为先判零杆，再求整体支座反力，最后按节点法或截面法求目标杆。节点上若仅有两根不共线杆且无外力，则两杆均为零杆；若三杆中两杆共线且节点无外力，则第三杆为零杆。
\end{notecard}

\section{一、零杆判别与基本节点法}

\begin{problemcard}{2-23 零杆数目}
图示桁架中零杆的数目为多少？不包括支座链杆。
\[
  \text{选项：}\quad
  \mathrm{A.}\ 6\text{ 根},\qquad
  \mathrm{B.}\ 7\text{ 根},\qquad
  \mathrm{C.}\ 8\text{ 根},\qquad
  \mathrm{D.}\ 9\text{ 根}.
\]
\begin{figure}[H]
\centering
\begin{tikzpicture}[scale=0.85,>=Stealth]
  \coordinate (A) at (0,1.6);
  \coordinate (B) at (0,0.6);
  \coordinate (C) at (0,-0.4);
  \coordinate (D) at (-1.1,-0.4);
  \coordinate (E) at (1.1,-0.4);
  \coordinate (F) at (1.1,0.6);
  \coordinate (G) at (2.2,0.6);
  \coordinate (H) at (2.2,-0.4);
  \draw[member] (A)--(B)--(C)--(D) (C)--(E)--(H)--(G)--(F)--(E);
  \draw[member] (B)--(F)--(G) (D)--(B) (A)--(F) (C)--(H);
  \wallroller{0}{1.6}
  \roller{-1.1}{-0.4}
  \wallroller{2.2}{0.6}
  \wallroller{2.2}{-1.1}
  \foreach \p in {A,B,C,D,E,F,G,H} \node[smallpin] at (\p) {};
  \draw[Brand,-{Stealth[length=2mm]}] (0.75,0.1)--(0.25,-0.3) node[midway,below] {\(F_p\)};
  \draw[Brand,-{Stealth[length=2mm]}] (0.7,0.1)--(1.0,0.5) node[midway,right] {\(F_p\)};
  \node[right] at (0.62,1.1) {1};
  \node[above] at (0.55,0.6) {2};
  \node[right] at (0,0.1) {3};
  \draw[<->] (-1.1,-0.85)--(0,-0.85) node[midway,below] {\(a\)};
  \draw[<->] (0,-0.85)--(1.1,-0.85) node[midway,below] {\(a\)};
  \draw[<->] (1.1,-0.85)--(2.2,-0.85) node[midway,below] {\(a\)};
\end{tikzpicture}
\end{figure}
\end{problemcard}

\begin{solutioncard}
\answer{答案：C，零杆共 \(8\) 根。}

判零杆时不把支座链杆计入。按“无外力两杆节点”和“三杆中两杆共线”逐步识别：先从无外力端节点剥离，得到第一组 \(2\) 根零杆；相邻节点随之成为新的无外力两杆或三杆节点，继续剥离可得
\[
  2+2+1+2+1=8.
\]
这一题即使题目没有特别强调，也应排除支座链杆，只数桁架内部杆件。
\end{solutioncard}

\begin{problemcard}{2-24 指定杆 1、2 的轴力}
计算图示桁架中 1、2 杆件的轴力。
\begin{figure}[H]
\centering
\begin{tikzpicture}[scale=0.9,>=Stealth]
  \coordinate (A) at (0,0);
  \coordinate (D) at (1.3,0);
  \coordinate (B) at (2.6,0);
  \coordinate (L) at (0,2.6);
  \coordinate (T) at (1.3,2.6);
  \coordinate (R) at (2.6,2.6);
  \coordinate (C) at (1.3,1.3);
  \draw[member] (A)--(B)--(R)--(L)--(A) (D)--(T) (A)--(C)--(R) (L)--(C)--(B);
  \foreach \p in {A,D,B,L,T,R,C} \node[smallpin] at (\p) {};
  \wallroller{0}{0}
  \roller{0}{0}
  \roller{2.6}{0}
  \draw[Brand,-{Stealth[length=2mm]}] (D)+(0,0.9)--(D)+(0,0.08) node[midway,right] {\(F_p\)};
  \node[above left] at (0.8,2.0) {1};
  \node[right] at (2.6,1.3) {2};
  \node[below left] at (A) {\(A\)};
  \node[below] at (B) {\(B\)};
  \node[right] at (C) {\(C\)};
  \draw[<->] (0,-0.45)--(1.3,-0.45) node[midway,below] {\(l\)};
  \draw[<->] (1.3,-0.45)--(2.6,-0.45) node[midway,below] {\(l\)};
  \draw[<->] (3.05,0)--(3.05,1.3) node[midway,right] {\(l\)};
  \draw[<->] (3.05,1.3)--(3.05,2.6) node[midway,right] {\(l\)};
\end{tikzpicture}
\end{figure}
\end{problemcard}

\begin{solutioncard}
\answer{\(F_{N1}=\dfrac{\sqrt2}{2}F_p,\qquad F_{N2}=-\dfrac12F_p\)。}

整体关于中线受竖向集中力，左右竖向反力相等：
\[
  F_{Ay}=F_{By}=\frac{F_p}{2}.
\]
先取左半部分截面，截过 1 杆。1 杆与水平成 \(45^\circ\)，由竖向平衡
\[
  F_{N1}\sin45^\circ-\frac{F_p}{2}=0,
\]
故
\[
  F_{N1}=\frac{F_p}{2\sin45^\circ}=\frac{\sqrt2}{2}F_p.
\]
再取右侧隔离体，对左端点取矩：
\[
  F_{N2}\cdot 2l+F_p\cdot l=0,
\]
于是
\[
  F_{N2}=-\frac12F_p.
\]
负号表示 2 杆受压。
\end{solutioncard}

\section{二、截面法求指定杆}

\begin{problemcard}{2-25 指定杆 1、2 的轴力}
求图示结构中指定杆件 1、2 的轴力。
\begin{figure}[H]
\centering
\begin{tikzpicture}[scale=0.82,>=Stealth]
  \coordinate (A) at (0,0);
  \coordinate (B) at (3,0);
  \coordinate (C) at (0,2);
  \coordinate (D) at (1,1);
  \coordinate (E) at (2,1);
  \coordinate (F) at (3,2);
  \draw[member] (A)--(B)--(F)--(C)--(A) (C)--(D)--(F) (A)--(E)--(B) (D)--(E) (E)--(B);
  \foreach \p in {A,B,C,D,E,F} \node[smallpin] at (\p) {};
  \wallroller{0}{0}
  \roller{0}{0}
  \roller{3}{0}
  \draw[Brand,-{Stealth[length=2mm]}] (-0.8,2)--(-0.05,2) node[midway,above] {\(50\,\mathrm{kN}\)};
  \draw[Brand,-{Stealth[length=2mm]}] (D)+(0,0.85)--(D)+(0,0.08) node[midway,right] {\(30\,\mathrm{kN}\)};
  \node[below] at (1.55,1.0) {1};
  \node[right] at (2.55,0.55) {2};
  \draw[<->] (0,-0.45)--(1,-0.45) node[midway,below] {\(l\)};
  \draw[<->] (1,-0.45)--(2,-0.45) node[midway,below] {\(l\)};
  \draw[<->] (2,-0.45)--(3,-0.45) node[midway,below] {\(l\)};
  \draw[<->] (3.35,0)--(3.35,1) node[midway,right] {\(l\)};
  \draw[<->] (3.35,1)--(3.35,2) node[midway,right] {\(l\)};
\end{tikzpicture}
\end{figure}
\end{problemcard}

\begin{solutioncard}
\answer{\(F_{N1}=-50\,\mathrm{kN},\qquad F_{N2}=-\dfrac{50\sqrt2}{3}\,\mathrm{kN}\)。}

取穿过 1 杆的上部截面。该截面上水平外力为 \(50\,\mathrm{kN}\)，1 杆为水平杆，因此由水平方向平衡直接得
\[
  50+F_{N1}=0,\qquad F_{N1}=-50\,\mathrm{kN}.
\]
再在节点附近取截面，2 杆为斜杆，其方向余弦可由图中几何关系确定。代入节点平衡并消去同截面上的非目标杆，得
\[
  F_{N2}=-\frac{50\sqrt2}{3}\,\mathrm{kN}.
\]
两个结果均为负，说明指定 1、2 杆均受压。
\end{solutioncard}

\begin{problemcard}{2-26 指定杆 1、2 的轴力}
求图示结构中指定杆件 1、2 的轴力。
\begin{figure}[H]
\centering
\begin{tikzpicture}[scale=0.82,>=Stealth]
  \coordinate (A) at (0,0);
  \coordinate (C) at (0,1.5);
  \coordinate (D) at (1.4,0);
  \coordinate (E) at (1.4,1.5);
  \coordinate (F) at (2.8,0);
  \coordinate (G) at (2.8,1.5);
  \coordinate (I) at (4.2,0);
  \coordinate (J) at (4.2,1.5);
  \coordinate (B) at (0.2,2.7);
  \draw[member] (A)--(D)--(F)--(I)--(J)--(G)--(E)--(C)--(A);
  \draw[member] (D)--(E) (F)--(G) (C)--(D)--(G)--(I) (D)--(G) (G)--(J) (B)--(J);
  \draw[member] (B)--(A);
  \foreach \p in {A,C,D,E,F,G,I,J,B} \node[smallpin] at (\p) {};
  \wallroller{0}{0}
  \roller{0}{0}
  \draw[Brand,-{Stealth[length=2mm]}] (D)+(0,0.65)--(D)+(0,0.05) node[midway,right] {\(F_p\)};
  \draw[Brand,-{Stealth[length=2mm]}] (I)+(0,0.65)--(I)+(0,0.05) node[midway,right] {\(F_p\)};
  \node[below left] at ($(C)!0.5!(D)$) {1};
  \node[above] at ($(B)!0.55!(J)$) {2};
  \node[left] at (B) {\(B\)};
  \node[right] at (J) {\(J\)};
  \draw[<->] (0,-0.45)--(1.4,-0.45) node[midway,below] {\(l\)};
  \draw[<->] (1.4,-0.45)--(2.8,-0.45) node[midway,below] {\(l\)};
  \draw[<->] (2.8,-0.45)--(4.2,-0.45) node[midway,below] {\(l\)};
  \draw[<->] (4.65,0)--(4.65,1.5) node[midway,right] {\(l\)};
  \draw[<->] (4.65,1.5)--(4.65,2.7) node[midway,right] {\(l\)};
\end{tikzpicture}
\end{figure}
\end{problemcard}

\begin{solutioncard}
\answer{\(F_{N1}=\dfrac{4\sqrt2}{3}F_p,\qquad F_{N2}=\dfrac{2\sqrt{10}}{3}F_p\)。}

先取右侧附属部分求传至主体的端部作用力，再用截面法穿过 1、2 杆。对左端支座点取矩，原答案笔记给出的控制方程为
\[
  F_{N2}\frac{3}{\sqrt{10}}\cdot 2l-F_p l-F_p\cdot 3l=0,
\]
故
\[
  F_{N2}=\frac{2\sqrt{10}}{3}F_p.
\]
再对同一隔离体列水平方向平衡，斜杆 2 的水平分量与 1 杆分量共同平衡，得
\[
  F_{N1}=\frac{4\sqrt2}{3}F_p.
\]
两杆均为正，表示按受拉为正的约定均受拉。
\end{solutioncard}

\section{三、多杆桁架的对称、零杆与截面法}

\begin{problemcard}{2-27 指定杆 1、2 的轴力}
求图示桁架中 1、2 杆的轴力。
\begin{figure}[H]
\centering
\begin{tikzpicture}[scale=0.72,>=Stealth]
  \coordinate (A) at (0,0);
  \coordinate (B) at (1,0);
  \coordinate (C) at (2,0);
  \coordinate (D) at (3,0);
  \coordinate (E) at (4,0);
  \coordinate (F) at (1,-1.3);
  \coordinate (G) at (2,-1.3);
  \coordinate (H) at (3,-1.3);
  \draw[member] (A)--(E) (F)--(H) (B)--(F)--(C)--(G)--(D)--(H)--(E);
  \draw[member] (A)--(F) (B)--(G) (G)--(D) (H)--(E) (F)--(C) (D)--(G);
  \foreach \p in {A,B,C,D,E,F,G,H} \node[smallpin] at (\p) {};
  \wallroller{0}{0}
  \roller{0}{0}
  \roller{4}{0}
  \draw[Brand,-{Stealth[length=2mm]}] (B)+(0,0.85)--(B)+(0,0.05) node[midway,right] {\(F_p\)};
  \draw[Brand,-{Stealth[length=2mm]}] (D)+(0,0.85)--(D)+(0,0.05) node[midway,right] {\(F_p\)};
  \draw[Brand,-{Stealth[length=2mm]}] (G)+(0,-0.85)--(G)+(0,-0.05) node[midway,right] {\(2F_p\)};
  \node[right] at ($(B)!0.5!(G)$) {1};
  \node[above] at ($(A)!0.55!(B)$) {2};
  \draw[<->] (0,-1.85)--(1,-1.85) node[midway,below] {\(d\)};
  \draw[<->] (1,-1.85)--(2,-1.85) node[midway,below] {\(d\)};
  \draw[<->] (2,-1.85)--(3,-1.85) node[midway,below] {\(d\)};
  \draw[<->] (3,-1.85)--(4,-1.85) node[midway,below] {\(d\)};
\end{tikzpicture}
\end{figure}
\end{problemcard}

\begin{solutioncard}
\answer{\(F_{N1}=-\sqrt2F_p,\qquad F_{N2}=0\)。}

结构和荷载关于中线对称。左端上弦 2 杆所在节点在水平方向没有需要它传递的外力，结合相邻两杆共线的零杆判据，得
\[
  F_{N2}=0.
\]
对左半跨作截面，截过 1 杆并利用中线对称条件，斜杆 1 的竖向分量平衡一侧集中力：
\[
  F_{N1}\sin45^\circ+F_p=0,
\]
故
\[
  F_{N1}=-\sqrt2F_p.
\]
负号表示 1 杆受压。
\end{solutioncard}

\begin{problemcard}{2-28 指定杆 \(a,b,c\) 的轴力}
作图示桁架中 \(a,b,c\) 三杆的轴力。
\begin{figure}[H]
\centering
\begin{tikzpicture}[scale=0.68,>=Stealth]
  \coordinate (A) at (0,2);
  \coordinate (B) at (0,1);
  \coordinate (C) at (1.3,2);
  \coordinate (D) at (1.3,1);
  \coordinate (E) at (2.6,2);
  \coordinate (F) at (2.6,1);
  \coordinate (G) at (3.9,2);
  \coordinate (H) at (3.9,1);
  \coordinate (I) at (2.6,0);
  \coordinate (J) at (2.6,-1);
  \coordinate (K) at (3.9,0);
  \coordinate (L) at (3.9,-1);
  \draw[member] (A)--(G)--(H)--(B)--(A) (C)--(D) (E)--(F) (I)--(J)--(L)--(K)--(I);
  \draw[member] (B)--(I)--(L) (B)--(J) (D)--(I) (F)--(K) (D)--(G) (A)--(D) (F)--(G);
  \foreach \p in {A,B,C,D,E,F,G,H,I,J,K,L} \node[smallpin] at (\p) {};
  \wallroller{0}{1}
  \wallroller{3.9}{2}
  \roller{2.6}{-1}
  \draw[Brand,-{Stealth[length=2mm]}] (C)+(0,0.65)--(C)+(0,0.05) node[midway,right] {\(F_p\)};
  \node[right] at ($(F)!0.5!(K)$) {\(a\)};
  \node[below left] at ($(B)!0.55!(J)$) {\(b\)};
  \node[above] at ($(C)!0.5!(E)$) {\(c\)};
  \draw[<->] (0,-1.45)--(1.3,-1.45) node[midway,below] {\(d\)};
  \draw[<->] (1.3,-1.45)--(2.6,-1.45) node[midway,below] {\(d\)};
  \draw[<->] (2.6,-1.45)--(3.9,-1.45) node[midway,below] {\(d\)};
  \draw[<->] (4.35,-1)--(4.35,0) node[midway,right] {\(d\)};
  \draw[<->] (4.35,0)--(4.35,1) node[midway,right] {\(d\)};
  \draw[<->] (4.35,1)--(4.35,2) node[midway,right] {\(d\)};
\end{tikzpicture}
\end{figure}
\end{problemcard}

\begin{solutioncard}
\answer{\(F_{Na}=-\dfrac{\sqrt2}{2}F_p,\quad F_{Nb}=\dfrac{\sqrt2}{2}F_p,\quad F_{Nc}=\dfrac12F_p\)。}

先由整体平衡求支座反力，再依次取三个截面。原答案提示的截面顺序为 \(I-I\)、\(II-II\)、\(III-III\)：由 \(I-I\) 截面按 \(x\) 轴投影求 \(F_{Na}\)，由 \(II-II\) 截面求 \(F_{Nc}\)，由 \(III-III\) 截面求 \(F_{Nb}\)。相应平衡式可整理为
\[
  F_{Na}\cos45^\circ+\frac12F_p=0,\qquad
  F_{Nc}-\frac12F_p=0,
\]
\[
  F_{Nb}\cos45^\circ-\frac12F_p=0.
\]
故得到上述三式。\(a\) 杆为受压，\(b,c\) 杆为受拉。
\end{solutioncard}

\begin{problemcard}{2-29 指定杆 \(a,b\) 的轴力}
试求图示桁架中 \(a,b\) 杆件的内力。
\begin{figure}[H]
\centering
\begin{tikzpicture}[scale=0.72,>=Stealth]
  \coordinate (A) at (0,0);
  \coordinate (B) at (1,1.2);
  \coordinate (C) at (2,2.4);
  \coordinate (D) at (3,1.2);
  \coordinate (E) at (4,0);
  \coordinate (F) at (1,0);
  \coordinate (G) at (2,0);
  \coordinate (H) at (3,0);
  \draw[member] (A)--(B)--(C)--(D)--(E) (A)--(E) (B)--(F)--(G)--(H)--(D) (B)--(G)--(D) (G)--(H);
  \foreach \p in {A,B,C,D,E,F,G,H} \node[smallpin] at (\p) {};
  \roller{0}{0}
  \roller{2}{0}
  \roller{4}{0}
  \wallroller{4}{0}
  \draw[Brand,-{Stealth[length=2mm]}] (F)+(0,0.65)--(F)+(0,0.05) node[midway,right] {\(F_p\)};
  \node[above] at ($(G)!0.5!(D)$) {\(a\)};
  \node[right] at ($(G)!0.55!(D)$) {\(b\)};
  \draw[<->] (0,-0.45)--(2,-0.45) node[midway,below] {\(2d\)};
  \draw[<->] (2,-0.45)--(4,-0.45) node[midway,below] {\(2d\)};
\end{tikzpicture}
\end{figure}
\end{problemcard}

\begin{solutioncard}
\answer{\(F_{Na}=-F_p,\qquad F_{Nb}=\dfrac{\sqrt2}{2}F_p\)。}

先对整体取矩求支座竖向反力，再截取含 \(a,b\) 杆的右半部分。由于 \(a\) 杆为水平受力杆，其轴力由水平方向平衡直接得到压缩值：
\[
  F_{Na}=-F_p.
\]
斜杆 \(b\) 与水平成 \(45^\circ\)，其竖向分量平衡截面处等效竖向力：
\[
  F_{Nb}\sin45^\circ-\frac12F_p=0,
\]
故
\[
  F_{Nb}=\frac{\sqrt2}{2}F_p.
\]
\end{solutioncard}

\begin{problemcard}{2-30 指定杆 \(a,b\) 的轴力}
求图示桁架中 \(a,b\) 杆件的内力，图中上部竖向荷载为 \(F_p=1\)。
\begin{figure}[H]
\centering
\begin{tikzpicture}[scale=0.68,>=Stealth]
  \coordinate (A) at (0,0);
  \coordinate (B) at (1,0);
  \coordinate (C) at (2,0);
  \coordinate (D) at (3,0);
  \coordinate (E) at (4,0);
  \coordinate (F) at (1,1);
  \coordinate (G) at (2,1);
  \coordinate (H) at (3,1);
  \coordinate (I) at (2,2);
  \draw[member] (A)--(E) (F)--(H) (A)--(F)--(I)--(H)--(E) (B)--(F) (C)--(G) (D)--(H) (B)--(G)--(D) (F)--(C)--(H);
  \foreach \p in {A,B,C,D,E,F,G,H,I} \node[smallpin] at (\p) {};
  \wallroller{0}{0}
  \roller{0}{0}
  \roller{3}{0}
  \draw[Brand,-{Stealth[length=2mm]}] (F)+(0,0.75)--(F)+(0,0.05) node[midway,left] {\(F_p=1\)};
  \node[above] at ($(C)!0.5!(D)$) {\(a\)};
  \node[below] at ($(C)!0.5!(H)$) {\(b\)};
  \draw[<->] (0,-0.45)--(4,-0.45) node[midway,below] {\(6\times1\)};
\end{tikzpicture}
\end{figure}
\end{problemcard}

\begin{solutioncard}
\answer{\(F_{Na}=0,\qquad F_{Nb}=\dfrac12F_p\)。若按图中 \(F_p=1\)，则 \(F_{Nb}=0.5\)。}

先判零杆：集中力右侧和悬臂部分的若干杆均为零杆。去掉零杆后，结构可化为原答案提示中的简化体系；由对称性或节点无外力条件可判定
\[
  F_{Na}=0.
\]
再对含 \(b\) 杆的节点列平衡。该节点剩余受力中，\(b\) 杆的轴力分量承担一半单位荷载，故
\[
  F_{Nb}=\frac12F_p.
\]
\end{solutioncard}

\begin{problemcard}{2-31 指定杆 \(a,b\) 的轴力}
计算图示桁架中 \(a,b\) 两杆的轴力。
\begin{figure}[H]
\centering
\begin{tikzpicture}[scale=0.72,>=Stealth]
  \coordinate (A) at (0,0);
  \coordinate (B) at (1,0);
  \coordinate (C) at (2,0);
  \coordinate (D) at (3,0);
  \coordinate (E) at (0,1.2);
  \coordinate (F) at (1,1.2);
  \coordinate (G) at (2,1.2);
  \coordinate (H) at (3,1.2);
  \draw[member] (A)--(D) (E)--(H) (A)--(E) (B)--(F) (C)--(G) (D)--(H);
  \draw[member] (A)--(F) (A)--(G) (A)--(H) (F)--(G)--(H);
  \foreach \p in {A,B,C,D,E,F,G,H} \node[smallpin] at (\p) {};
  \wallroller{0}{1.2}
  \roller{0}{0}
  \draw[Brand,-{Stealth[length=2mm]}] (B)+(0,0.75)--(B)+(0,0.05) node[midway,right] {\(F_p\)};
  \draw[Brand,-{Stealth[length=2mm]}] (C)+(0,0.75)--(C)+(0,0.05) node[midway,right] {\(F_p\)};
  \draw[Brand,-{Stealth[length=2mm]}] (D)+(0,0.75)--(D)+(0,0.05) node[midway,right] {\(F_p\)};
  \draw[Brand,-{Stealth[length=2mm]}] (D)+(0.05,0)--(D)+(0.75,0) node[midway,above] {\(F_p\)};
  \node[above] at ($(A)!0.5!(H)$) {\(a\)};
  \node[below] at ($(A)!0.5!(G)$) {\(b\)};
  \draw[<->] (0,-0.45)--(1,-0.45) node[midway,below] {\(l\)};
  \draw[<->] (1,-0.45)--(2,-0.45) node[midway,below] {\(l\)};
  \draw[<->] (2,-0.45)--(3,-0.45) node[midway,below] {\(l\)};
\end{tikzpicture}
\end{figure}
\end{problemcard}

\begin{solutioncard}
\answer{\(F_{Na}=-\sqrt5F_p,\qquad F_{Nb}=F_p\)。}

对右侧部分取截面，截过 \(a,b\) 及一根辅助杆。由几何关系，\(a\) 杆斜率对应方向余弦
\[
  \cos\alpha=\frac{2}{\sqrt5},\qquad \sin\alpha=\frac{1}{\sqrt5}.
\]
先对截面交点取矩消去辅助杆，可得 \(b\) 杆：
\[
  F_{Nb}=F_p.
\]
再列水平方向平衡，外力 \(F_p\) 与 \(b\) 杆分量已知，解得
\[
  F_{Na}=-\sqrt5F_p.
\]
负号表示 \(a\) 杆受压。
\end{solutioncard}

\begin{problemcard}{2-32 指定杆 1、2 的轴力}
求图示桁架中 1、2 杆件的轴力。
\begin{figure}[H]
\centering
\begin{tikzpicture}[scale=0.82,>=Stealth]
  \coordinate (A) at (0,0);
  \coordinate (B) at (1,1.4);
  \coordinate (C) at (2,2.4);
  \coordinate (D) at (3,1.4);
  \coordinate (E) at (4,0);
  \coordinate (L) at (0,1.4);
  \coordinate (R) at (4,1.4);
  \coordinate (M) at (2,1.0);
  \draw[member] (A)--(B)--(C)--(D)--(E) (A)--(E) (L)--(B) (D)--(R) (L)--(A) (R)--(E);
  \draw[member] (A)--(M)--(E) (B)--(M)--(D) (L)--(M)--(R) (C)--(M);
  \foreach \p in {A,B,C,D,E,L,R,M} \node[smallpin] at (\p) {};
  \wallroller{0}{0}
  \roller{0}{0}
  \roller{4}{0}
  \draw[Brand,-{Stealth[length=2mm]}] (L)+(-0.75,0)--(L)+(-0.05,0) node[midway,above] {\(F_p\)};
  \draw[Brand,-{Stealth[length=2mm]}] (R)+(0.75,0)--(R)+(0.05,0) node[midway,above] {\(F_p\)};
  \node[below left] at ($(B)!0.45!(M)$) {1};
  \node[left] at ($(C)!0.5!(M)$) {2};
  \node[below] at ($(A)!0.5!(M)$) {\(A\)};
  \draw[<->] (0,-0.45)--(2,-0.45) node[midway,below] {\(a\)};
  \draw[<->] (2,-0.45)--(4,-0.45) node[midway,below] {\(a\)};
  \draw[<->] (4.45,0)--(4.45,0.8) node[midway,right] {\(a/2\)};
  \draw[<->] (4.45,0.8)--(4.45,1.6) node[midway,right] {\(a/2\)};
\end{tikzpicture}
\end{figure}
\end{problemcard}

\begin{solutioncard}
\answer{\(F_{N1}=-\dfrac{\sqrt5}{2}F_p,\qquad F_{N2}=-F_p\)。}

原答案提示：先判断 \(A\) 点连接的两根斜杆为零杆，再用节点法或截面法。去掉零杆后，对穿过 1、2 杆的截面列平衡。由中部竖杆 2 的竖向平衡直接得
\[
  F_{N2}=-F_p.
\]
杆 1 的几何长度对应水平、竖向分量比为 \(2:1\)，故
\[
  \cos\alpha=\frac{2}{\sqrt5},\qquad \sin\alpha=\frac{1}{\sqrt5}.
\]
代入截面平衡可得
\[
  F_{N1}=-\frac{\sqrt5}{2}F_p.
\]
两杆均受压。
\end{solutioncard}

\begin{notecard}
第 5 次作业讲解笔记补充：桁架题不要急于列大方程。先用零杆规则删去明显不受力杆，再根据“只求少数杆”优先选截面法；若目标杆集中在某个节点附近，则用节点法更快。截面法的核心是选取矩心，让同一截面上尽量多的未知杆被消去。
\end{notecard}

\chapter{组合结构}

\begin{notecard}
组合结构同时含有受弯杆、桁架杆或链杆。求解时先把结构拆成“梁段、刚架段、二力杆段”，能单独平衡的部分先算；链杆只承受轴力，受弯杆则可能同时有轴力、剪力和弯矩，不能把梁式杆漏算成纯桁架杆。
\end{notecard}

\section{一、轴力判断与受弯杆内力}

\begin{problemcard}{2-33 杆 \(AB\) 的轴力}
图示结构中杆 \(AB\) 的轴力为多少？
\[
  \text{选项：}\quad
  \mathrm{A.}\ 0,\quad
  \mathrm{B.}\ -0.894F_p,\quad
  \mathrm{C.}\ -0.5F_p,\quad
  \mathrm{D.}\ -F_p.
\]
\begin{figure}[H]
\centering
\begin{tikzpicture}[scale=0.72,>=Stealth]
  \coordinate (L) at (0,0);
  \coordinate (P) at (1.4,0);
  \coordinate (A) at (2.8,0);
  \coordinate (R) at (4.2,0);
  \coordinate (S) at (5.6,0);
  \coordinate (T) at (2.8,1.4);
  \coordinate (B) at (2.8,-1.4);
  \draw[member] (L)--(S) (P)--(T)--(R) (P)--(B)--(R) (A)--(B);
  \foreach \p in {L,P,A,R,S,T,B} \node[smallpin] at (\p) {};
  \wallroller{0}{0}
  \roller{0}{0}
  \roller{5.6}{0}
  \draw[Brand,-{Stealth[length=2mm]}] (T)+(0,0.75)--(T)+(0,0.05) node[midway,right] {\(F_p\)};
  \node[right] at (A) {\(A\)};
  \node[right] at (B) {\(B\)};
  \draw[<->] (0,-1.85)--(1.4,-1.85) node[midway,below] {\(a\)};
  \draw[<->] (1.4,-1.85)--(2.8,-1.85) node[midway,below] {\(a\)};
  \draw[<->] (2.8,-1.85)--(4.2,-1.85) node[midway,below] {\(a\)};
  \draw[<->] (4.2,-1.85)--(5.6,-1.85) node[midway,below] {\(a\)};
\end{tikzpicture}
\end{figure}
\end{problemcard}

\begin{solutioncard}
\answer{答案：C，即 \(F_{NAB}=-0.5F_p\)。}

本题容易错在把梁式水平杆当成普通桁架杆。由整体对称性，左右支座竖向反力均为
\[
  F_{Ly}=F_{Sy}=\frac12F_p.
\]
上部三角形节点处两根斜杆方向对称，设斜杆轴力为 \(F_N\)，则
\[
  2F_N\sin45^\circ=F_p,\qquad F_N=\frac{F_p}{\sqrt2}.
\]
将中部 \(A\) 点附近隔离，竖向杆 \(AB\) 需平衡相邻斜杆传来的竖向分量，故
\[
  F_{NAB}=-F_N\sin45^\circ=-\frac12F_p.
\]
负号表示 \(AB\) 杆受压。
\end{solutioncard}

\begin{problemcard}{2-34 受弯杆弯矩图与链杆轴力}
对图示结构，作受弯杆件的弯矩图，并算出各链杆的轴力，标注在杆件旁边。
\begin{figure}[H]
\centering
\begin{tikzpicture}[scale=0.76,>=Stealth]
  \coordinate (A) at (0,0);
  \coordinate (B) at (1.4,1.4);
  \coordinate (C) at (2.8,1.4);
  \coordinate (D) at (4.2,0);
  \coordinate (E) at (1.4,2.8);
  \coordinate (F) at (2.8,2.8);
  \coordinate (G) at (1.4,4.2);
  \coordinate (H) at (2.8,4.2);
  \draw[member] (A)--(B)--(C)--(D) (B)--(E)--(G)--(H)--(F)--(C) (E)--(F) (B)--(F);
  \foreach \p in {A,B,C,D,E,F,G,H} \node[smallpin] at (\p) {};
  \wallroller{0}{0}
  \roller{0}{0}
  \roller{4.2}{0}
  \foreach \x in {1.55,1.85,...,2.65}
    \draw[Brand,-{Stealth[length=1.3mm]}] (\x,4.75)--(\x,4.25);
  \draw[Brand] (1.4,4.75)--(2.8,4.75);
  \node[above] at (2.1,4.75) {\(q\)};
  \draw[<->] (0,-0.45)--(1.4,-0.45) node[midway,below] {\(a\)};
  \draw[<->] (1.4,-0.45)--(2.8,-0.45) node[midway,below] {\(a\)};
  \draw[<->] (2.8,-0.45)--(4.2,-0.45) node[midway,below] {\(a\)};
  \draw[<->] (4.55,0)--(4.55,1.4) node[midway,right] {\(a\)};
  \draw[<->] (4.55,1.4)--(4.55,2.8) node[midway,right] {\(a\)};
  \draw[<->] (4.55,2.8)--(4.55,4.2) node[midway,right] {\(a\)};
\end{tikzpicture}
\end{figure}
\end{problemcard}

\begin{solutioncard}
先取顶层受弯杆，均布荷载合力为 \(qa\)，由对称性传到两侧竖杆的竖向作用为 \(qa/2\)。顶梁弯矩图为简支梁抛物线：
\[
  M_{\max}=\frac18qa^2.
\]
中层水平杆无横向荷载，且两端铰结传力，弯矩为零。两根竖向受弯杆受侧向链杆作用，其端部控制弯矩为
\[
  M=-\frac12qa.
\]
底部两根斜链杆与水平成 \(45^\circ\)，由节点竖向平衡得其轴力大小为
\[
  F_N=\frac{qa/2}{\sin45^\circ}=\frac{\sqrt2}{2}qa.
\]
\begin{figure}[H]
\centering
\begin{tikzpicture}[scale=0.76,>=Stealth]
  \begin{scope}
    \draw[member] (0,0)--(1.4,1.4)--(2.8,1.4)--(4.2,0) (1.4,1.4)--(1.4,4.2)--(2.8,4.2)--(2.8,1.4) (1.4,2.8)--(2.8,2.8);
    \draw[Brand,thick] (1.4,4.45) parabola bend (2.1,4.85) (2.8,4.45);
    \node[above] at (2.1,4.85) {\(qa^2/8\)};
    \node[left] at (1.15,3.5) {\(-qa/2\)};
    \node[right] at (3.05,3.5) {\(-qa/2\)};
    \node[above] at (2.1,2.8) {\(0\)};
    \node[below] at (2.1,-0.35) {\(M\) 图};
  \end{scope}
  \begin{scope}[shift={(6,0)}]
    \draw[member] (0,0)--(1.4,1.4)--(2.8,1.4)--(4.2,0);
    \node[above left] at (0.7,0.7) {\(\frac{\sqrt2}{2}qa\)};
    \node[above right] at (3.5,0.7) {\(\frac{\sqrt2}{2}qa\)};
    \node[above] at (2.1,1.4) {\(0\)};
    \node[below] at (2.1,-0.35) {链杆轴力};
  \end{scope}
\end{tikzpicture}
\end{figure}
\end{solutioncard}

\begin{problemcard}{2-35 弯矩图与杆 1 轴力}
对图示结构作弯矩图，并求杆件 1 的轴力。
\begin{figure}[H]
\centering
\begin{tikzpicture}[scale=0.78,>=Stealth]
  \draw[member] (0,0)--(2,0)--(4,0)--(4,2)--(2,2)--(2,0);
  \draw[member] (0,0)--(4,2);
  \wallroller{0}{0}
  \roller{4}{0}
  \foreach \x in {0.25,0.55,...,3.75}
    \draw[Brand,-{Stealth[length=1.2mm]}] (\x,2.55)--(\x,2.05);
  \draw[Brand] (0,2.55)--(4,2.55);
  \node[above] at (2,2.55) {\(q\)};
  \node[right] at (2,1.0) {1};
  \draw[<->] (0,-0.45)--(2,-0.45) node[midway,below] {\(l\)};
  \draw[<->] (2,-0.45)--(4,-0.45) node[midway,below] {\(l\)};
  \draw[<->] (4.4,0)--(4.4,2) node[midway,right] {\(l\)};
\end{tikzpicture}
\end{figure}
\end{problemcard}

\begin{solutioncard}
\answer{\(F_{N1}=-2ql\)。}

顶梁均布荷载合力为 \(2ql\)，由组合结构传至中竖杆。杆 1 为竖向链杆，只传轴力；由节点竖向平衡可得
\[
  F_{N1}+2ql=0,\qquad F_{N1}=-2ql.
\]
受弯杆的弯矩图控制值按答案图整理为
\[
  M_{\text{左端}}=2ql^2,\qquad M_{\text{顶梁中点}}=\frac12ql^2.
\]
\begin{figure}[H]
\centering
\begin{tikzpicture}[scale=0.78,>=Stealth]
  \draw[member] (0,0)--(2,0)--(4,0)--(4,2)--(2,2)--(2,0) (0,0)--(4,2);
  \draw[Brand,thick] (0,0.9)--(2,0)--(4,0) (2,2.4) parabola bend (3,2.75) (4,2.4);
  \node[left] at (0,0.9) {\(2ql^2\)};
  \node[above] at (3,2.75) {\(ql^2/2\)};
  \node[right] at (2,1.0) {\(F_{N1}=-2ql\)};
  \node[below] at (2,-0.35) {\(M\) 图与杆 1 轴力};
\end{tikzpicture}
\end{figure}
\end{solutioncard}

\section{二、内力图综合}

\begin{problemcard}{2-36 弯矩图、剪力图与轴力图}
试作图示结构的弯矩图、剪力图、轴力图。
\begin{figure}[H]
\centering
\begin{tikzpicture}[scale=0.75,>=Stealth]
  \draw[member] (0,0)--(0,3)--(3,3)--(3,0);
  \draw[member] (0,0)--(3,3);
  \wallroller{0}{0}
  \wallroller{0}{3}
  \roller{0}{0}
  \draw[Brand,-{Stealth[length=2mm]}] (3,4.0)--(3,3.1) node[midway,right] {\(20\,\mathrm{kN}\)};
  \foreach \x in {0.35,0.75,...,2.65}
    \draw[Brand,-{Stealth[length=1.2mm]}] (\x,3.55)--(\x,3.05);
  \draw[Brand] (0.2,3.55)--(2.8,3.55);
  \node[above] at (1.5,3.55) {\(20\,\mathrm{kN/m}\)};
  \draw[<->] (0,-0.45)--(3,-0.45) node[midway,below] {\(4m+2m\)};
  \draw[<->] (3.4,0)--(3.4,1.5) node[midway,right] {\(3m\)};
  \draw[<->] (3.4,1.5)--(3.4,3) node[midway,right] {\(3m\)};
\end{tikzpicture}
\end{figure}
\end{problemcard}

\begin{solutioncard}
按答案图，弯矩、剪力和轴力主要控制值为：
\[
  M:\ 40,\ 140\quad(\mathrm{kN\cdot m});
\]
\[
  F_Q:\ \frac{140}{3},\ 30,\ 50,\ 20\quad(\mathrm{kN});
\]
\[
  F_N:\ 30,\ \frac{280}{3},\ \frac{350}{3}\quad(\mathrm{kN}).
\]
\begin{figure}[H]
\centering
\begin{tikzpicture}[scale=0.78,>=Stealth]
  \begin{scope}
    \draw[member] (0,0)--(0,3)--(3,3)--(3,0) (0,0)--(3,3);
    \draw[Brand,thick] (0.25,0)--(0.25,3)--(1.6,2.55)--(3.2,3);
    \node[above] at (1.6,2.55) {\(140\)};
    \node[above] at (2.4,3.22) {\(40\)};
    \node[below] at (1.5,-0.35) {\(M\) 图};
  \end{scope}
  \begin{scope}[shift={(5.2,0)}]
    \draw[member] (0,0)--(0,3)--(3,3)--(3,0) (0,0)--(3,3);
    \node[left] at (0,2.2) {\(140/3\)};
    \node[above] at (1.5,3) {\(30\)};
    \node[right] at (3,2.2) {\(20\)};
    \node[below] at (1.5,-0.35) {\(F_Q,F_N\) 控制值};
  \end{scope}
\end{tikzpicture}
\end{figure}
检查时应把斜链杆作为二力杆处理，受弯杆则按梁段分别作 \(M,Q,N\) 图。
\end{solutioncard}

\begin{problemcard}{2-37 作弯矩图}
作图示组合结构的弯矩图。
\begin{figure}[H]
\centering
\begin{tikzpicture}[scale=0.72,>=Stealth]
  \coordinate (A) at (0,0);
  \coordinate (B) at (4,0);
  \coordinate (D) at (1,0.25);
  \coordinate (E) at (1,1.7);
  \coordinate (G) at (1,3);
  \coordinate (F) at (2.6,3);
  \coordinate (C) at (2.6,0.5);
  \draw[member] (A)--(B) (D)--(G)--(F) (D)--(E)--(F) (E)--(C) (G)--(D);
  \foreach \p in {A,B,D,E,G,F,C} \node[smallpin] at (\p) {};
  \roller{0}{0}
  \roller{4}{0}
  \draw[Brand,->,thick] (0.7,1.8) arc[start angle=120,end angle=-240,radius=0.35];
  \node[left] at (0.55,1.7) {\(F_pl\)};
  \node[below] at (A) {\(A\)};
  \node[below] at (B) {\(B\)};
\end{tikzpicture}
\end{figure}
\end{problemcard}

\begin{solutioncard}
答案弯矩图的主要控制值为底梁跨中下侧
\[
  M=\frac{2F_pl}{3}.
\]
上部桁架链杆只传轴力，弯矩图仅画在受弯梁杆上。
\begin{figure}[H]
\centering
\begin{tikzpicture}[scale=0.72,>=Stealth]
  \draw[member] (0,0)--(4,0) (1,0.25)--(1,3)--(2.6,3) (1,1.7)--(2.6,3) (1,1.7)--(2.6,0.5);
  \draw[Brand,thick] (0,0)--(2,-0.65)--(4,0);
  \node[below] at (2,-0.65) {\(\frac{2F_pl}{3}\)};
  \node[below] at (2,-1.05) {\(M\) 图};
\end{tikzpicture}
\end{figure}
\end{solutioncard}

\begin{problemcard}{2-38 作弯矩图}
作图示结构的弯矩图。
\begin{figure}[H]
\centering
\begin{tikzpicture}[scale=0.72,>=Stealth]
  \draw[member] (0,0)--(0,3)--(2,3)--(4,3)--(4,0);
  \draw[member] (0,0)--(2,3) (0,1.4)--(2,3) (0,0)--(4,3) (4,0)--(4.8,-0.4);
  \wallroller{0}{0}
  \fixedbase{4}{0}
  \node[smallpin] at (0,0) {};
  \node[smallpin] at (0,1.4) {};
  \node[smallpin] at (0,3) {};
  \node[smallpin] at (2,3) {};
  \node[smallpin] at (4,3) {};
  \draw[<->] (0.3,1.4)--(1.3,2.3) node[midway,right] {\(4\)};
  \draw[<->] (0.3,0)--(1.6,2.0) node[midway,below] {\(24\)};
\end{tikzpicture}
\end{figure}
\end{problemcard}

\begin{solutioncard}
答案图中弯矩控制值为
\[
  4,\qquad 24,\qquad 32,\qquad 48\quad(\mathrm{kN\cdot m}).
\]
\begin{figure}[H]
\centering
\begin{tikzpicture}[scale=0.72,>=Stealth]
  \draw[member] (0,0)--(0,3)--(2,3)--(4,3)--(4,0);
  \draw[Brand,thick] (0.2,0)--(0.8,1.4)--(2,2.35)--(4,3.45)--(4.65,0);
  \node[left] at (0.8,1.4) {\(24\)};
  \node[above] at (2,2.35) {\(32\)};
  \node[right] at (4.65,0) {\(48\)};
  \node[below] at (2,-0.35) {\(M\) 图};
\end{tikzpicture}
\end{figure}
\end{solutioncard}

\begin{problemcard}{2-39 作弯矩图并求桁架杆轴力}
作图示结构的弯矩图，并求桁架杆件的轴力。
\begin{figure}[H]
\centering
\begin{tikzpicture}[scale=0.72,>=Stealth]
  \coordinate (A) at (0,0);
  \coordinate (C) at (0,1.5);
  \coordinate (D) at (0,3);
  \coordinate (B) at (2,3.1);
  \draw[member] (A)--(D)--(B) (A)--(B) (C)--(B);
  \foreach \p in {A,C,D,B} \node[smallpin] at (\p) {};
  \wallroller{0}{0}
  \roller{2}{3.1}
  \draw[Brand,->,thick] (0.05,0.35) arc[start angle=180,end angle=-120,radius=0.33];
  \node[left] at (0,0.4) {\(2F_pl\)};
  \draw[Brand,->,thick] (0.05,2.7) arc[start angle=180,end angle=-120,radius=0.33];
  \node[left] at (0,2.7) {\(F_pl\)};
  \node[below] at (A) {\(A\)};
  \node[right] at (B) {\(B\)};
\end{tikzpicture}
\end{figure}
\end{problemcard}

\begin{solutioncard}
答案弯矩图如图示，斜向 \(AB\) 杆所在桁架杆轴力为零：
\[
  F_{N,AB}=0.
\]
底部受弯杆端部控制力偶为 \(2F_pl\)，上端控制力偶为 \(F_pl\)。
\begin{figure}[H]
\centering
\begin{tikzpicture}[scale=0.72,>=Stealth]
  \draw[member] (0,0)--(0,3)--(2,3.1) (0,0)--(2,3.1) (0,1.5)--(2,3.1);
  \draw[Brand,thick] (0.32,0)--(0.32,3)--(2,3.1);
  \node[left] at (0.32,0.5) {\(2F_pl\)};
  \node[left] at (0.32,2.5) {\(F_pl\)};
  \node[right] at (1.25,2.05) {\(0\)};
  \node[below] at (1,-0.35) {\(M\) 图};
\end{tikzpicture}
\end{figure}
\end{solutioncard}

\begin{notecard}
第 6 次作业讲解笔记补充：组合结构的关键是“谁能弯、谁只能拉压”。链杆、桁架杆和两端铰接无横向荷载的杆只传轴力；连续梁段和刚接杆件必须画弯矩图。作图时先把二力杆替换为节点力，再对受弯杆列平衡。
\end{notecard}

\chapter{三铰拱}

\begin{notecard}
三铰拱在竖向荷载作用下，通常先按同跨度简支梁求竖向反力，再利用中间铰弯矩为零求水平推力。任意截面的弯矩可写成
\[
  M=M^{0}-Hy,
\]
其中 \(M^{0}\) 是相应简支梁在同截面处的弯矩，\(H\) 为水平推力，\(y\) 为该截面到两拱脚连线的竖向高度。求剪力和轴力时，应把截面处的合力沿拱轴切线与法线方向分解。
\end{notecard}

\section{三铰拱水平推力与截面内力}

\begin{problemcard}{2-40 截面 \(K\) 的弯矩}
图 2-188 所示三铰拱，拱轴线方程为
\[
  y=\frac{4f}{l^2}(lx-x^2).
\]
求截面 \(K\) 的弯矩，规定下拉为正。（广西大学 2010）
\begin{figure}[H]
\centering
\begin{tikzpicture}[scale=0.58,>=Stealth]
  \coordinate (A) at (0,0);
  \coordinate (C) at (6,4);
  \coordinate (K) at (10,2.222);
  \coordinate (B) at (12,0);
  \draw[member,domain=0:12,samples=100] plot (\x,{4/144*(12*\x-\x*\x)});
  \node[smallpin] at (A) {};
  \node[smallpin] at (C) {};
  \node[smallpin] at (K) {};
  \node[smallpin] at (B) {};
  \wallroller{0}{0}
  \roller{0}{0}
  \roller{12}{0}
  \wallroller{12.55}{0}
  \node[below left] at (A) {\(A\)};
  \node[above] at (C) {\(C\)};
  \node[above right] at (K) {\(K\)};
  \node[below right] at (B) {\(B\)};
  \foreach \x in {0.4,1.2,...,5.6}
    \draw[Brand,-{Stealth[length=1.4mm]}] (\x,5.25)--(\x,4.55);
  \draw[Brand] (0.2,5.25)--(5.8,5.25);
  \node[above] at (3,5.25) {\(q=20\,\mathrm{kN/m}\)};
  \draw[Brand,-{Stealth[length=2mm]}] (K)+(0,1.45)--(K)+(0,0.18) node[midway,right] {\(48\,\mathrm{kN}\)};
  \draw[<->] (0,-0.85)--(6,-0.85) node[midway,below] {\(6\,\mathrm{m}\)};
  \draw[<->] (6,-0.85)--(12,-0.85) node[midway,below] {\(6\,\mathrm{m}\)};
  \draw[<->] (0,-1.45)--(10,-1.45) node[midway,below] {\(10\,\mathrm{m}\)};
  \draw[<->] (12.9,0)--(12.9,4) node[midway,right] {\(4\,\mathrm{m}\)};
  \figtitle{图 2-188：三铰拱截面 \(K\)}
\end{tikzpicture}
\end{figure}
\end{problemcard}

\begin{solutioncard}
\answer{答案：\(M_K=13.33\,\mathrm{kN\cdot m}\)。}

本题先求右拱脚的竖向反力和水平反力。左半跨均布荷载合力为
\[
  q\cdot6=20\times6=120\,\mathrm{kN},
\]
作用点距 \(A\) 为 \(3\,\mathrm{m}\)。对 \(A\) 点取矩：
\[
  12F_{By}-48\times10-120\times3=0,
\]
故
\[
  F_{By}=70\,\mathrm{kN}.
\]
再利用中铰 \(C\) 的弯矩为零。取右半拱，关于 \(C\) 点取矩：
\[
  4F_{Bx}+48\times4-6F_{By}=0,
\]
代入 \(F_{By}=70\,\mathrm{kN}\)，得
\[
  F_{Bx}=57\,\mathrm{kN}.
\]
截面 \(K\) 的横坐标为 \(x_K=10\,\mathrm{m}\)，拱高为
\[
  y_K=\frac{4\times4}{12^2}(12\times10-10^2)
      =\frac{20}{9}\,\mathrm{m}.
\]
取 \(K\) 到 \(B\) 的右段隔离体，对 \(K\) 点取矩：
\[
  M_K+F_{Bx}y_K-F_{By}(12-10)=0.
\]
因此
\[
  M_K=70\times2-57\times\frac{20}{9}
      =\frac{40}{3}
      =13.33\,\mathrm{kN\cdot m}.
\]
\end{solutioncard}

\begin{problemcard}{2-41 水平支座反力}
图 2-189 所示三铰拱的水平支座反力是多少？（武汉大学 2010）
\begin{figure}[H]
\centering
\begin{tikzpicture}[scale=0.64,>=Stealth]
  \coordinate (A) at (0,0);
  \coordinate (P) at (3,1.5);
  \coordinate (C) at (6,2);
  \coordinate (B) at (12,0);
  \draw[member,domain=0:12,samples=100] plot (\x,{2/36*(12*\x-\x*\x)});
  \foreach \p in {A,C,B} \node[smallpin] at (\p) {};
  \wallroller{0}{0}
  \roller{0}{0}
  \roller{12}{0}
  \wallroller{12.55}{0}
  \draw[Brand,-{Stealth[length=2mm]}] (P)+(0,1.2)--(P)+(0,0.1) node[midway,left] {\(20\,\mathrm{kN}\)};
  \draw[Brand,-{Stealth[length=2mm]}] (C)+(0,1.3)--(C)+(0,0.12) node[midway,right] {\(20\,\mathrm{kN}\)};
  \draw[<->] (0,-0.65)--(3,-0.65) node[midway,below] {\(3\,\mathrm{m}\)};
  \draw[<->] (3,-0.65)--(6,-0.65) node[midway,below] {\(3\,\mathrm{m}\)};
  \draw[<->] (6,-0.65)--(12,-0.65) node[midway,below] {\(6\,\mathrm{m}\)};
  \draw[<->] (12.9,0)--(12.9,2) node[midway,right] {\(2\,\mathrm{m}\)};
  \node[below left] at (A) {\(A\)};
  \node[above] at (C) {\(C\)};
  \node[below right] at (B) {\(B\)};
  \figtitle{图 2-189：两集中荷载三铰拱}
\end{tikzpicture}
\end{figure}
\end{problemcard}

\begin{solutioncard}
\answer{答案：水平支座反力为 \(45\,\mathrm{kN}\)。}

先由整体竖向平衡求右支座竖向反力。对 \(A\) 点取矩：
\[
  12F_{By}-20\times3-20\times6=0,
\]
因此
\[
  F_{By}=15\,\mathrm{kN},\qquad F_{Ay}=40-15=25\,\mathrm{kN}.
\]
再取左半拱 \(A\sim C\)。中铰 \(C\) 处弯矩为零，关于 \(C\) 点取矩：
\[
  2H-25\times6+20\times3=0.
\]
于是
\[
  H=\frac{150-60}{2}=45\,\mathrm{kN}.
\]
左右拱脚水平反力大小相等、方向相反，题问水平支座反力的大小即为 \(45\,\mathrm{kN}\)。
\end{solutioncard}

\begin{problemcard}{2-42 链杆 \(AB\) 的轴力}
图 2-190 所示结构中，链杆 \(AB\) 的轴力为多少？以受拉为正。（中南大学 2010）
\begin{figure}[H]
\centering
\begin{tikzpicture}[scale=0.78,>=Stealth]
  \coordinate (L) at (0,0);
  \coordinate (A) at (2.15,2);
  \coordinate (T) at (4,3);
  \coordinate (B) at (5.85,2);
  \coordinate (R) at (8,0);
  \draw[member,domain=0:8,samples=100] plot (\x,{3/16*(8*\x-\x*\x)});
  \draw[member] (A)--(B);
  \draw[member] (T)--(4,2.55);
  \foreach \p in {L,A,T,B,R} \node[smallpin] at (\p) {};
  \wallroller{0}{0}
  \roller{0}{0}
  \roller{8}{0}
  \wallroller{8.45}{0}
  \draw[Brand,-{Stealth[length=2mm]}] (T)+(0,1.1)--(T)+(0,0.12) node[midway,left] {\(F_p\)};
  \node[left] at (A) {\(A\)};
  \node[right] at (B) {\(B\)};
  \draw[<->] (0,-0.65)--(4,-0.65) node[midway,below] {\(l/2\)};
  \draw[<->] (4,-0.65)--(8,-0.65) node[midway,below] {\(l/2\)};
  \draw[<->] (8.7,0)--(8.7,2) node[midway,right] {\(l/4\)};
  \draw[<->] (8.7,2)--(8.7,3) node[midway,right] {\(l/4\)};
  \figtitle{图 2-190：含水平链杆的三铰拱}
\end{tikzpicture}
\end{figure}
\end{problemcard}

\begin{solutioncard}
\answer{答案：\(F_{NAB}=-F_p\)，即链杆 \(AB\) 受压。}

结构和荷载关于中线对称，故两拱脚竖向反力相等：
\[
  F_{Ly}=F_{Ry}=\frac{F_p}{2}.
\]
取右半结构，由右拱脚到跨中铰的水平距离为 \(l/2\)。对左侧拱脚或由整体对称可知，右拱脚竖向反力为 \(F_p/2\)。

再取右半拱中包含链杆 \(AB\) 的隔离体，对右拱脚 \(R\) 点取矩。链杆 \(AB\) 的作用线距 \(R\) 的竖向距离为 \(l/4\)，跨中竖向荷载在右半隔离体中对应的力臂为 \(l/2\)，于是
\[
  F_{NAB}\cdot\frac{l}{4}
  +\frac{F_p}{2}\cdot\frac{l}{2}=0.
\]
化简得
\[
  F_{NAB}=-F_p.
\]
按题设受拉为正，负号表示链杆 \(AB\) 受压。
\end{solutioncard}

\begin{problemcard}{2-43 截面 \(D\) 的弯矩与剪力}
图 2-191 所示三铰拱，拱轴线方程为
\[
  y=\frac{4f}{l^2}(lx-x^2),
\]
其中 \(f=4\,\mathrm{m}\)，\(l=16\,\mathrm{m}\)。求截面 \(D\) 的弯矩 \(M_D\) 与剪力 \(F_{QD}\)。（浙江大学 2012）
\begin{figure}[H]
\centering
\begin{tikzpicture}[scale=0.45,>=Stealth]
  \coordinate (A) at (0,0);
  \coordinate (C) at (8,4);
  \coordinate (D) at (12,3);
  \coordinate (B) at (16,0);
  \draw[member,domain=0:16,samples=120] plot (\x,{1*\x-\x*\x/16});
  \foreach \p in {A,C,D,B} \node[smallpin] at (\p) {};
  \wallroller{0}{0}
  \roller{0}{0}
  \roller{16}{0}
  \wallroller{16.7}{0}
  \foreach \x in {0.8,2.0,...,7.2}
    \draw[Brand,-{Stealth[length=1.4mm]}] (\x,5.4)--(\x,4.55);
  \draw[Brand] (0.4,5.4)--(7.6,5.4);
  \node[above] at (4,5.4) {\(q\)};
  \node[below left] at (A) {\(A\)};
  \node[above] at (C) {\(C\)};
  \node[above right] at (D) {\(D\)};
  \node[below right] at (B) {\(B\)};
  \draw[<->] (0,-0.8)--(8,-0.8) node[midway,below] {\(8\,\mathrm{m}\)};
  \draw[<->] (8,-0.8)--(12,-0.8) node[midway,below] {\(4\,\mathrm{m}\)};
  \draw[<->] (12,-0.8)--(16,-0.8) node[midway,below] {\(4\,\mathrm{m}\)};
  \draw[<->] (17.3,0)--(17.3,3) node[midway,right] {\(3\,\mathrm{m}\)};
  \draw[<->] (17.3,3)--(17.3,4) node[midway,right] {\(1\,\mathrm{m}\)};
  \figtitle{图 2-191：抛物线三铰拱截面 \(D\)}
\end{tikzpicture}
\end{figure}
\end{problemcard}

\begin{solutioncard}
\answer{答案：\(M_D=4q\)（上侧受拉），\(F_{QD}=0\)。}

左半跨均布荷载合力为 \(8q\)，作用点距 \(A\) 为 \(4\,\mathrm{m}\)。对 \(A\) 点取矩：
\[
  16F_{By}-8q\times4=0,
\]
得
\[
  F_{By}=2q.
\]
由竖向平衡还可得
\[
  F_{Ay}=8q-2q=6q.
\]
再取中铰 \(C\) 的弯矩为零。对左半拱 \(A\sim C\) 关于 \(C\) 点取矩：
\[
  4H-6q\times8+8q\times4=0,
\]
因此
\[
  H=4q.
\]
按右半拱受力方向取截面，支座 \(B\) 处反力为
\[
  F_{Bx}=4q,\qquad F_{By}=2q,
\]
与讲解笔记中的右段隔离体一致。

截面 \(D\) 位于 \(x=12\,\mathrm{m}\)，由拱轴线
\[
  y=x-\frac{x^2}{16}
\]
得
\[
  y_D=12-\frac{12^2}{16}=3\,\mathrm{m}.
\]
取右段 \(D\sim B\)，对 \(D\) 点取矩：
\[
  M_D+4q\times3-2q\times4=0.
\]
若按代数弯矩取下侧受拉为正，则 \(M_D=-4q\)。题图答案按受拉侧标注，因此写为
\[
  M_D=4q,
\]
且标注为上侧受拉。

最后求剪力。拱轴线导数为
\[
  y'=1-\frac{2x}{16}=1-\frac{x}{8},
\]
在 \(D\) 点
\[
  y'_D=1-\frac{12}{8}=-\frac12.
\]
因此截面切线方向的斜率大小为 \(1/2\)，方向余弦可取
\[
  \cos\theta=\frac{2}{\sqrt5},\qquad \sin\theta=\frac{1}{\sqrt5}.
\]
右段外力合力 \((4q,\,2q)\) 与截面切线平行，其法向分量为零，所以
\[
  F_{QD}=0.
\]
\end{solutioncard}

\begin{notecard}
第 7 次作业讲解笔记补充：三铰拱题应优先找“中铰”。中铰不传弯矩，常用来求水平推力；截面弯矩则可用 \(M=M^0-Hy\) 快速复核。若要求剪力和轴力，不要直接把竖向剪力当作 \(F_Q\)，必须按截面切线方向重新分解。
\end{notecard}

\chapter{静定结构影响线}

\begin{notecard}
静定结构影响线可用静力法或机动法绘制。对梁式结构，剪力影响线在截面处常有突变，弯矩影响线在截面处通常形成折点；移动均布荷载要求把荷载布置在影响线同号面积最大的区段上，移动集中荷载组则比较各集中力落在影响线峰值附近时的控制组合。
\end{notecard}

\section{移动荷载与梁的影响线}

\begin{problemcard}{3-1 截面 \(K\) 的最大剪力}
图 3-48 所示结构在任意长度可移动均布荷载 \(20\,\mathrm{kN/m}\) 作用下，求截面 \(K\) 的剪力最大值 \(F_{QK,\max}\)。（天津大学 2010）
\begin{figure}[H]
\centering
\begin{tikzpicture}[scale=0.55,>=Stealth]
  \coordinate (A) at (0,0);
  \coordinate (K) at (4,0);
  \coordinate (H1) at (5,0);
  \coordinate (B) at (13,0);
  \coordinate (H2) at (15,0);
  \coordinate (C) at (21,0);
  \draw[member] (A)--(C);
  \draw[ground] (A)+(0,-0.55)--+(0,0.55);
  \foreach \y in {-0.45,-0.25,-0.05,0.15,0.35}
    \draw[ground] ($(A)+(0,\y)$)--($(A)+(-0.22,\y-0.16)$);
  \node[smallpin] at (K) {};
  \node[smallpin] at (H1) {};
  \node[smallpin] at (B) {};
  \node[smallpin] at (H2) {};
  \node[smallpin] at (C) {};
  \roller{13}{0}
  \roller{21}{0}
  \node[above] at (K) {\(K\)};
  \foreach \x in {7,8,9,10,11,12}
    \draw[Brand,-{Stealth[length=1.2mm]}] (\x,1.15)--(\x,0.15);
  \draw[Brand] (6.6,1.15)--(12.4,1.15);
  \node[above] at (9.5,1.15) {\(20\,\mathrm{kN/m}\)};
  \draw[<->] (0,-0.8)--(4,-0.8) node[midway,below] {\(4\,\mathrm{m}\)};
  \draw[<->] (4,-0.8)--(5,-0.8) node[midway,below] {\(1\,\mathrm{m}\)};
  \draw[<->] (5,-0.8)--(13,-0.8) node[midway,below] {\(8\,\mathrm{m}\)};
  \draw[<->] (13,-0.8)--(15,-0.8) node[midway,below] {\(2\,\mathrm{m}\)};
  \draw[<->] (15,-0.8)--(21,-0.8) node[midway,below] {\(6\,\mathrm{m}\)};
  \figtitle{图 3-48：多跨静定梁及截面 \(K\)}
\end{tikzpicture}
\end{figure}
\end{problemcard}

\begin{solutioncard}
\answer{答案：\(F_{QK,\max}=100\,\mathrm{kN}\)。}

用机动法去掉截面 \(K\) 处的剪力约束，并令相对竖向位移方向与正剪力方向一致。由于 \(K\) 左侧到固定端为刚片，\(K\) 右侧经内部铰与支座形成折线，影响线在 \(K\) 右侧的正值区可整理为一个三角形。其控制底长为
\[
  1+9=10\,\mathrm{m},
\]
最大竖标为 \(1\)。为了使剪力取最大正值，任意长度均布荷载应只布置在该正值三角形区域上。因此
\[
  F_{QK,\max}
  =q\cdot A_{+}
  =20\times\frac12\times10\times1
  =100\,\mathrm{kN}.
\]
\begin{figure}[H]
\centering
\begin{tikzpicture}[scale=0.56,>=Stealth]
  \draw[member] (0,0)--(21,0);
  \foreach \x in {4,5,13,15,21} \node[smallpin] at (\x,0) {};
  \roller{13}{0}
  \roller{21}{0}
  \draw[Brand,thick] (4,0)--(5,1)--(14,0);
  \fill[Brand!10] (4,0)--(5,1)--(14,0)--cycle;
  \node[above] at (5,1) {\(1\)};
  \node[below] at (4,0) {\(K\)};
  \draw[<->] (4,-0.75)--(5,-0.75) node[midway,below] {\(1\)};
  \draw[<->] (5,-0.75)--(14,-0.75) node[midway,below] {\(9\)};
  \node[below] at (10.5,-1.15) {\(F_{QK}\) 正值影响线};
\end{tikzpicture}
\end{figure}
\end{solutioncard}

\begin{problemcard}{3-2 截面 \(K\) 的最大弯矩}
仍取图 3-48 所示结构，在任意长度可移动均布荷载 \(20\,\mathrm{kN/m}\) 作用下，求截面 \(K\) 的最大弯矩绝对值 \(|M_{K,\max}|\)。（华南理工大学 2013）
\end{problemcard}

\begin{solutioncard}
\answer{答案：\(|M_{K,\max}|=90\,\mathrm{kN\cdot m}\)。}

去掉截面 \(K\) 处弯矩约束，令两侧刚片发生与正弯矩一致的相对转角。由于左端固定，弯矩影响线在固定端一侧为零；在 \(K\) 右侧形成负值三角形控制区。由答案图的控制竖标和长度，该三角形面积为
\[
  A_M=\frac12\times9\times1=\frac92.
\]
任意长度均布荷载应布满该同号区域，所以
\[
  |M_{K,\max}|
  =qA_M
  =20\times\frac92
  =90\,\mathrm{kN\cdot m}.
\]
\begin{figure}[H]
\centering
\begin{tikzpicture}[scale=0.56,>=Stealth]
  \draw[member] (0,0)--(21,0);
  \foreach \x in {4,5,13,15,21} \node[smallpin] at (\x,0) {};
  \roller{13}{0}
  \roller{21}{0}
  \draw[Brand,thick] (4,0)--(5,-1)--(14,0);
  \fill[Brand!10] (4,0)--(5,-1)--(14,0)--cycle;
  \node[below] at (5,-1) {\(-1\)};
  \node[below] at (4,0) {\(K\)};
  \draw[<->] (5,-1.55)--(14,-1.55) node[midway,below] {\(9\,\mathrm{m}\)};
  \node[below] at (10.5,-2.0) {\(M_K\) 控制影响线};
\end{tikzpicture}
\end{figure}
\end{solutioncard}

\begin{problemcard}{3-3 截面 \(K\) 的剪力影响线与移动集中荷载}
图 3-49 所示多跨梁，求截面 \(K\) 的剪力影响线在 \(K\) 点的竖标 \(y_{K\text{左}}\)、\(y_{K\text{右}}\)；并求图示移动荷载作用下截面 \(K\) 的剪力最大值。（浙江大学 2009）
\begin{figure}[H]
\centering
\begin{tikzpicture}[scale=0.52,>=Stealth]
  \coordinate (A) at (0,0);
  \coordinate (H1) at (5,0);
  \coordinate (B) at (13,0);
  \coordinate (K) at (15,0);
  \coordinate (H2) at (17,0);
  \coordinate (C) at (25,0);
  \draw[member] (A)--(C);
  \draw[ground] (A)+(0,-0.55)--+(0,0.55);
  \foreach \y in {-0.45,-0.25,-0.05,0.15,0.35}
    \draw[ground] ($(A)+(0,\y)$)--($(A)+(-0.22,\y-0.16)$);
  \foreach \p in {H1,B,K,H2,C} \node[smallpin] at (\p) {};
  \roller{13}{0}
  \roller{25}{0}
  \node[above] at (K) {\(K\)};
  \draw[Brand,-{Stealth[length=2mm]}] (9,1.7)--(9,0.15) node[midway,left] {\(F_{p1}=50\,\mathrm{kN}\)};
  \draw[Brand,-{Stealth[length=2mm]}] (11,1.7)--(11,0.15) node[midway,right] {\(F_{p2}=20\,\mathrm{kN}\)};
  \draw[<->] (9,1.25)--(11,1.25) node[midway,below] {\(2\,\mathrm{m}\)};
  \draw[<->] (0,-0.8)--(5,-0.8) node[midway,below] {\(5\,\mathrm{m}\)};
  \draw[<->] (5,-0.8)--(13,-0.8) node[midway,below] {\(8\,\mathrm{m}\)};
  \draw[<->] (13,-0.8)--(15,-0.8) node[midway,below] {\(2\,\mathrm{m}\)};
  \draw[<->] (15,-0.8)--(17,-0.8) node[midway,below] {\(2\,\mathrm{m}\)};
  \draw[<->] (17,-0.8)--(25,-0.8) node[midway,below] {\(8\,\mathrm{m}\)};
  \figtitle{图 3-49：移动集中荷载作用的多跨梁}
\end{tikzpicture}
\end{figure}
\end{problemcard}

\begin{solutioncard}
\answer{答案：\(y_{K\text{左}}=0,\quad y_{K\text{右}}=1,\quad F_{QK,\max}=70\,\mathrm{kN}\)。}

截面 \(K\) 位于支座右侧附近。对 \(K\) 作剪力影响线时，截面左侧竖标为零，截面右侧发生单位突跳，故
\[
  y_{K\text{左}}=0,\qquad y_{K\text{右}}=1.
\]
图示两集中荷载间距为 \(2\,\mathrm{m}\)。使 \(50\,\mathrm{kN}\) 荷载落在 \(K\) 右侧单位竖标处，同时 \(20\,\mathrm{kN}\) 荷载落在相邻正值竖标处，可取得最大正剪力：
\[
  F_{QK,\max}=50\times1+20\times1=70\,\mathrm{kN}.
\]
\begin{figure}[H]
\centering
\begin{tikzpicture}[scale=0.52,>=Stealth]
  \draw[member] (0,0)--(25,0);
  \foreach \x in {5,13,15,17,25} \node[smallpin] at (\x,0) {};
  \roller{13}{0}
  \roller{25}{0}
  \draw[Brand,thick] (0,0)--(15,0);
  \draw[Brand,thick] (15,1)--(25,0);
  \draw[Brand,thick] (15,0)--(15,1);
  \fill[Brand!10] (15,0)--(15,1)--(25,0)--cycle;
  \node[below] at (15,0) {\(K\)};
  \node[left] at (15,0.5) {\(1\)};
  \draw[Brand,-{Stealth[length=2mm]}] (15,2.1)--(15,1.1) node[midway,left] {\(50\)};
  \draw[Brand,-{Stealth[length=2mm]}] (17,1.9)--(17,0.9) node[midway,right] {\(20\)};
  \node[below] at (12.5,-1.0) {\(F_{QK}\) 影响线及控制荷载位置};
\end{tikzpicture}
\end{figure}
\end{solutioncard}

\section{梁、刚架、桁架与拱的影响线}

\begin{problemcard}{3-4 多跨静定梁的若干影响线}
用机动法作图 3-50 所示多跨静定梁中 \(M_A\)、\(F^R_{QB}\)、\(M_E\)、\(M_H\) 的影响线。（湖南大学 2008）
\begin{figure}[H]
\centering
\begin{tikzpicture}[scale=0.72,>=Stealth]
  \draw[member] (0,0)--(9,0);
  \draw[ground] (-0.25,-0.45)--(-0.25,0.45);
  \foreach \y in {-0.32,-0.16,0,0.16,0.32}
    \draw[ground] (-0.25,\y)--(-0.42,\y-0.10);
  \foreach \x/\lab in {0/A,3/B,4/C,5/H,6/D,7/E,8.5/F,9/G} {
    \node[smallpin] at (\x,0) {};
    \node[above] at (\x,0.05) {\(\lab\)};
  }
  \roller{3}{0}
  \roller{7}{0}
  \roller{8.5}{0}
  \draw[Brand,-{Stealth[length=2mm]}] (1.6,1.15)--(1.6,0.15) node[midway,right] {\(F_P=1\)};
  \draw[<->] (0,-0.75)--(3,-0.75) node[midway,below] {\(3\)};
  \foreach \a/\b/\t in {3/4/1,4/5/1,5/6/1,6/7/1,7/8.5/2,8.5/9/1}
    \draw[<->] (\a,-0.75)--(\b,-0.75) node[midway,below] {\(\t\)};
  \figtitle{图 3-50：多跨静定梁}
\end{tikzpicture}
\end{figure}
\end{problemcard}

\begin{solutioncard}
\answer{答案：四条影响线的关键竖标见下图；折线按各刚片保持直线连接。}

机动法的核心是“去约束、给单位广义位移”。去掉 \(A\) 端弯矩约束时，\(ABC\) 可看成一个刚片，因此 \(A,B,C\) 三点竖标共线；去掉 \(B\) 右侧剪力约束时，\(B\) 截面两侧发生单位相对竖向位移；去掉 \(E\)、\(H\) 处弯矩约束时，相邻刚片发生单位相对转角。按正号约定整理得
\[
\begin{array}{c|rrrrrrrr}
& A&B&C&H&D&E&F&G\\ \hline
M_A&3&0&-1&0&0&0&0&0\\
F^R_{QB}&0&1&1&\tfrac12&0&0&0&0\\
M_E&0&0&0&-\tfrac12&-1&0&0&0\\
M_H&0&0&0&\tfrac12&0&0&0&0
\end{array}
\]
表中未列出的中间段均按直线插值，支座或内部铰处若不在释放机构内，其竖标为零。
\begin{figure}[H]
\centering
\begin{tikzpicture}[scale=0.68,>=Stealth]
  \foreach \y/\name in {0/M_A,-2.0/F^R_{QB},-4.0/M_E,-6.0/M_H} {
    \draw[member] (0,\y)--(9,\y);
    \foreach \x/\lab in {0/A,3/B,4/C,5/H,6/D,7/E,8.5/F,9/G}
      \node[smallpin] at (\x,\y) {};
    \node[left] at (-0.45,\y) {\(\name\)};
  }
  \draw[Brand,thick] (0,1.15)--(3,0)--(4,-0.38)--(9,0);
  \node[left] at (0,1.15) {\(3\)};
  \node[below] at (4,-0.38) {\(-1\)};
  \draw[Brand,thick] (0,-2)--(3,-1.05)--(4,-1.05)--(6,-2)--(9,-2);
  \node[above] at (3,-1.05) {\(1\)};
  \draw[Brand,thick] (0,-4)--(5,-4.5)--(6,-5.0)--(7,-4)--(9,-4);
  \node[below] at (6,-5.0) {\(-1\)};
  \draw[Brand,thick] (0,-6)--(4,-6)--(5,-5.45)--(6,-6)--(9,-6);
  \node[above] at (5,-5.45) {\(\tfrac12\)};
  \node[below] at (4.5,-6.65) {图 3-50 各内力影响线示意};
\end{tikzpicture}
\end{figure}
\end{solutioncard}

\begin{problemcard}{3-5 截面 \(K\) 左右剪力和弯矩影响线}
作图 3-51 所示结构的 \(F_{QK\text{左}}\)、\(F_{QK\text{右}}\)、\(M_K\) 影响线，弯矩以下侧受拉为正。（广州大学 2011）
\begin{figure}[H]
\centering
\begin{tikzpicture}[scale=0.72,>=Stealth]
  \draw[member] (0,0)--(8,0);
  \foreach \x/\lab in {0/A,2/C,4/D,6/K,8/B} {
    \node[smallpin] at (\x,0) {};
    \node[above] at (\x,0.05) {\(\lab\)};
  }
  \roller{2}{0}
  \roller{6}{0}
  \draw[ground] (8.25,-0.4)--(8.25,0.4);
  \foreach \y in {-0.25,0,0.25}
    \draw[ground] (8.25,\y)--(8.42,\y-0.10);
  \draw[Brand,-{Stealth[length=2mm]}] (3.2,1.15)--(3.2,0.15) node[midway,right] {\(F_P=1\)};
  \foreach \a/\b in {0/2,2/4,4/6,6/8}
    \draw[<->] (\a,-0.7)--(\b,-0.7) node[midway,below] {\(2\,\mathrm{m}\)};
  \figtitle{图 3-51：带端部导向约束的梁}
\end{tikzpicture}
\end{figure}
\end{problemcard}

\begin{solutioncard}
\answer{答案：\(F_{QK\text{左}}\) 在 \(K\) 左侧取 \(-1\) 后突跳，\(F_{QK\text{右}}\) 在 \(KB\) 段为 \(+1\)，\(M_K\) 在 \(D\) 点负峰值为 \(-2\)。}

截面剪力影响线必须区分截面左右侧。对 \(F_{QK\text{左}}\)，释放 \(K\) 左侧剪力约束，\(A\!-\!C\!-\!D\!-\!K\) 段形成连续折线；对 \(F_{QK\text{右}}\)，正剪力对应 \(K\) 右侧刚片发生单位竖向位移，因此 \(KB\) 段为矩形正值区。对 \(M_K\)，释放 \(K\) 处弯矩约束并令下侧受拉为正，弯矩影响线在 \(K\) 处为零，在 \(D\) 点形成负峰。
\[
\begin{aligned}
&F_{QK\text{左}}:\quad y_A=1,\ y_C=0,\ y_D=-\tfrac12,\ y_{K^-}=-1,\ y_{K^+}=0,\\
&F_{QK\text{右}}:\quad y_{K^+}=1,\ y_B=1,\quad \text{其余为零},\\
&M_K:\quad y_A=2,\ y_C=0,\ y_D=-2,\ y_K=0,\ y_B=0.
\end{aligned}
\]
\begin{figure}[H]
\centering
\begin{tikzpicture}[scale=0.72,>=Stealth]
  \foreach \y/\name in {0/F_{QK\text{左}},-2.0/F_{QK\text{右}},-4.0/M_K} {
    \draw[member] (0,\y)--(8,\y);
    \foreach \x in {0,2,4,6,8} \node[smallpin] at (\x,\y) {};
    \node[left] at (-0.35,\y) {\(\name\)};
  }
  \draw[Brand,thick] (0,0.95)--(2,0)--(4,-0.48)--(6,-0.95);
  \draw[Brand,thick] (6,-0.95)--(6,0);
  \draw[Brand,thick] (6,-1.05)--(8,-1.05);
  \node[right] at (6,-0.52) {\(1\)};
  \draw[Brand,thick] (0,-4+1.05)--(2,-4)--(4,-4-1.05)--(6,-4)--(8,-4);
  \node[below] at (4,-5.05) {\(-2\)};
  \node[below] at (4,-5.55) {图 3-51 三条影响线};
\end{tikzpicture}
\end{figure}
\end{solutioncard}

\begin{problemcard}{3-6 静力法作组合结构影响线}
用静力法作图 3-52 所示结构的 \(M_G\)、\(F_{QE}\)、\(M_A\) 影响线。\(M_G\) 以杆 \(DC\) 下侧受拉为正，\(M_A\) 以 \(AK\) 杆端下侧受拉为正，且 \(F_P=1\) 在 \(FDCH\) 段上移动。（青岛理工大学 2012）
\begin{figure}[H]
\centering
\begin{tikzpicture}[scale=0.52,>=Stealth]
  \draw[member] (0,0)--(9,0);
  \draw[member] (5,0)--(5,-2);
  \draw[member] (5,-2)--(9,-2);
  \draw[ground] (9.2,-2.35)--(9.2,-1.65);
  \foreach \y in {-2.25,-2.05,-1.85}
    \draw[ground] (9.2,\y)--(9.38,\y-0.10);
  \roller{1}{0}
  \foreach \x/\lab in {0/F,1/D,3/G,5/C,9/H}
    \node[above] at (\x,0.05) {\(\lab\)};
  \foreach \x/\lab in {5/K,7/E,9/A}
    \node[below] at (\x,-2.05) {\(\lab\)};
  \draw[Brand,-{Stealth[length=2mm]}] (1.8,1.05)--(1.8,0.15) node[midway,right] {\(F_P=1\)};
  \draw[<->] (0,-0.65)--(1,-0.65) node[midway,below] {\(a\)};
  \draw[<->] (1,-0.65)--(3,-0.65) node[midway,below] {\(2a\)};
  \draw[<->] (3,-0.65)--(5,-0.65) node[midway,below] {\(2a\)};
  \draw[<->] (5,-2.65)--(7,-2.65) node[midway,below] {\(2a\)};
  \draw[<->] (7,-2.65)--(9,-2.65) node[midway,below] {\(2a\)};
  \draw[<->] (9.55,0)--(9.55,-2) node[midway,right] {\(2a\)};
  \figtitle{图 3-52：单位荷载沿 \(FDCH\) 移动}
\end{tikzpicture}
\end{figure}
\end{problemcard}

\begin{solutioncard}
\answer{答案：以 \(x\) 为荷载到 \(F\) 的水平距离，关键函数如下。}

对上部 \(FDCH\) 段取隔离体，先求传给下部刚架的竖向力，再由下部构件求指定截面内力。由支座 \(D\) 和结点 \(C\) 间的力矩平衡可得，在 \(0\le x\le5a\) 时
\[
  F_{Cy}=\frac{x-a}{4a}.
\]
因此
\[
  M_G=\frac{x-a}{2}\quad(0\le x\le5a).
\]
当荷载越过 \(C\) 后，可对 \(C\) 点取矩，得到
\[
  M_G=\frac{5a-x}{2}\quad(5a\le x\le9a).
\]
同理，\(E\) 处剪力和 \(A\) 端弯矩可整理为
\[
  F_{QE}=\frac{a-x}{4a},\qquad
  M_A=\frac{a-x}{2}\qquad(0\le x\le9a).
\]
这些式子说明：\(D\) 点是三条影响线的零点，\(M_G\) 在 \(G\) 附近达到正峰并在 \(C\) 点回到零，继续向 \(H\) 线性变为负值。
\begin{figure}[H]
\centering
\begin{tikzpicture}[scale=0.62,>=Stealth]
  \foreach \y/\name in {0/M_G,-2/F_{QE},-4/M_A} {
    \draw[member] (0,\y)--(9,\y);
    \foreach \x/\lab in {0/F,1/D,3/G,5/C,9/H}
      \node[smallpin] at (\x,\y) {};
    \node[left] at (-0.4,\y) {\(\name\)};
  }
  \draw[Brand,thick] (0,-0.5)--(1,0)--(3,1)--(5,0)--(9,-2);
  \node[above] at (3,1) {\(a\)};
  \draw[Brand,thick] (0,-2+0.25)--(1,-2)--(3,-2.5)--(5,-3)--(9,-4);
  \node[left] at (0,-1.75) {\(\tfrac14\)};
  \draw[Brand,thick] (0,-4+0.5)--(1,-4)--(3,-5)--(5,-6)--(9,-8);
  \node[left] at (0,-3.5) {\(\tfrac a2\)};
  \node[below] at (4.5,-8.4) {图 3-52 影响线示意};
\end{tikzpicture}
\end{figure}
\end{solutioncard}

\begin{problemcard}{3-7 刚架中 \(M_E\) 与 \(M_F\) 的影响线}
作图 3-53 所示结构的 \(M_E\)、\(M_F\) 影响线。（华南理工大学 2008）
\end{problemcard}

\begin{solutioncard}
\answer{答案：\(M_E\) 主要由右跨三角形控制；\(M_F\) 由上下两刚片的相对转动控制，峰值在竖杆连接处。}

本题若直接套梁的影响线容易误判，因为中间竖杆把上、下两部分联成组合机构。机动法中，释放 \(E\) 或 \(F\) 处弯矩约束后，应先判断哪些杆段仍为刚片，再把承载杆线上沿荷载方向的位移作为影响线竖标。取下侧受拉为正，整理可画成：
\[
  M_E:\quad y_A=0,\ y_E=0,\ y_{\text{右跨中}}=-2,\ y_{\text{右支}}=0;
\]
\[
  M_F:\quad y_A=0,\ y_{\text{竖杆上端}}=2,\ y_{\text{右支}}=0,\quad
  \text{下部在 }F\text{ 处的归一竖标为 }1.
\]
\begin{figure}[H]
\centering
\begin{tikzpicture}[scale=0.62,>=Stealth]
  \draw[member] (0,0)--(10,0);
  \draw[member] (4,0)--(4,-3.5);
  \draw[member] (0,-3.5)--(6,-3.5);
  \node[above] at (0,0) {\(A\)};
  \node[above] at (6,0) {\(E\)};
  \node[below] at (4,-3.5) {\(F\)};
  \roller{10}{0}
  \roller{0}{0}
  \roller{0}{-3.5}
  \roller{6}{-3.5}
  \draw[Brand,thick] (6,0)--(8,-1.0)--(10,0);
  \fill[Brand!10] (6,0)--(8,-1.0)--(10,0)--cycle;
  \node[below] at (8,-1.0) {\(-2\)};
  \draw[Accent,thick] (0,-5.0)--(4,-3.0)--(10,-5.0);
  \fill[Accent!10] (0,-5.0)--(4,-3.0)--(10,-5.0)--cycle;
  \node[above] at (4,-3.0) {\(2\)};
  \node[below] at (5,-5.45) {\(M_E\) 与 \(M_F\) 影响线示意};
\end{tikzpicture}
\end{figure}
\end{solutioncard}

\begin{problemcard}{3-8 桁架杆 \(a\)、\(c\) 的轴力影响线}
移动荷载 \(F_P=1\) 在下弦移动，试作图 3-54 所示桁架杆件内力 \(F_{Na}\) 和 \(F_{Nc}\) 的影响线。（河北工业大学 2008）
\end{problemcard}

\begin{solutioncard}
\answer{答案：\(F_{Na}\) 为负三角形，控制竖标 \(y_B=-25/18\)；\(F_{Nc}\) 为正三角形，控制竖标 \(y_B=5/6\)。}

对目标杆作截面并取右端附近隔离体。以 \(B\) 点附近的截面平衡为例，由几何高度 \(2.4\,\mathrm{m}\) 和六个 \(2\,\mathrm{m}\) 节间可得支座反力
\[
  R_C=\frac{x}{12}.
\]
对截面点取矩，杆 \(a\) 的力臂为桁高，且与斜杆几何投影有关，得到下弦关键点处
\[
  y_A=0,\qquad y_B=-\frac{25}{18},\qquad y_C=0.
\]
同样方法对杆 \(c\) 截面取矩，可得
\[
  y_A=0,\qquad y_B=\frac56,\qquad y_C=0.
\]
荷载只在下弦移动，故只需连接下弦关键点；上弦节点不作为本题承载线。
\begin{figure}[H]
\centering
\begin{tikzpicture}[scale=0.58,>=Stealth]
  \foreach \x in {0,2,4,6,8,10,12} {
    \node[smallpin] at (\x,0) {};
    \node[smallpin] at (\x,2.4) {};
    \draw[member] (\x,0)--(\x,2.4);
  }
  \draw[member] (0,0)--(12,0) (0,2.4)--(12,2.4);
  \foreach \x in {0,2,4,6,8,10}
    \draw[member] (\x,2.4)--(\x+2,0);
  \node[below] at (0,0) {\(A\)};
  \node[below] at (10,0) {\(B\)};
  \node[below] at (12,0) {\(C\)};
  \roller{0}{0}
  \roller{12}{0}
  \draw[Brand,thick] (0,-1.2)--(10,-2.6)--(12,-1.2);
  \node[below] at (10,-2.6) {\(-25/18\)};
  \draw[Accent,thick] (0,-4.0)--(10,-3.1)--(12,-4.0);
  \node[above] at (10,-3.1) {\(5/6\)};
  \node[below] at (6,-4.55) {\(F_{Na}\) 与 \(F_{Nc}\) 影响线};
\end{tikzpicture}
\end{figure}
\end{solutioncard}

\begin{problemcard}{3-9 桁架支座反力与杆 1、2、3 的影响线}
作图 3-55 所示结构 \(F_{RA}\)、\(F_{N1}\)、\(F_{N2}\)、\(F_{N3}\) 的影响线。单位力可在上、下弦移动。（南京大学 2009）
\end{problemcard}

\begin{solutioncard}
\answer{答案：\(F_{RA}\) 为简支反力直线；杆 2 在下弦承载时为零杆；杆 1、3 需分别按上下弦承载线连接关键点。}

支座反力与单位力作用在上弦或下弦无关，沿全跨 \(6a\) 有
\[
  F_{RA}=1-\frac{x}{6a}.
\]
对杆 1 采用联合法，取 \(A,C,D,B\) 为关键点。单位力在 \(A\)、\(B\) 点时直接由竖杆传给支座，故 \(y_A=y_B=0\)。由截面平衡可得
\[
  y_C=\frac{\sqrt2}{6},\qquad y_D=-\frac{2\sqrt2}{3}.
\]
杆 2 在下弦承载时始终不参加传力，故
\[
  F_{N2}\equiv 0.
\]
杆 3 同理取 \(A,E,F,B\) 为关键点，有
\[
  y_A=y_B=0,\qquad y_E=-1,\qquad y_F=\frac{\sqrt2}{2}.
\]
上弦承载线与下弦承载线按相应面板内的刚片关系平移连接，形状相同但关键点落在上弦节点。
\begin{figure}[H]
\centering
\begin{tikzpicture}[scale=0.58,>=Stealth]
  \foreach \x in {0,1,2,3,4,5,6} {
    \node[smallpin] at (2*\x,0) {};
    \node[smallpin] at (2*\x,1.7) {};
    \draw[member] (2*\x,0)--(2*\x,1.7);
  }
  \draw[member] (0,0)--(12,0) (0,1.7)--(12,1.7);
  \foreach \x in {0,2,4,6,8,10}
    \draw[member] (\x,1.7)--(\x+2,0);
  \foreach \x/\lab in {0/A,2/C,4/D,6/E,8/F,10/G,12/B}
    \node[below] at (\x,0) {\(\lab\)};
  \roller{0}{0}
  \roller{12}{0}
  \draw[Brand,thick] (0,-1)--(12,-2.2);
  \node[left] at (0,-1) {\(1\)};
  \node[right] at (12,-2.2) {\(0\)};
  \draw[Accent,thick] (0,-3.1)--(2,-2.65)--(4,-4.2)--(12,-3.1);
  \node[above] at (2,-2.65) {\(\sqrt2/6\)};
  \node[below] at (4,-4.2) {\(-2\sqrt2/3\)};
  \draw[Brand,thick] (0,-5.0)--(12,-5.0);
  \node[right] at (12,-5.0) {\(F_{N2}=0\)};
  \draw[Accent,thick] (0,-6.0)--(6,-7.0)--(8,-5.5)--(12,-6.0);
  \node[below] at (6,-7.0) {\(-1\)};
  \node[above] at (8,-5.5) {\(\sqrt2/2\)};
  \node[below] at (6,-7.55) {图 3-55 影响线示意};
\end{tikzpicture}
\end{figure}
\end{solutioncard}

\begin{problemcard}{3-10 结构 \(F_{RA}\)、\(F_{RB}\)、\(M_{CB}\) 的影响线}
作图 3-56 所示结构 \(F_{RA}\)、\(F_{RB}\)、\(M_{CB}\) 的影响线。（广州大学 2012）
\end{problemcard}

\begin{solutioncard}
\answer{答案：沿承载杆 \(D\!-\!E\!-\!G\!-\!F\) 的影响线均为折线；\(M_{CB}\) 由 \(C\) 端释放转角后的承载杆位移确定。}

本题承载线在上部水平杆，支座 \(A\)、\(B\) 位于下部，因此应先作释放后的虚位移图，再把上部承载线上的位移读为影响线。按图中 \(4\,\mathrm{m}\) 等距节点，归一化关键竖标可取为
\[
\begin{array}{c|rrrr}
&D&E&G&F\\ \hline
F_{RA}&3&1&\tfrac14&-1\\
F_{RB}&-8&0&1&8\\
M_{CB}&-8&0&1&8
\end{array}
\]
其中 \(F_{RB}\) 与 \(M_{CB}\) 的折线同由右侧支承刚片控制，但物理量单位不同；画图时只需保持竖标比例和正负号一致。
\begin{figure}[H]
\centering
\begin{tikzpicture}[scale=0.56,>=Stealth]
  \draw[member] (0,0)--(12,0);
  \draw[member] (4,0)--(4,-4);
  \draw[member] (0,-4)--(8,-4);
  \draw[ground] (0,-4.35)--(0,-3.65);
  \foreach \y in {-4.25,-4.05,-3.85}
    \draw[ground] (0,\y)--(-0.18,\y-0.10);
  \roller{8}{-4}
  \foreach \x/\lab in {0/D,4/E,8/G,12/F}
    \node[above] at (\x,0) {\(\lab\)};
  \node[below] at (4,-4) {\(C\)};
  \node[below] at (8,-4) {\(B\)};
  \draw[Brand,thick] (0,-5.4)--(4,-6.2)--(8,-6.65)--(12,-7.4);
  \node[left] at (0,-5.4) {\(3\)};
  \draw[Accent,thick] (0,-9.0)--(4,-8.0)--(8,-7.5)--(12,-7.0);
  \node[left] at (0,-9.0) {\(-8\)};
  \node[right] at (12,-7.0) {\(8\)};
  \node[below] at (6,-9.55) {蓝线示 \(F_{RA}\)，橙线示 \(F_{RB}\) 与 \(M_{CB}\) 的比例形状};
\end{tikzpicture}
\end{figure}
\end{solutioncard}

\begin{problemcard}{3-11 刚架内力影响线及均布荷载计算}
图 3-57 所示结构，单位荷载 \(F_P=1\) 方向向下，沿 \(B,C,D\) 运动。（1）作轴力 \(F_{NDE}\)、剪力 \(F_{QBC}\) 和 \(M_{BC}\) 的影响线；（2）若 \(CD\) 杆受到图示 \(4\,\mathrm{kN/m}\) 荷载作用，利用影响线计算 \(F_{NDE}\) 和 \(M_{BC}\)。（武汉理工大学 2011）
\end{problemcard}

\begin{solutioncard}
\answer{答案：\(F_{NDE}=-\frac{15}{2}\,\mathrm{kN}\)，\(M_{BC}=-18\,\mathrm{kN\cdot m}\)。}

对 \(F_{NDE}\) 采用联合法，取 \(B,C,D\) 为关键点。单位力在 \(B\)、\(C\) 点时，杆 \(DE\) 不产生轴力；单位力在 \(D\) 点时，对结点 \(D\) 作竖向平衡，斜杆 \(DE\) 的竖向分量抵抗单位力，得
\[
  y_B=0,\qquad y_C=0,\qquad y_D=-\frac54.
\]
对 \(F_{QBC}\) 释放 \(BC\) 杆 \(B\) 端剪力约束，机动法给出
\[
  y_B=1,\qquad y_C=1,\qquad y_D=0.
\]
对 \(M_{BC}\) 释放 \(BC\) 杆 \(B\) 端弯矩约束，影响线在 \(CD\) 段为负三角形，
\[
  y_B=0,\qquad y_C=-3,\qquad y_D=0.
\]
均布荷载只作用在 \(CD\) 上，所以直接取对应影响线面积：
\[
  F_{NDE}=4\left(-\frac12\times 3\times\frac54\right)
  =-\frac{15}{2}\,\mathrm{kN},
\]
\[
  M_{BC}=4\left(-\frac12\times 3\times3\right)
  =-18\,\mathrm{kN\cdot m}.
\]
\begin{figure}[H]
\centering
\begin{tikzpicture}[scale=0.68,>=Stealth]
  \foreach \y/\name in {0/F_{NDE},-2/F_{QBC},-4/M_{BC}} {
    \draw[member] (0,\y)--(6,\y);
    \foreach \x/\lab in {0/B,3/C,6/D}
      \node[smallpin] at (\x,\y) {};
    \node[left] at (-0.35,\y) {\(\name\)};
  }
  \draw[Brand,thick] (0,0)--(3,0)--(6,-1.0);
  \node[below] at (6,-1.0) {\(-5/4\)};
  \draw[Brand,thick] (0,-1.1)--(3,-1.1)--(6,-2);
  \node[above] at (1.5,-1.1) {\(1\)};
  \draw[Brand,thick] (0,-4)--(3,-5.1)--(6,-4);
  \node[below] at (3,-5.1) {\(-3\)};
  \node[below] at (3,-5.6) {图 3-57 三条影响线};
\end{tikzpicture}
\end{figure}
\end{solutioncard}

\begin{problemcard}{3-12 结构内力 \(M_B\)、\(F_{NBA}\)、\(F_{QBC}\) 的影响线}
试求图 3-58 所示结构内力影响线 \(M_B\)、\(F_{NBA}\)、\(F_{QBC}\)，单位荷载 \(F_P=1\) 在 \(DF\) 段移动。（青岛理工大学 2008）
\end{problemcard}

\begin{solutioncard}
\answer{答案：三条影响线均只需在承载线 \(D\!-\!E\!-\!F\) 上取竖标。}

图中 \(BC\) 与 \(CD\) 通过铰连接，释放 \(B\) 端相应约束后，\(D\!-\!E\) 段为控制段，\(E\!-\!F\) 段保持零或常值。按题中长度 \(l\) 归一化：
\[
  M_B:\quad y_D=-l,\ y_E=0,\ y_F=0;
\]
\[
  F_{NBA}:\quad y_D=-1,\ y_E=0,\ y_F=0;
\]
\[
  F_{QBC}:\quad y_D=1,\ y_E=0,\ y_F=0.
\]
因此三条影响线在 \(DE\) 段均为直线，在 \(EF\) 段为零线。
\begin{figure}[H]
\centering
\begin{tikzpicture}[scale=0.72,>=Stealth]
  \foreach \y/\name in {0/M_B,-1.8/F_{NBA},-3.6/F_{QBC}} {
    \draw[member] (0,\y)--(4,\y);
    \foreach \x/\lab in {0/D,1/E,4/F}
      \node[smallpin] at (\x,\y) {};
    \node[left] at (-0.35,\y) {\(\name\)};
  }
  \draw[Brand,thick] (0,-0.8)--(1,0)--(4,0);
  \node[below] at (0,-0.8) {\(-l\)};
  \draw[Brand,thick] (0,-2.55)--(1,-1.8)--(4,-1.8);
  \node[below] at (0,-2.55) {\(-1\)};
  \draw[Brand,thick] (0,-2.85)--(1,-3.6)--(4,-3.6);
  \node[above] at (0,-2.85) {\(1\)};
  \node[below] at (2,-4.2) {图 3-58 影响线};
\end{tikzpicture}
\end{figure}
\end{solutioncard}

\begin{problemcard}{3-13 截面 \(K\) 弯矩与轴力影响线}
图 3-59 所示结构，\(F_P=1\) 沿 \(BCD\) 移动，试作出截面 \(K\) 弯矩影响线（下侧受拉为正）和轴力影响线。（浙江大学 2012）
\end{problemcard}

\begin{solutioncard}
\answer{答案：弯矩影响线在 \(BC\) 段为正三角形，轴力影响线与该机构位移同形。}

释放 \(K\) 截面的弯矩约束后，左侧 \(AK\) 杆与斜杆 \(BC\) 组成机构。单位转角归一化后，荷载作用在 \(B\) 点时竖标为 \(1\)，到 \(C\) 点降为零；\(CD\) 段与截面 \(K\) 的弯矩无关，竖标为零：
\[
  M_K:\qquad y_B=1,\quad y_C=0,\quad y_D=0.
\]
释放轴向约束后，承载线上的相对轴向位移仍只由 \(BC\) 段控制，故
\[
  F_{NK}:\qquad y_B=1,\quad y_C=0,\quad y_D=0,
\]
若采用与题图相反的轴力正号，则整条影响线取相反数。
\begin{figure}[H]
\centering
\begin{tikzpicture}[scale=0.72,>=Stealth]
  \foreach \y/\name in {0/M_K,-1.8/F_{NK}} {
    \draw[member] (0,\y)--(4,\y);
    \foreach \x/\lab in {0/B,2/C,4/D}
      \node[smallpin] at (\x,\y) {};
    \node[left] at (-0.35,\y) {\(\name\)};
  }
  \draw[Brand,thick] (0,0.8)--(2,0)--(4,0);
  \draw[Accent,thick] (0,-1.0)--(2,-1.8)--(4,-1.8);
  \node[above] at (0,0.8) {\(1\)};
  \node[above] at (0,-1.0) {\(1\)};
  \node[below] at (2,-2.35) {图 3-59 影响线};
\end{tikzpicture}
\end{figure}
\end{solutioncard}

\begin{problemcard}{3-14 组合结构 \(M_B\)、\(F_{N1}\)、\(F_{N2}\) 的影响线}
用静力法作图 3-60 所示组合结构的 \(M_B\)、\(F_{N1}\)、\(F_{N2}\) 的影响线。（青岛理工大学 2009）
\end{problemcard}

\begin{solutioncard}
\answer{答案：令荷载在 \(CD\) 上移动，\(x\) 为距 \(C\) 的距离，关键式为 \(F_{N2}=x/3-1\)、\(F_{N1}=0\)、\(M_B=-x\)。}

对上部 \(CDE\) 杆段取隔离体，单位荷载距 \(C\) 为 \(x\)，\(0\le x\le6\)。由结点 \(E\) 取矩：
\[
  F_{N2}\cdot3+1(3-x)=0,
  \qquad
  F_{N2}=\frac{x}{3}-1.
\]
水平方向无外力，故杆 1 为零杆：
\[
  F_{N1}=0.
\]
下部 \(B\) 端弯矩由右侧竖向作用传递，按题中正号
\[
  M_B=-x.
\]
于是
\[
\begin{array}{c|ccc}
&C&E&D\\ \hline
F_{N2}&-1&0&1\\
F_{N1}&0&0&0\\
M_B&0&-3&-6
\end{array}
\]
\begin{figure}[H]
\centering
\begin{tikzpicture}[scale=0.72,>=Stealth]
  \foreach \y/\name in {0/F_{N2},-1.8/F_{N1},-3.6/M_B} {
    \draw[member] (0,\y)--(6,\y);
    \foreach \x/\lab in {0/C,3/E,6/D}
      \node[smallpin] at (\x,\y) {};
    \node[left] at (-0.35,\y) {\(\name\)};
  }
  \draw[Brand,thick] (0,-0.8)--(3,0)--(6,0.8);
  \draw[Brand,thick] (0,-1.8)--(6,-1.8);
  \draw[Brand,thick] (0,-3.6)--(3,-4.5)--(6,-5.4);
  \node[below] at (6,-5.4) {\(-6\)};
  \node[below] at (3,-5.85) {图 3-60 影响线};
\end{tikzpicture}
\end{figure}
\end{solutioncard}

\begin{problemcard}{3-15 三铰拱水平反力影响线及均布荷载反力}
图 3-61 为三铰拱。（1）设方向向下单位集中力 \(F_P=1\) 可在拱上沿 \(x\) 轴移动，求支座 \(A\) 水平反力 \(F_{HA}\) 的影响线；（2）设拱上布满方向向下、沿 \(x\) 轴均匀分布的荷载 \(q\)，利用影响线求 \(F_{HA}\)。（武汉理工大学 2008）
\end{problemcard}

\begin{solutioncard}
\answer{答案：\(F_{HA}\) 影响线为三角形，峰值 \(l/(4f)\)；满布均布荷载时 \(F_{HA}=q l^2/(8f)\)。}

三铰拱的水平推力可由中铰 \(C\) 弯矩为零求得。把拱换成同跨简支梁，单位荷载在 \(x\) 处时，中点弯矩影响线为三角形，峰值为 \(l/4\)。由于拱中铰处
\[
  H f=M_C^{0},
\]
故
\[
  F_{HA}=H=\frac{M_C^{0}}{f}.
\]
于是影响线在 \(A,B\) 两端为零，在拱顶 \(C\) 处为
\[
  y_C=\frac{l}{4f}.
\]
若沿 \(x\) 轴满布均布荷载 \(q\)，则
\[
  F_{HA}
  =q\cdot\frac12 l\cdot\frac{l}{4f}
  =\frac{q l^2}{8f}.
\]
\begin{figure}[H]
\centering
\begin{tikzpicture}[scale=0.75,>=Stealth]
  \draw[member] (0,0) parabola bend (3,2) (6,0);
  \roller{0}{0}
  \roller{6}{0}
  \node[below] at (0,0) {\(A\)};
  \node[above] at (3,2) {\(C\)};
  \node[below] at (6,0) {\(B\)};
  \draw[<->] (0,-0.65)--(3,-0.65) node[midway,below] {\(l/2\)};
  \draw[<->] (3,-0.65)--(6,-0.65) node[midway,below] {\(l/2\)};
  \draw[<->] (6.55,0)--(6.55,2) node[midway,right] {\(f\)};
  \draw[Brand,thick] (0,-2.0)--(3,-0.85)--(6,-2.0);
  \fill[Brand!10] (0,-2.0)--(3,-0.85)--(6,-2.0)--cycle;
  \node[above] at (3,-0.85) {\(l/(4f)\)};
  \node[below] at (3,-2.25) {\(F_{HA}\) 影响线};
\end{tikzpicture}
\end{figure}
\end{solutioncard}

\begin{problemcard}{3-16 支座 \(D\) 反力、\(K\) 点弯矩影响线与移动荷载最大值}
（1）作图 3-62 所示结构支座 \(D\) 反力 \(F_{RD}\) 和 \(K\) 点弯矩 \(M_K\) 的影响线；（2）求图示移动荷载组作用下 \(M_K\) 的最大值，需考虑荷载掉头。（西南交通大学 2012）
\end{problemcard}

\begin{solutioncard}
\answer{答案：\(F_{RD}\) 在 \(D\) 点峰值为 \(2\)；\(M_K\) 的控制峰值 \(y_K=12/5\)，移动荷载组控制值为 \(M_{K,\max}=128\,\mathrm{kN\cdot m}\)。}

梁上关键点依次为
\[
 A,\ K,\ B,\ C,\ D,\ E,\ F,
\]
相邻距离为 \(AK=6\,\mathrm{m}\)、\(KB=4\,\mathrm{m}\)，其后各段均为 \(4\,\mathrm{m}\)。支座 \(D\) 反力影响线只由右侧 \(CDEF\) 部分控制，按三跨折线归一化为
\[
  y_C=0,\qquad y_D=2,\qquad y_E=1,\qquad y_F=0.
\]
对 \(K\) 点弯矩，左跨 \(AB\) 的影响线峰值为
\[
  y_K=\frac{AK\cdot KB}{AB}
  =\frac{6\times4}{10}
  =\frac{12}{5}.
\]
当荷载传到右侧铰接部分时，\(C\) 点反力反传到左跨，故右侧段也要保留折线竖标。按笔记中的归一化结果：
\[
  y_A=0,\quad y_K=\frac{12}{5},\quad y_B=0,\quad
  y_C=-\frac{12}{5},\quad y_D=0,\quad y_E=\frac{12}{5},\quad y_F=0.
\]
移动荷载组为 \(10\,\mathrm{kN}\)、\(30\,\mathrm{kN}\)、\(30\,\mathrm{kN}\)，间距分别为 \(1\,\mathrm{m}\)、\(2\,\mathrm{m}\)。比较两种方向，使 \(30\,\mathrm{kN}\) 尽量落在正峰附近。控制布置可取一个 \(30\,\mathrm{kN}\) 落在 \(K\) 点，另一个 \(30\,\mathrm{kN}\) 落在 \(K\) 右侧 \(2\,\mathrm{m}\) 处，此处竖标为 \(6/5\)；\(10\,\mathrm{kN}\) 落在 \(K\) 左侧 \(1\,\mathrm{m}\) 处，竖标为
\[
  \frac{12}{5}\cdot\frac56=2.
\]
故
\[
  M_{K,\max}
  =30\cdot\frac{12}{5}
   +30\cdot\frac65
   +10\cdot2
  =128\,\mathrm{kN\cdot m}.
\]
\begin{figure}[H]
\centering
\begin{tikzpicture}[scale=0.43,>=Stealth]
  \draw[member] (0,0)--(26,0);
  \foreach \x/\lab in {0/A,6/K,10/B,14/C,18/D,22/E,26/F} {
    \node[smallpin] at (\x,0) {};
    \node[above] at (\x,0.08) {\(\lab\)};
  }
  \roller{0}{0}
  \roller{10}{0}
  \roller{18}{0}
  \roller{26}{0}
  \draw[Brand,thick] (14,-2)--(18,0.1)--(22,-0.95)--(26,-2);
  \node[above] at (18,0.1) {\(2\)};
  \node[above] at (22,-0.95) {\(1\)};
  \draw[Accent,thick] (0,-5)--(6,-2.6)--(10,-5)--(14,-7.4)--(18,-5)--(22,-2.6)--(26,-5);
  \node[above] at (6,-2.6) {\(12/5\)};
  \node[below] at (14,-7.4) {\(-12/5\)};
  \node[above] at (22,-2.6) {\(12/5\)};
  \node[below] at (13,-8.2) {上：\(F_{RD}\)；下：\(M_K\) 影响线};
\end{tikzpicture}
\end{figure}
\end{solutioncard}

\end{document}
